\documentclass[11pt]{article}
\usepackage[a4paper,margin=27mm]{geometry}
\usepackage[T1]{fontenc}
\usepackage{lmodern,microtype,amsmath,amssymb,amsthm,mathtools,enumitem}
\usepackage[hidelinks,hypertexnames=false]{hyperref}
\newtheorem{theorem}{Theorem}
\newtheorem{lemma}[theorem]{Lemma}
\newtheorem{proposition}[theorem]{Proposition}
\newtheorem{corollary}[theorem]{Corollary}
\theoremstyle{definition}\newtheorem{definition}[theorem]{Definition}
\theoremstyle{remark}\newtheorem{remark}[theorem]{Remark}
\newcommand{\R}{\mathbb R}
\newcommand{\Dscr}{\mathcal D}
\newcommand{\supp}{\operatorname{supp}}
\newcommand{\debt}{\mathrm{debt}}
\newcommand{\cred}{\mathrm{credit}}
\newcommand{\sil}{\mathrm{silence}}
\newcommand{\Band}{\Xi}
\newcommand{\Phihat}{\widehat\Phi}
\newcommand{\Psitil}{\widetilde\Psi}
\newcommand{\psitil}{\tilde\psi}
\newcommand{\ftil}{\tilde f}
\newcommand{\htil}{\tilde h}
\newcommand{\one}{\mathbf 1}
\newcommand{\esssup}{\operatorname*{ess\,sup}}
\newcommand{\essinf}{\operatorname*{ess\,inf}}
\title{Two Additive Worlds Without Completeness}
\author{}\date{11 October 2026}
\begin{document}\maketitle

\noindent
This note continues IC (\emph{Independence Chains and Two Additive Worlds}),
DT (\emph{Dense Thresholds and Debt Reduction}), ECS (\emph{Edge Cycles and
Separability}) and CCF (\emph{Closed Choice Forces Fineness}).
It also uses CC (\emph{What Closed Competitive Choice Requires of
Independent Rankings}), CBOI (\emph{Closed Choice Beyond Ordinary Integrals}),
EDL (\emph{Edge Dichotomy and Leverage}) and FAED (\emph{Finitely Additive Edge
Dichotomy}). It settles the
three-world case in which two worlds carry one-sided additive weak
extensions, with no completeness assumption at the third world.

\begin{enumerate}[leftmargin=*,itemsep=3pt]
\item \emph{Main theorem} (Theorem~\ref{thm:main}). Let $n=3$ and assume (G)
at every world and (S). Let world~$1$ be any preorder satisfying (IB);
completeness is not assumed. Let worlds $2,3$ be preorders that have
one-sided weak extensions \eqref{eq:W} of CBOI's finitely additive form, with
arbitrary finitely additive charges. Then $V$ is dominated. For $n=3$ this
removes the completeness hypothesis from CCF's Theorem~18.
\item \emph{Two special cases with shorter proofs.} (I) Worlds $2,3$ are
ranked exactly by ordinary additive losses (Theorem~\ref{thm:I}). This is the
sub-case of IC, DT and ECS, which was open there for an incomplete world~$1$.
(II) Worlds $2,3$ are preorders satisfying (G) with one-sided ordinary
additive extensions (Theorem~\ref{thm:II}).
\item \emph{Routing primaries} (Theorems~\ref{thm:debtprimary}
and~\ref{thm:creditprimary}). World~$1$ compares debt sets, and separately
credit sets, strictly according to explicit additive weights $\Phi$ and
$\Psi$, whenever the lighter of the two sets is bounded (and, in the
finitely additive case, clean); the heavier set is arbitrary. These
comparisons are derived from moves at the edge selections alone. No
completeness is involved.
\item \emph{Silent targets} (Section~\ref{sec:silent}). Near the edge
$[e_2,e_3]$ there are priors whose Bayes model, for an auxiliary world-1
weight, never gives credit. At a jump corner of the separable debt boundary
this Bayes model is not coherent. A KKM dominator of it (a selected
model weakly better at every world, given by (K1), the KKM-type lemma of
CBOI; Section~\ref{sec:setting}) then yields a pure
trade of debt for credit at a ratio that the forced trades of the edge
selections forbid.
\item \emph{The finitely additive case} (Sections~\ref{sec:faclean}
to~\ref{sec:main}). Closedness of the section at the \emph{other} world of an
edge excludes the germs of the charges along the edges
(Theorems~\ref{thm:rightgerm}, \ref{thm:zerogerm}, \ref{thm:rightgerm22}
and~\ref{thm:crossgerm}). The edge $[e_2,e_3]$
has the endpoint property, and Pareto bounds at its selections control the
purely finitely additive parts there (Theorem~\ref{thm:pureratio}). Pure
separability (Theorem~\ref{thm:puresep}) gives a debt primary that weighs
finitely additive mass exactly (Theorem~\ref{thm:dirtyprimary}). With these,
the fixed KKM targets of the ordinary proof work again, and no splicing of
targets is needed.
\end{enumerate}
Completeness is used at no world. (IB) is used only at world~$1$. At
worlds~$2,3$ we use transitivity, strict (G), closedness of their sections
and \eqref{eq:W}, but neither completeness nor (IB).

\emph{What remains open.} First, the general conjecture for incomplete
rankings when no world carries an additive extension. Second, the case of
complete rankings at all worlds with no additive world. By CCF's
Theorem~16 that case reduces to the endpoint property on every edge; it has
been reduced further in separate work. Within the present setting a few
questions are left open because the proof does not need them; they are
listed in Remark~\ref{rem:notneeded}.

The arguments were checked, in the working notes on which this note is
based, by independent referees in four parts: the routing primaries and their
consequences (Sections~\ref{sec:routing} and~\ref{sec:totals}), the
silent-target argument (Section~\ref{sec:silent}), the one-sided ordinary
case (Section~\ref{sec:onesided}), and the finitely additive case
(Sections~\ref{sec:faclean}--\ref{sec:main}). Their corrections are
incorporated. Results taken from those working notes are proved in full
below.

\section{Setting and inputs}\label{sec:setting}

\subsection{Problems, models and regions}

We use the setting of CC, CBOI, EDL, FAED, IC, DT, ECS and CCF. There
are three worlds. Put $\Delta\coloneqq\{p\in\R^3:p\ge0,\
\sum_ip_i=1\}$, with vertices $e_1,e_2,e_3$; for $S\subseteq\{1,2,3\}$,
$\Delta_S$ is the face of $\Delta$ spanned by $\{e_i:i\in S\}$, and
$\supp p\coloneqq\{i:p_i>0\}$. Binary problems $\{g,-g\}$ are represented by
$g\in H\coloneqq\{g\in\R^3:g_1>0\}$, with the normalized Gaussian
measure~$\mu$. Models are measurable and identified modulo $\mu$-null sets.
At world $i$ the advice of a model $\Dscr$ at $g$ has a guidance rank
$r_i(g,\Dscr)\in\{0,1,2,3\}$: good, silent, bad, contradictory. A model is
\emph{regular} if it is never contradictory and correct on unanimous
problems. A coherent model $\Dscr_K$, with $K\subseteq\Delta$ nonempty,
compact and convex, accepts $g$ if and only if $\min_{m\in K}m\cdot g>0$
(and accepts every nonzero $g\ge0$); coherent models are regular, and
$V\coloneqq\Dscr_\Delta$. The ideal model $D_i\coloneqq\{g:g_i>0\}$ of
world~$i$ gives the advice good at~$i$ on every problem. Write $H'$ for the mixed problems in $H$ and
\[
 P_1\coloneqq\{g_2<0,\ g_3<0\},\qquad P_2\coloneqq\{g_2<0<g_3\},\qquad
 P_3\coloneqq\{g_2>0>g_3\}
\]
(inside $H$), so $H'=P_1\cup P_2\cup P_3$ up to null sets. As in IC, $o_K$
denotes the option that is good exactly at the worlds in $K$:
\begin{center}
\begin{tabular}{c|c|c}
 & accept $g$ (credit) & accept $-g$ (debt)\\\hline
$P_1$ & $o_1$ & $o_{23}$\\
$P_2$ & $o_{13}$ & $o_2$\\
$P_3$ & $o_{12}$ & $o_3$
\end{tabular}
\end{center}
At world~$1$, accepting $g$ is \emph{credit} (rank $0$), silence has rank
$1$, and accepting $-g$ is \emph{debt} (rank $2$). On $H'$ a regular model is
an advice function with these three values. For a regular $M$ and a
measurable $Q\subseteq H'$, $M+\debt(Q)$, $M+\cred(Q)$ and $M+o(Q)$ agree
with $M$ except that they give debt, credit, or the option $o$ on $\pm Q$;
``$M$ silenced on $Q$'' gives silence there. On $P_2\cup P_3$ write
$\omega_2$ for the option good at world~$2$ and bad at world~$3$ ($o_2$ on
$P_2$, $o_{12}$ on $P_3$), and $\omega_3$ for the option good at $3$ and bad
at $2$ ($o_{13}$ on $P_2$, $o_3$ on $P_3$).

Coordinates are
\[
 x\coloneqq\frac{g_1}{|g_2|}\ \text{on }D_x\coloneqq P_1\cup P_2,\qquad
 y\coloneqq\frac{g_1}{|g_3|}\ \text{on }D_y\coloneqq P_1\cup P_3,\qquad
 t\coloneqq\frac{|g_2|}{|g_3|}\ \text{on }P_2\cup P_3,
\]
that is, $x=t_{12}$, $y=t_{13}$ and $t=t_{23}$ in IC's notation. The law of
$(x,y)$ under $\mu|_{P_1}$ has the same null sets as Lebesgue measure on
$(0,\infty)^2$ (ECS, Lemma~1). The laws of $x$ on $D_x$, of $y$ on $D_y$, of
$t$ on $P_2\cup P_3$, and the joint laws of $(x,t)$ on $P_2$ and of $(y,t)$
on $P_3$ have a.e.\ positive densities (IC, Section~1). Level sets of these
coordinates are null. For sets defined by coordinates we use the shorthand
$D_x\cap[a,b)\coloneqq D_x\cap\{a\le x<b\}$, and similarly for $y$ on $D_y$
and $t$ on $P_2\cup P_3$; on $P_2\cup P_3$ we also write $(a,b)$ for
$(P_2\cup P_3)\cap\{a<t<b\}$.

\paragraph{Symmetry.} Exchanging worlds $2$ and $3$ exchanges $P_2$ and
$P_3$, $x$ and $y$, $D_x$ and $D_y$, the edges $[e_1,e_2]$ and $[e_1,e_3]$,
and replaces $t$ by $1/t$. All hypotheses below are invariant. A statement
obtained in this way is called the \emph{mirror} of the original.

\subsection{Axioms and their first consequences}

Each $\preceq_i$ is a preorder on models (Definition~\ref{def:H});
smaller is better, and $X\prec_iY$ means that $X\preceq_iY$ and not
$Y\preceq_iX$. (G) at world~$i$ is guidance monotonicity: a weak pointwise
improvement in the rank $r_i$ gives $\preceq_i$, and an improvement on a
non-null set gives $\prec_i$; the second half is \emph{strict (G)}. (IB) is
independence under replacement of shared advice: replacing advice that
two models share on a measurable set of problems by other shared advice
leaves their comparison unchanged. (S) supplies coherent undominated models $C_p$,
$p\in\Delta$, competitive on $\supp p$, with closed sections
$A_i(X)\coloneqq\{p:C_p\preceq_iX\}$ for every model $X$ and every world $i$.
A model $C$ is \emph{competitive} on $S$ if for every model $X$ some $i\in S$
has $C\preceq_iX$. A model $M$ is \emph{dominated} if some model $X$ has
$X\preceq_iM$ at every world and $X\prec_iM$ at some world, and
\emph{undominated} otherwise. We use:
\begin{enumerate}[label=(K\arabic*),leftmargin=*,itemsep=2pt]
\item every model $X$ has a selected $C_q$ with $C_q\preceq_iX$ for
$i=1,2,3$ (CBOI, Lemma~1; its proof uses Brouwer's theorem and no
completeness); this is the KKM-type lemma of CBOI, and such a $C_q$ is
called a \emph{KKM dominator} of~$X$;
\item every undominated model is tied at every world with some selected
$C_p$ and is competitive on $\supp p$ (CBOI, Lemma~1);
\item the face lemma (CC, Lemma~7, as used in IC and CCF; see also FAED,
Section~1): the credal set of a coherent model competitive on $S$ lies in the
face $\Delta_S$, by strict (G); in particular the credal set of $C_q$ lies in
$\Delta_{\supp q}$, and $C_{e_i}=D_i$;
\item if $V$ is undominated, it is competitive on $\{1,2,3\}$ (IC,
Section~1; by (K2) and the face lemma applied to $V$, whose credal set is
$\Delta$).
\end{enumerate}
By (IB), whether $M\preceq_1M'$ holds, for regular models that differ exactly
on $\pm Q$, depends only on their advice on $Q$. As in IC, Section~1, a change
is \emph{accepted} at world~$1$ in a context if the model before the change
is $\preceq_1$ the model after it; an accepted change transfers to every
context that has the same advice on the change set. A \emph{chain} is a
finite sequence of accepted changes, each applied to a model that has the
base's advice on the change set. By transitivity the start of a chain is
$\preceq_1$ its end.

\subsection{Losses at worlds 2 and 3}

Following CBOI, for $k=2,3$ let $\nu_{k1},\nu_{k2},\nu_{k3}$ be finite
nonnegative finitely additive measures (\emph{charges}) on the measurable
subsets of $H$ with
\begin{equation}
 \nu_{kl}(A)=0\text{ if }\mu(A)=0,\qquad\nu_{kl}(A)>0\text{ if }\mu(A)>0,
 \label{eq:positive}
\end{equation}
and put $L_k(\Dscr)\coloneqq\sum_{l=1}^3\nu_{kl}\{g:r_k(g,\Dscr)\ge l\}$. The
one-sided bound is
\begin{equation}
 B\text{ coherent undominated},\ X\text{ regular},\ B\preceq_kX
 \ \Longrightarrow\ L_k(B)\le L_k(X).
 \tag{W$_k$}\label{eq:W}
\end{equation}

\begin{definition}\label{def:H}
Hypothesis (H) consists of: $n=3$; (G) at every world; $\preceq_1$ is a
preorder satisfying (IB); $\preceq_2,\preceq_3$ are preorders with losses
$L_2,L_3$ of the above form satisfying \eqref{eq:positive} and
\eqref{eq:W}; and (S). Hypothesis (H$_{\rm II}$) is (H) with charges of the
form $\nu_{kl}=c_{kl}\,\mu$, $c_{kl}$ integrable; then $L_k$ is an ordinary
additive loss with strictly ordered local penalties and an integrable
envelope. Hypothesis (H$_{\rm I}$) is (H$_{\rm II}$) with
$X\preceq_kY\iff L_k(X)\le L_k(Y)$ for $k=2,3$.
\end{definition}

Under (H$_{\rm I}$), (G) at worlds $2,3$ and \eqref{eq:W} are automatic, so
(H$_{\rm I}$) is the standing hypothesis (SH) of IC, Section~3, without the
assumption that $V$ is undominated.

\paragraph{Ranks at worlds 2 and 3.} On $D_x$ (where $g_2<0$) world~$2$ ranks
debt $0$, silence $1$ and credit $2$; on $P_3$ it ranks credit $0$, silence
$1$, debt $2$. On $D_y$ world~$3$ ranks debt $0$, silence $1$, credit $2$; on
$P_2$ it ranks credit $0$, silence $1$, debt $2$. For a regular $M$ put
$R^k_l(M)\coloneqq\{g\in H':r_k(g,M)\ge l\}$. Regular models have rank $0$ on
unanimous problems and never rank $3$, so
$L_k(M)=\nu_{k1}(R^k_1(M))+\nu_{k2}(R^k_2(M))$.

\begin{lemma}\label{lem:diff}
If regular $M,M'$ agree off a measurable $Q\subseteq H'$, then for $k=2,3$
\[
 L_k(M')-L_k(M)=\sum_{l=1}^2\bigl[\nu_{kl}(Q\cap R^k_l(M'))
 -\nu_{kl}(Q\cap R^k_l(M))\bigr].
\]
\end{lemma}
\begin{proof}
$R^k_l(M')$ is the disjoint union of $Q\cap R^k_l(M')$ and
$R^k_l(M)\setminus Q$; likewise for $M$. Apply finite additivity.
\end{proof}

Per unit of charge, the changes from silence are as follows (the reverse
change flips both signs, and a change between credit and debt is the sum of
two such changes):
\begin{equation}
\begin{tabular}{c|c|c}
 & silence $\to$ credit & silence $\to$ debt\\\hline
$P_1$ & $(+\nu_{22},\,+\nu_{32})$ & $(-\nu_{21},\,-\nu_{31})$\\
$P_2$ & $(+\nu_{22},\,-\nu_{31})$ & $(-\nu_{21},\,+\nu_{32})$\\
$P_3$ & $(-\nu_{21},\,+\nu_{32})$ & $(+\nu_{22},\,-\nu_{31})$
\end{tabular}
\label{eq:table}
\end{equation}
The entries are the changes of $(L_2,L_3)$. On $P_2\cup P_3$, silence
$\to\omega_2$ changes them by $(-\nu_{21},+\nu_{32})$, and silence
$\to\omega_3$ by $(+\nu_{22},-\nu_{31})$.

\begin{lemma}\label{lem:YH}
Assume (H). There are integrable functions $c_{kl}>0$ a.e.\ ($k=2,3$,
$l=1,2$) and measurable sets $Z_1\supseteq Z_2\supseteq\cdots$ with
$\mu(Z_m)\to0$ such that $\pi_{kl}\coloneqq\nu_{kl}-c_{kl}\,\mu$ is a
nonnegative charge vanishing on $\mu$-null sets, with $\pi_{kl}(H\setminus
Z_m)=0$ for all $m,k,l$. Under (H$_{\rm II}$) one can take $Z_m=\varnothing$.
\end{lemma}
\begin{proof}
CBOI, Lemma~3, gives the Yosida--Hewitt decomposition
$\nu_{kl}=c_{kl}\,\mu+\pi_{kl}$ with $c_{kl}\ge0$ integrable, $\pi_{kl}\ge0$
purely finitely additive, sets $Z'_m$ with $\mu(Z'_m)\to0$ and
$\pi_{kl}(H\setminus Z'_m)=0$ \cite{YH,PP}; and $c_{kl}>0$ a.e.\ by
\eqref{eq:positive}, as in the proof of CBOI, Lemma~3: if $c_{kl}=0$ on a
non-null $A$, then $\nu_{kl}(A\setminus Z'_m)=c_{kl}(A\setminus Z'_m)=0$ while
$\mu(A\setminus Z'_m)>0$ for large $m$. Since $0\le\pi_{kl}\le\nu_{kl}$, $\pi_{kl}$ vanishes on null
sets. Put $Z_m\coloneqq Z'_1\cap\dots\cap Z'_m$. Then $\mu(Z_m)\le\mu(Z'_m)$
and $\pi_{kl}(H\setminus Z_m)\le\sum_{i\le m}\pi_{kl}(H\setminus Z'_i)=0$.
\end{proof}

Fix these objects. Write
\[
 c_2\coloneqq c_{21},\quad v\coloneqq c_{22},\quad c_3\coloneqq c_{31},\quad
 u\coloneqq c_{32}.
\]
Under (H$_{\rm II}$) these are IC's and ECS's densities: $c_k$ is the silence
penalty and $v=b_2-c_2$, $u=b_3-c_3$. For an integrable $w\ge0$ and
measurable $A$ write $w(A)\coloneqq\int_Aw\,d\mu$. A measurable set $Q$ is
\emph{clean} if $Q\cap Z_m=\varnothing$ for some $m$; then
$\nu_{kl}(Q')=c_{kl}(Q')$ for all measurable $Q'\subseteq Q$ and all $k,l$.
Under (H$_{\rm II}$) every set is clean. Always
$\nu_{kl}(Q)\ge c_{kl}(Q)$. Since $\mu(Z_m)\to0$, $w(A\setminus Z_m)\to w(A)$
for every integrable $w\ge0$. Each $w\,d\mu$ is atomless, so for
$A_0\subseteq A$ every value in $[w(A_0),w(A)]$ is $w(A')$ for some
$A_0\subseteq A'\subseteq A$ \cite{Sierpinski}. Hence any finite list of
strict requirements on the masses $w(A')$ of a set $A'\subseteq A$ that can
be met inside $A$ can be met by a clean $A'$: work inside $A\setminus Z_m$
for large $m$. We say that such sets are \emph{chosen clean}.

\begin{remark}\label{rem:cleanside}
\emph{Clean-side rule.} By Lemma~\ref{lem:diff} and \eqref{eq:table}, the
change of $L_k$ under a change of advice on a finite union of disjoint sets
is a signed sum of charges of these sets. If every set that enters with the
sign $+$ is clean, then the change is at most the same signed sum with every
charge replaced by its density mass.
\end{remark}

\subsection{Forced comparisons}

\begin{lemma}\label{lem:forced}
Assume (H). Let $B$ be coherent and undominated, competitive on a set $S$ of
worlds with $1\in S$, and let $X$ be regular with $L_k(X)<L_k(B)$ for every
$k\in S\setminus\{1\}$. Then $B\preceq_1X$. If $B$ and $X$ differ exactly on
$\pm Q$, the change from $B$'s advice to $X$'s advice on $Q$ is accepted at
world~$1$ in every context having $B$'s advice on $Q$. If $B$ is competitive
on $\{2,3\}$, no regular $X$ has both $L_2(X)<L_2(B)$ and $L_3(X)<L_3(B)$.
\end{lemma}
\begin{proof}
By \eqref{eq:W} in contrapositive form, $B\not\preceq_kX$ for $k\in
S\setminus\{1\}$. Competitiveness gives some $i\in S$ with $B\preceq_iX$, so
$i=1$. The transfer is (IB) at world~$1$. The last statement is the same
argument with $S=\{2,3\}$.
\end{proof}

We call the last statement the \emph{Pareto property}. Lemma~\ref{lem:forced}
and the inequality $L_k(C)\le L_k(X)$ for a selected $C\preceq_kX$ (KKM
accounting) are the only ways in which the losses $L_2,L_3$ enter the proofs
below, apart from the results quoted from ECS and DT in
Section~\ref{sec:onesided}, the inputs from CCF in
Lemma~\ref{lem:edges}(a) and from EDL and FAED in Lemma~\ref{lem:FG}, and
Lemma~\ref{lem:otherworld}. These also use the losses only in the same two
forms: \eqref{eq:W} for a selected model, in contrapositive form or as
$L_k(C)\le L_k(X)$.

For measurable $B,C\subseteq H'$ define, as in CCF, Section~3,
\[
 B\preceq_dC\ :\Longleftrightarrow\ V+\debt(B)\preceq_1V+\debt(C),\quad
 B\preceq_cC\ :\Longleftrightarrow\ V+\cred(C)\preceq_1V+\cred(B).
\]
So $B\preceq_dC$ says that debt on $B$ is no worse than debt on $C$, and
$B\preceq_cC$ that credit on $B$ is worth no more than credit on $C$. Strict
relations are $\prec_d,\prec_c$.

\begin{lemma}\label{lem:F}
Assume (G) and (IB) at world~$1$.
\begin{enumerate}[label=(F\arabic*),leftmargin=*,itemsep=1pt]
\item $\preceq_d$ and $\preceq_c$ are preorders; chains may pass through
overlapping sets.
\item $B\subseteq C$ implies $B\preceq_dC$ and $B\preceq_cC$, strictly if
$C\setminus B$ is non-null.
\item $B\preceq_dC$ if and only if $B\setminus C\preceq_dC\setminus B$, and
likewise for $\prec_d$, $\preceq_c$, $\prec_c$.
\item If $B_i\preceq_dC_i$ for $i\le r$, where the $B_i$ are pairwise
disjoint, the $C_i$ are pairwise disjoint and $B_i\cap C_j=\varnothing$ for
all $i,j$, then $\bigcup_iB_i\preceq_d\bigcup_iC_i$. Likewise for
$\preceq_c$.
\item If $N$ is regular and silent on $B\cup C$, then $B\preceq_dC$ if and
only if $N+\debt(B)\preceq_1N+\debt(C)$; likewise for $\preceq_c$, $\prec_d$,
$\prec_c$.
\end{enumerate}
\end{lemma}
\begin{proof}
(F1) is transitivity of $\preceq_1$ in the single context $V$. (F2) is (G):
$V+\debt(C)$ is a pointwise world-1 worsening of $V+\debt(B)$, and
$V+\cred(C)$ an improvement of $V+\cred(B)$. (F3) and (F5) are (IB): the two
pairs differ on the same set, $B\triangle C$ in (F3), with the same advice
there. (F4): at step $k$ compare $V+\debt(E\cup B_k)$ and
$V+\debt(E\cup C_k)$ with
$E\coloneqq\bigcup_{i<k}C_i\cup\bigcup_{i>k}B_i$, which is disjoint from
$B_k\cup C_k$; by (F5) this comparison is $B_k\preceq_dC_k$. Transitivity
finishes.
\end{proof}

\subsection{Edge selections}

On $[e_1,e_2]$ the families of IC, Section~1, are: $E_1(a)$, $a\in[0,\infty]$,
with credit on $D_x\cap\{x>a\}$ and silence on the rest of $D_x$; $E_2(b)$,
$b\in[0,\infty]$, with debt on $D_x\cap\{x<b\}$ and silence on the rest of
$D_x$. Both give credit ($o_{12}$) on $P_3$. Then $E_1(0)=D_1$,
$E_2(\infty)=D_2$ and $E_1(\infty)=E_2(0)=V_{12}\coloneqq\Dscr_{[e_1,e_2]}$.
On $[e_1,e_3]$ the families $E_1'(a)$ and $E_3(b)$ are defined in the same
way with $y$ on $D_y$, and give credit ($o_{13}$) on $P_2$;
$V_{13}\coloneqq\Dscr_{[e_1,e_3]}$. In CCF's notation, $E_1(a)=E^2_1(a)$ and
$E_2(b)=E_2(b)$, and CCF parametrizes both families on $[e_1,e_2]$ by one
parameter: $M_s\coloneqq E_1(-1/s)$ for $s<0$, $M_{-\infty}\coloneqq D_1$,
$M_s\coloneqq E_2(s)$ for $s\ge0$. By strict (G) (CCF, (2)), $s<s'$ implies
$M_s\prec_1M_{s'}$ and $M_{s'}\prec_2M_s$.

\begin{lemma}\label{lem:edges}
Assume (H) and that $V$ is undominated.
\begin{enumerate}[label=(\alph*)]
\item Every selection $C_q$ with $q\in[e_1,e_2]$ has $e_1$ or $e_2$ in its
credal set; likewise on $[e_1,e_3]$.
\item There is a continuous map $s\colon[e_1,e_2]\to[-\infty,\infty]$ with
$C_q=M_{s(q)}$, $s(e_1)=-\infty$, $s(e_2)=\infty$. Every $M_s$ is selected;
for finite $s$ only at points of the open edge, so that $V_{12}$, $E_1(a)$
and $E_2(b)$, $a,b\in(0,\infty)$, are competitive on $\{1,2\}$.
\item If $s_n\to s^*$ in $[-\infty,\infty]$, with $s_n<s^*$ for all $n$ or
$s_n>s^*$ for all $n$, then, after passing to a subsequence of
$(s_n)$, there are indices $q_n\to q^*$ on the edge with
$C_{q_n}=M_{s_n}$ and $C_{q^*}=M_{s^*}$.
\item In the situation of (c), if $Y$ is regular and $M_{s^*}\not\preceq_1Y$,
then there is $n_0$ such that $M_s\not\preceq_1Y$ for every $s$ between
$s_{n_0}$ and $s^*$ (inclusive). The same holds with world~$2$ in place of
world~$1$.
\end{enumerate}
The mirror statements of (b)--(d) hold on $[e_1,e_3]$.
\end{lemma}
\begin{proof}
(a) Selections at $e_1,e_2$ are $D_1,D_2$ (K3). A selection at an open-edge
point is coherent, undominated and competitive on $\{1,2\}$, so CCF,
Proposition~17, applies; its hypotheses are (G), (S), (IB) at world~$1$, $V$
undominated, and losses at worlds $2,3$ with \eqref{eq:positive} and
\eqref{eq:W}, and it uses no completeness. (In the ordinary case this is IC,
Proposition~3; both rest on IC's chain lemma, IC, Lemma~2.) (b) is CCF, Lemma~3, which needs
(G) at worlds $1$ and $2$, (S), and (a). (c) Choose $q_n$ with
$s(q_n)=s_n$, by (b), and pass to a convergent subsequence (so (c) holds
along a subsequence of $(s_n)$, which suffices for every use below);
continuity gives $s(q^*)=s^*$. (d) CCF, Lemma~4, gives $q^*$ with $s(q^*)=s^*$, hence
$C_{q^*}=M_{s^*}$, such that every neighbourhood $N$ of $q^*$ has the property:
for some $n_0$ every $s$ between $s_{n_0}$ and $s^*$ is $s(q)$ for some $q\in
N$. Since $q^*\notin A_1(Y)$ and $A_1(Y)$ is closed, some neighbourhood $N$
misses $A_1(Y)$, and every $M_s$ with $s$ between $s_{n_0}$ and $s^*$ is
$C_q$ for some $q\in N$. For world~$2$ use $A_2(Y)$.
\end{proof}

\paragraph{Forced moves.} For the next lemma recall that $E_2(s)$ has debt on
$D_x\cap\{x<s\}$ and silence on $D_x\cap\{x\ge s\}$, and $E_1(s)$ credit on
$D_x\cap\{x>s\}$ and silence on $D_x\cap\{x\le s\}$.

\begin{lemma}\label{lem:moves}
Assume (H) and that $V$ is undominated. Let $s\in(0,\infty)$.
\begin{enumerate}[label=(M\arabic*),leftmargin=*,itemsep=1pt]
\item If $Q\subseteq D_x\cap\{x<s\}$, $W\subseteq D_x\cap\{x\ge s\}$ and
$\nu_{21}(Q)<\nu_{21}(W)$, then $Q\preceq_dW$. In particular this holds if
$Q$ is clean and $c_2(Q)<c_2(W)$.
\item The mirror: $Q\subseteq D_y\cap\{y<s\}$, $W\subseteq D_y\cap\{y\ge s\}$,
$\nu_{31}(Q)<\nu_{31}(W)$ imply $Q\preceq_dW$.
\item If $T\subseteq D_x\cap\{x>s\}$, $T'\subseteq D_x\cap\{x\le s\}$ and
$\nu_{22}(T')<\nu_{22}(T)$, then $T'\preceq_cT$. In particular this holds if
$T'$ is clean and $v(T')<v(T)$.
\item The mirror: $T\subseteq D_y\cap\{y>s\}$, $T'\subseteq D_y\cap\{y\le
s\}$, $\nu_{32}(T')<\nu_{32}(T)$ imply $T'\preceq_cT$.
\item If $A,T\subseteq D_x$ are disjoint and $\nu_{22}(T)<\nu_{21}(A)$, then
$V\preceq_1V+\debt(A)+\cred(T)$. In particular this holds if $T$ is clean
and $v(T)<c_2(A)$. The mirror holds on $D_y$ with $\nu_{32},\nu_{31}$.
\end{enumerate}
\end{lemma}
\begin{proof}
(M1) The base $E_2(s)$ is selected at an open-edge index, hence competitive
on $\{1,2\}$ (Lemma~\ref{lem:edges}). It has debt on $Q$ and is silent on $W$.
Let $X$ be $E_2(s)$ with silence on $Q$ and debt on $W$. On $D_x$, debt has
world-2 rank $0$ and silence rank $1$, so by Lemma~\ref{lem:diff}
$L_2(X)-L_2(E_2(s))=\nu_{21}(Q)-\nu_{21}(W)<0$. Lemma~\ref{lem:forced} gives
$E_2(s)\preceq_1X$, and (IB) transfers the change to the context $V$. The
clean case follows from $\nu_{21}(Q)=c_2(Q)$ and $\nu_{21}(W)\ge c_2(W)$.
(M3) The base $E_1(s)$ has credit on $T$ and is silent on $T'$; moving the
credit from $T$ to $T'$ changes $L_2$ by $\nu_{22}(T')-\nu_{22}(T)<0$, so
$V+\cred(T)\preceq_1V+\cred(T')$, that is, $T'\preceq_cT$. (M5) The base
$V_{12}$ is competitive on $\{1,2\}$ and silent on $D_x$; adding debt on $A$
and credit on $T$ changes $L_2$ by $-\nu_{21}(A)+\nu_{22}(T)<0$. $V_{12}$ and
$V$ differ only on $P_3$, so (IB) transfers the accepted change to $V$. (M2),
(M4) and the mirror of (M5) are the same with $E_3$, $E_1'$, $V_{13}$ and
$L_3$.
\end{proof}

In (M1)--(M4) the sets are disjoint because of the threshold. In (M1) and
(M2) only the \emph{source} $Q$ needs to be clean for the density version,
in (M3) and (M4) only the \emph{destination} $T'$, and in (M5) only the credit
set $T$. A \emph{debt chain} is a finite sequence of sets in which
consecutive sets are related by (M1) or (M2); its first set is $\preceq_d$ its
last. Each link is a statement in the context $V$, so only consecutive sets
need to be disjoint. Credit chains are defined in the same way with (M3),
(M4).

\subsection{Further inputs}

\begin{lemma}\label{lem:FG}
Assume (H) and that $V$ is undominated. There are nondecreasing functions
$F,G\colon(0,\infty)\to(0,\infty)$ with
\[
 c_2=F(t)\,u,\qquad v=G(t)\,c_3\qquad\text{a.e.\ on }P_2\cup P_3,
\]
and $d\coloneqq\sup F\le g_0\coloneqq\inf G$; $d$ and $g_0$ are positive and
finite.
\end{lemma}
\begin{proof}
Under (H$_{\rm II}$) this is EDL, Theorem~8, and its proof (Steps 2, 5
and~6), as recorded in IC, Section~3. Under (H) the same holds for the
densities $c_{kl}$: FAED, Proposition~6, proves Steps 1 and~3 of EDL's proof
for the ordinary parts, and Steps 2, 4, 5 and~6 use only the ordinary
penalties. EDL's $F$ and $G$ take values in $[0,\infty]$. Since $c_2/u$ and
$v/c_3$ are a.e.\ positive and finite, $\{F=0\}$ and $\{G=\infty\}$ are
intervals of Lebesgue measure zero at the ends of $(0,\infty)$, hence empty;
and $F\le\kappa^*\le G$ with $0<\kappa^*<\infty$, where $\kappa^*$ is
EDL's constant (Step~4 of the proof of EDL, Theorem~8; the inequality is
EDL's Step~5).
\end{proof}

We also use the following results as stated: IC, Lemma~14 (composition), which
uses only (IB) at world~$1$; IC, Lemma~11 (shape of a good base), which is a
statement about coherent models only; CBOI, Lemma~4 (compactness): the
coherent models are sequentially compact for Hausdorff convergence of credal
sets, along which advice converges a.e., and ordinary integral losses are
continuous along it.

\section{Preliminaries valid in all three cases}\label{sec:prelim}

Throughout Sections~\ref{sec:prelim}--\ref{sec:totals} we assume (H) and
that $V$ is undominated; the results are steps towards a contradiction. All
of them hold in the three cases (H$_{\rm I}$), (H$_{\rm II}$) and (H).

\subsection{Non-comparisons from closed sections}

\begin{lemma}\label{lem:noncomp}
Let $0<a<b<\infty$.
\begin{enumerate}[label=(\alph*)]
\item If $B_0\subseteq D_x\cap\{x<a\}$ is non-null, there is $S_0$ with
$B_0\not\preceq_dD_x\cap\{x\ge s\}$ for all $s\ge S_0$.
\item If $B_0\subseteq D_x\cap\{x<a\}$ is non-null, there is $\delta>0$ with
$B_0\not\preceq_dD_x\cap[s,b)$ for all $s\in(b-\delta,b)$.
\item If $Z\subseteq D_x\cap\{x\ge b\}$ is non-null, there is $\delta>0$ with
$V\not\preceq_1V+\debt(D_x\cap[s,b))+\cred(Z)$ for all $s\in(b-\delta,b)$.
\item If $T_0\subseteq D_x\cap\{x<a\}$ is non-null, there is $S_0$ with
$T_0\not\preceq_cD_x\cap\{x>s\}$ for all $s\ge S_0$.
\item If $T_0\subseteq D_x\cap\{x<a\}$ is non-null, there is $\delta>0$ with
$T_0\not\preceq_cD_x\cap\{s<x\le b\}$ for all $s\in(b-\delta,b)$.
\end{enumerate}
The mirror statements hold on $D_y$ with $y$, $E_3$, $E_1'$ and $V_{13}$.
\end{lemma}
\begin{proof}
In each case we name a regular $Y$ and a model $M_{s^*}$ with
$Y\prec_1M_{s^*}$ by strict (G), so that $M_{s^*}\not\preceq_1Y$, and apply
Lemma~\ref{lem:edges}(d) to an approach $M_s\to M_{s^*}$. Then we identify
$M_s\preceq_1Y$ with the displayed comparison by (IB).

(a) $Y\coloneqq D_2$ silenced on $B_0$, and $D_2=E_2(\infty)$ is approached by
$E_2(s)$, $s\to\infty$. For $s\ge a$, $E_2(s)$ and $Y$ differ exactly on
$B_0$ (debt, silence) and on $D_x\cap\{x\ge s\}$ (silence, debt). So
$E_2(s)\preceq_1Y$ if and only if $B_0\preceq_dD_x\cap\{x\ge s\}$.

(b) $Y\coloneqq E_2(b)$ silenced on $B_0$, approached by $E_2(s)$, $s\uparrow
b$. For $a\le s<b$ the two models differ exactly on $B_0$ (debt, silence)
and on $D_x\cap[s,b)$ (silence, debt).

(c) $Y\coloneqq E_2(b)+\cred(Z)$; $E_2(b)$ is silent on $Z$. Approach by
$E_2(s)$, $s\uparrow b$. The two models differ exactly on $D_x\cap[s,b)$
(silence, debt) and on $Z$ (silence, credit), as do $V$ and
$V+\debt(D_x\cap[s,b))+\cred(Z)$.

(d) $Y\coloneqq V_{12}+\cred(T_0)$, and $V_{12}=M_0$ is approached by
$E_1(s)=M_{-1/s}$, $s\to\infty$. For $s\ge a$, $E_1(s)$ and $Y$ differ
exactly on $D_x\cap\{x>s\}$ (credit, silence) and on $T_0$ (silence, credit).
So $E_1(s)\preceq_1Y$ if and only if
$V+\cred(D_x\cap\{x>s\})\preceq_1V+\cred(T_0)$, that is,
$T_0\preceq_cD_x\cap\{x>s\}$.

(e) $Y\coloneqq E_1(b)+\cred(T_0)$; $E_1(b)$ is silent on $T_0$. Approach by
$E_1(s)=M_{-1/s}$, $s\uparrow b$, so that the parameter increases to
$-1/b$. For $a\le s<b$, $E_1(s)$ and $Y$ differ exactly on $D_x\cap\{s<x\le
b\}$ (credit, silence) and on $T_0$ (silence, credit).
\end{proof}

\subsection{Closing, ratio criterion, and \texorpdfstring{$d<g_0$}{d<g0}}

For $b'\in(0,\infty)$ and $n>1/b'$ put
$\Band_n\coloneqq D_y\cap\{b'-1/n\le y<b'\}$, and let $\debt(\Band_n)$ be
debt there ($o_{23}$ on $P_1$, $o_3$ on $P_3$); this is the advice of
$E_3(b')$ there.

\begin{lemma}\label{lem:closing}
There are no $b'\in(0,\infty)$ and non-null $S\subseteq P_3\cap\{y\ge b'\}$
such that $V\preceq_1V+o_{12}(S)+\debt(\Band_n)$ for all large $n$. The mirror
statement holds on $[e_1,e_2]$, with $S\subseteq P_2\cap\{x\ge b'\}$, $o_{13}$
and the bands $D_x\cap\{b'-1/n\le x<b'\}$.
\end{lemma}
\begin{proof}
This is IC, Lemma~13. Its proof identifies the limit selection through
continuity and injectivity of $L_3$ on the edge, which needs exact
representation; we replace that step by the threshold map. Put
$X\coloneqq E_3(b')+o_{12}(S)$. For each $n$, $E_3(b'-1/n)$ and $X$ differ
exactly on $\Band_n$ (silence, debt) and on $S$ (silence, $o_{12}$), as do the
two models in the hypothesis. So by (IB) the hypothesis says
$E_3(b'-1/n)\preceq_1X$. By Lemma~\ref{lem:edges}(c), after passing to a
subsequence, there are indices $q_n\to q^*$ with $C_{q_n}=E_3(b'-1/n)$ and $C_{q^*}=E_3(b')$. Closedness of
$A_1(X)$ gives $E_3(b')\preceq_1X$. But $X$ improves $E_3(b')$ at world~$1$ by
credit on its silent set $S$, so $X\prec_1E_3(b')$ by strict (G). This is a
contradiction.
\end{proof}

\begin{lemma}\label{lem:ratio}
Let $S_0\subseteq P_3\cap\{y\ge1\}$ be clean and non-null, and let
$0<\ell_0<\infty$. Suppose that for every $R>0$ there is $s(R)>0$ with the
following property: for every $S'\subseteq S_0$ with $0<u(S')\le s(R)$ and
every $W\subseteq P_1\cap\{x\ge\ell_0\}$ with $c_2(W)\ge u(S')/R$, some
$D'\subseteq W$ has $c_2(D')\le u(S')/R$ and
$V\preceq_1V+o_{12}(S')+o_{23}(D')$. Then (H) and the undominatedness of $V$
are contradictory. The mirror statement holds with $S_0\subseteq P_2\cap\{x\ge
1\}$, $o_{13}$, $v$, $c_3$ and $W\subseteq P_1\cap\{y\ge\ell_0\}$.
\end{lemma}
\begin{proof}
This is IC, Lemma~15, whose statement contains a good base only to place
$S_0$; its proof does not use one. Take $b'=1$ and $S\coloneqq S_0$. Given
$n\ge2$, the set $W_n\coloneqq\Band_n\cap P_1\cap\{x\ge\ell_0\}$ has
$c_2(W_n)>0$. Choose $R$ with $u(S)/R<c_2(W_n)/2$ and split $S$ into disjoint
pieces $S_j$ with $u(S_j)\le s(R)$. Choose $D_j$ successively inside
$W_n\setminus(D_1\cup\dots\cup D_{j-1})$; this set has $c_2$-mass at least
$c_2(W_n)-\sum_iu(S_i)/R\ge c_2(W_n)/2\ge u(S_j)/R$. IC, Lemma~14, gives
$V\preceq_1V+o_{12}(S)+o_{23}(\bigcup_jD_j)$. Adding debt on the rest of
$\Band_n$ is a world-1 worsening, and $\Band_n\cap S=\varnothing$ since
$S\subseteq\{y\ge1\}$. So $V\preceq_1V+o_{12}(S)+\debt(\Band_n)$ for every
$n\ge2$, contrary to Lemma~\ref{lem:closing}. The mirror is the same on
$[e_1,e_2]$.
\end{proof}

\begin{proposition}\label{prop:dg}
$d<g_0$.
\end{proposition}
\begin{proof}
The chain below is run for small prize pieces and fed into
Lemma~\ref{lem:ratio}. (With a fixed prize set the auxiliary sets $T_1,T_2$
below would have to carry a fixed mass while the thresholds $t,t'$ escape to
$\infty$ and $0$, which is impossible when $F<d$ and $G>g_0$ everywhere.)
Suppose $d=g_0$ (Lemma~\ref{lem:FG} gives $d\le
g_0$). Then $F$ increases to $d$ at $\infty$ and $G$ decreases to $g_0$ at
$0+$. Fix $a'>1$, a clean non-null $S_0\subseteq P_3\cap\{1\le y\le
a'\}$, and $\ell>0$. Given $R>0$ choose $0<t'<t$ with
$\eta\coloneqq G(t'+)-F(t-)<1/(3R)$; note $\eta\ge0$. Put
\[
 s(R)\coloneqq\tfrac12\min\bigl\{u(P_3\cap\{t_{23}>t,\ y>a'\}),\
 c_3(P_2\cap\{t_{23}<t'\}),\ \tfrac R2\,c_2(P_1\cap\{x<\ell,\ y<\tfrac12\})
 \bigr\}.
\]
(Here $t_{23}$ is the coordinate and $t,t'$ are numbers.) Let $S'\subseteq
S_0$ with $s\coloneqq u(S')\le s(R)$, and $W\subseteq P_1\cap\{x\ge\ell\}$
with $c_2(W)\ge s/R$. Choose clean sets
\[
 T_1\subseteq P_3\cap\{t_{23}>t,\ y>a'\},\quad
 T_2\subseteq P_2\cap\{t_{23}<t'\},\quad
 Q\subseteq P_1\cap\{x<\ell,\ y<\tfrac12\},
\]
with $u(T_1)=c_3(T_2)=u_1\in(s,2s)$ and $c_2(Q)\in(\eta u_1,s/R)$; the
interval is nonempty because $\eta u_1<2s/(3R)$. Choose $D'\subseteq W$ with
$c_2(D')\in(c_2(Q),s/R]$. The chain:
\begin{enumerate}[label=\arabic*.,leftmargin=*,itemsep=2pt]
\item At $V$: add $o_{12}(T_1)$, $o_{13}(T_2)$, $o_{23}(Q)$. On $T_1$,
$c_2=F(t_{23})u\ge F(t-)u$; on $T_2$, $v=G(t_{23})c_3\le G(t'+)c_3$. By
\eqref{eq:table}, with $T_1,T_2$ clean,
\begin{gather*}
 \Delta L_2\le-F(t-)u_1+G(t'+)u_1-c_2(Q)=\eta u_1-c_2(Q)<0,\\
 \Delta L_3\le u(T_1)-c_3(T_2)-c_3(Q)=-c_3(Q)<0.
\end{gather*}
$V$ is competitive on $\{1,2,3\}$ (K4), so the change is accepted
(Lemma~\ref{lem:forced}).
\item Base $E_1'(a')$, with credit on $D_y\cap\{y>a'\}$ and silence on
$D_y\cap\{y\le a'\}$: move the credit from $T_1$ to $S'$. Then $\Delta
L_3=-\nu_{32}(T_1)+u(S')\le s-u_1<0$ ($S'$ is clean). The current model has
the base's advice on $T_1\cup S'$. This is (M4).
\item By (G): $T_2\colon o_{13}\to$ silence.
\item Base $E_2(\ell)$: move the debt from $Q$ to $D'$. This is (M1): $Q$ is
clean and $c_2(D')>c_2(Q)$.
\end{enumerate}
The end is $V+o_{12}(S')+o_{23}(D')$, with $c_2(D')\le u(S')/R$. All sets
involved are pairwise disjoint: $T_1$, $S'$ lie in $P_3$ and are separated
by $y$; $T_2\subseteq P_2$; $Q$ and $D'$ lie in $P_1$ and are separated by
$x$. So the hypothesis of Lemma~\ref{lem:ratio} holds with $\ell_0=\ell$,
which is a contradiction.
\end{proof}

\subsection{The edge \texorpdfstring{$[e_2,e_3]$}{[e2,e3]} and full-support selections}

\begin{lemma}\label{lem:noedge}
No selection $C_q$, at any index $q$, has a credal set contained in
$[e_2,e_3]$ that contains neither $e_2$ nor $e_3$. In particular every
selection on the edge $[e_2,e_3]$ has $e_2$ or $e_3$ in its credal set.
\end{lemma}
\begin{proof}
Let $C=\Dscr_K$ with $K\subseteq[e_2,e_3]$, $e_2,e_3\notin K$. Then every
$m\in K$ has $m_2,m_3>0$, and by compactness
$0<\tau_{\rm lo}\coloneqq\min_Km_3/m_2\le\tau_{\rm hi}\coloneqq\max_Km_3/m_2<\infty$.
On $P_2$, $m\cdot g=g_3(m_3-m_2t)$, so $C$ gives $o_2$ on $\{t>\tau_{\rm
hi}\}$; on $P_3$, $m\cdot g=|g_3|(m_2t-m_3)$, so $C$ gives $o_3$ on
$\{t<\tau_{\rm lo}\}$. Fix $\epsilon>0$ with $d(1+\epsilon)<g_0$
(Proposition~\ref{prop:dg}), choose a clean non-null $W_2\subseteq
P_3\cap\{t<\tau_{\rm lo}\}$, and then a clean $W_1\subseteq P_2\cap\{t>\tau_{\rm
hi}\}$ with $u(W_1)=(1+\epsilon)c_3(W_2)$. Let $X$ be $C$ silenced on
$W_1\cup W_2$. By \eqref{eq:table} and Lemma~\ref{lem:FG},
\begin{align*}
 L_2(X)-L_2(C)&=c_2(W_1)-\nu_{22}(W_2)\le d\,u(W_1)-g_0c_3(W_2)
 =\bigl(d(1+\epsilon)-g_0\bigr)c_3(W_2)<0,\\
 L_3(X)-L_3(C)&=-\nu_{32}(W_1)+c_3(W_2)\le-\epsilon\,c_3(W_2)<0.
\end{align*}
So $C\not\preceq_2X$ and $C\not\preceq_3X$ by \eqref{eq:W}. Also $X\prec_1C$
by strict (G), since $X$ removes debt. This contradicts competitiveness of
$C$ on $\supp q$.
\end{proof}

\begin{lemma}\label{lem:noboth}
No full-support selection gives both $o_{13}$ and $o_2$ on non-null parts of
$P_2$.
\end{lemma}
\begin{proof}
This is IC, Lemma~10. Its Case~2 used IC's tied-level lemma (Lemma~9), which
needs exact representation; here it cannot occur, by
Proposition~\ref{prop:dg}. We give the argument of Case~1 with clean sets.
Let $C$ be such a selection, $A_+\coloneqq\{C=o_{13}\}\cap P_2$ and
$A_-\coloneqq\{C=o_2\}\cap P_2$, both non-null. Choose clean non-null
$T'\subseteq A_+$ and $B'\subseteq A_-$. Then
$a\coloneqq\esssup_{B'}F(t)\le d<g_0\le b\coloneqq\essinf_{T'}G(t)$. Shrink
$B'$ (replacing it by a non-null subset) so that $\tau\coloneqq\essinf_{B'}x>0$
and
$c_2(B')<c_2(P_1\cap\{x<\tau\})$, and so that $u(B')$ is small enough for the
choices below. Fix $\alpha_0>0$ with $P_3\cap\{F(t)>\alpha_0\}$ non-null.

Take nested $T_\lambda\subseteq T'$, $\lambda\in[0,1]$, with
$c_3(T_\lambda)=\lambda u(B')$ (this needs $u(B')\le c_3(T')$), and
put $\chi(\lambda)\coloneqq v(T_\lambda)-c_2(B')$. It is continuous and
increasing, $\chi(0)<0$, and $\chi(\lambda)\ge(b\lambda-a)u(B')>0$ for
$\lambda>a/b$, because $v=Gc_3\ge bc_3$ on $T'$ and $c_2=Fu\le au$ on $B'$.
Let $\lambda_0\in(0,a/b]$ be its smallest zero; $a/b<1$. For $\lambda\in
(\lambda_0,1)$ close to $\lambda_0$, $\delta\coloneqq\chi(\lambda)>0$ is as
small as we like; put $T\coloneqq T_\lambda$. Choose $\eta>0$ small, a clean
$B\subseteq P_1\cap\{x<\tau\}$ with $c_2(B)=c_2(B')-\eta$, and a clean
$S\subseteq P_3\cap\{F(t)>\alpha_0\}$ with $c_2(S)=2(\delta+\eta)$. Then
$u(S)\le c_2(S)/\alpha_0<c_3(T)$ once $\delta,\eta$ are small, since
$c_3(T)\ge\lambda_0u(B')$. The chain:
\begin{itemize}[leftmargin=*,itemsep=2pt]
\item At $V$: add $o_{13}(T)$, $o_{23}(B)$, $o_{12}(S)$. With $T$ and $S$
clean, $\Delta L_2\le v(T)-c_2(B)-c_2(S)=-\delta-\eta<0$ and $\Delta
L_3\le-c_3(T)-c_3(B)+u(S)<0$. Accepted, by (K4) and Lemma~\ref{lem:forced}.
\item Base $E_2(\tau)$, with debt on $B\subseteq\{x<\tau\}$ and silence on
$B'\subseteq\{x\ge\tau\}$: move the debt from $B$ to $B'$. This is (M1), since
$B$ is clean and $c_2(B')>c_2(B)$.
\item Base $C$, which gives $o_{13}$ on $T$ and $o_2$ on $B'$: silence both.
With $T$ and $B'$ clean, $\Delta L_2\le-v(T)+c_2(B')=-\delta<0$ and $\Delta
L_3\le c_3(T)-u(B')=(\lambda-1)u(B')<0$. $C$ is competitive on $\{1,2,3\}$,
so the change is accepted.
\end{itemize}
The end is $V+o_{12}(S)$, a strict world-1 improvement of $V$. This
contradicts strict (G).
\end{proof}

\subsection{Separability and positivity}

\begin{theorem}\label{thm:sep}
There are nondecreasing functions
$f,h,\ftil,\htil\colon(0,\infty)\to(0,\infty)$ such that, a.e.\ on $P_1$,
\[
 \frac{c_3}{c_2}=\frac{f(x)}{h(y)},\qquad\frac uv=\frac{\ftil(x)}{\htil(y)}.
\]
\end{theorem}
\begin{proof}
Under (H$_{\rm I}$) this is ECS, Theorem~14. ECS uses its standing
hypothesis only through its Lemma~2 (moves along the edges). That lemma cites
IC, Lemma~7, for the selection of the edge models, which is
Lemma~\ref{lem:edges}(b) here, and uses the losses only to conclude
$E\not\preceq_kE'$ from $L_k(E')<L_k(E)$; so its four parts are (M1)--(M4)
for clean sources, resp.\ clean destinations. In the proofs of ECS's
Lemmas~4 and~5 every set that is moved, and every set to which advice is
moved, is chosen inside one of the sets of a cycle, and the contradiction is
strict (G) at world~$1$. Given a cycle $(X_i,Y_i)$ of negative weight,
replace each of its sets by its intersection with $H\setminus Z_m$, with $m$
so large that all of them are non-null. By the remark after ECS,
Definition~3, the result is a cycle of no larger weight, and now every chosen
set is clean. So ECS's Lemmas~4 and~5 hold for $\zeta_\rho=\log(c_3/c_2)$ and
$\zeta_\sigma=\log(u/v)$, and ECS, Proposition~13, which is pure measure
theory, gives the factorizations.
\end{proof}

By ECS, Remark~17, the one-sided limits and the multiplicative jump ratios
$f(a+)/f(a-)$, $h(a+)/h(a-)$ are determined by $c_3/c_2$. Put
\[
 \varphi\coloneqq f(x)\,c_2\ \text{on }D_x,\quad
 \varphi\coloneqq h(y)\,c_3\ \text{on }P_3;\qquad
 \psi\coloneqq\ftil(x)\,v\ \text{on }D_x,\quad
 \psi\coloneqq\htil(y)\,u\ \text{on }P_3.
\]
On $P_1$, $f(x)c_2=h(y)c_3$ and $\ftil(x)v=\htil(y)u$ a.e., so
$\varphi=h(y)c_3$ and $\psi=\htil(y)u$ a.e.\ on $D_y$ as well. For
measurable $Q\subseteq H'$ put $\Phi(Q)\coloneqq\int_Q\varphi\,d\mu$ and
$\Psi(Q)\coloneqq\int_Q\psi\,d\mu$, with values in $[0,\infty]$. Both
densities are a.e.\ positive on $H'$.

\begin{lemma}\label{lem:pos}
$f(0+)>0$ and $h(0+)>0$.
\end{lemma}
\begin{proof}
Suppose $h(0+)=0$. On
$P_1\cap\{\tfrac12<x<1\}$ we have $c_3/c_2\ge f(\tfrac12)/h(y)$, so for every
$R'>0$ the set $P_1\cap\{x<1,\ y<1,\ c_3>R'c_2\}$ is non-null. Fix a clean
non-null $S_0\subseteq P_3\cap\{y\ge1\}$; we verify the hypothesis of
Lemma~\ref{lem:ratio} with $\ell_0=1$. Given $R$, put $R'\coloneqq3R$ and
$s(R)\coloneqq\frac14c_3(P_1\cap\{x<1,\ y<1,\ c_3>R'c_2\})$. Let
$S'\subseteq S_0$ with $s\coloneqq u(S')\le s(R)$, and $W\subseteq
P_1\cap\{x\ge1\}$ with $c_2(W)\ge s/R$. Choose a clean $B\subseteq
P_1\cap\{x<1,\ y<1,\ c_3>R'c_2\}$ with $c_3(B)=2s$, which exists since
$2s$ is less than the $c_3$-mass of this set; then $c_2(B)\le2s/R'$.
\begin{enumerate}[label=\arabic*.,leftmargin=*,itemsep=1pt]
\item At $V$: add $o_{12}(S')$ and $o_{23}(B)$. With $S'$ clean, $\Delta
L_2\le-c_2(S')-c_2(B)<0$ and $\Delta L_3\le s-2s<0$. Accepted.
\item Base $E_2(1)$: move the debt from $B$ to $D'\subseteq W$ with
$c_2(D')=\frac43c_2(B)$; this is possible since $\frac43c_2(B)\le
\frac{8s}{9R}<c_2(W)$. It is (M1).
\end{enumerate}
The end is $V+o_{12}(S')+o_{23}(D')$ with $c_2(D')\le8s/(9R)\le s/R$.
Lemma~\ref{lem:ratio} gives a contradiction. If $f(0+)=0$, the mirror
argument (exchange worlds $2,3$; mirror of Lemma~\ref{lem:ratio}) applies.
\end{proof}

\section{Routing primaries}\label{sec:routing}

The two theorems of this section show that world~$1$ compares debt sets, and
credit sets, strictly according to $\Phi$ and $\Psi$, provided the
lighter set is bounded and clean; the heavier set is arbitrary. The proofs route advice through
$P_1$ by edge moves only.

For $b\in(0,\infty)$ put
\[
 R_b\coloneqq(P_2\cap\{x<b\})\cup(P_3\cap\{y<b\})\cup
 (P_1\cap(\{x<b\}\cup\{y<b\})),\qquad F_b\coloneqq H'\setminus R_b .
\]
So $F_b=(P_2\cap\{x\ge b\})\cup(P_3\cap\{y\ge b\})\cup(P_1\cap\{x\ge b,\
y\ge b\})$. A set is \emph{bounded} if it lies in some $R_b$. Since
$R_b\subseteq(D_x\cap\{x<b\})\cup(D_y\cap\{y<b\})$, and $\varphi\le f(b)c_2$,
$\psi\le\ftil(b)v$ on $D_x\cap\{x<b\}$, and $\varphi\le h(b)c_3$, $\psi\le
\htil(b)u$ on $D_y\cap\{y<b\}$, we have $\Phi(R_b),\Psi(R_b)<\infty$.

\begin{lemma}\label{lem:windows}
Let $w\colon(0,\infty)\to(0,\infty)$ be nondecreasing, $0<\alpha<\beta<\infty$,
$\epsilon>0$ and $\theta>0$. There are finitely many disjoint intervals
$U_l=[a_l,b_l)\subseteq[\alpha,\beta)$ with $w(b_l-)\le(1+\epsilon)w(a_l)$
(\emph{$\epsilon$-windows} of $w$) whose union covers $[\alpha,\beta)$ up to a
set of Lebesgue measure less than $\theta$.
\end{lemma}
\begin{proof}
$g\coloneqq\log w$ is nondecreasing and finite on $[\alpha,\beta]$, with
$\Gamma\coloneqq g(\beta)-g(\alpha)<\infty$. Put
$\eta\coloneqq\log(1+\epsilon)$. Only finitely many points (at most
$2\Gamma/\eta+2$, the endpoints $\alpha,\beta$ included) of
$[\alpha,\beta]$ carry a jump $g(z+)-g(z-)\ge\eta/2$. Delete relatively open
intervals of total length less than $\theta$ around them. What remains is a
finite union of closed intervals $[p,q]$, at every point of which
$g(z+)-g(z)<\eta/2$. On such an interval put $z_0\coloneqq p$ and
\[
 z_{k+1}\coloneqq\min\bigl(q,\ \inf\{z>z_k:g(z)\ge g(z_k)+\eta/2\}\bigr).
\]
Then $z_{k+1}>z_k$, because $g<g(z_k)+\eta/2$ on a right neighbourhood of
$z_k$. On $[z_k,z_{k+1})$ we have $g<g(z_k)+\eta/2$, so
$w(z_{k+1}-)\le e^{\eta/2}w(z_k)\le(1+\epsilon)w(z_k)$. If $z_{k+1}<q$ then
$g(z)\ge g(z_k)+\eta/2$ for all $z>z_{k+1}$, so $g(z_{k+2})\ge
g(z_k)+\eta/2$. Hence $z_k=q$ after at most $4\Gamma/\eta+2$ steps. The
intervals $[z_k,z_{k+1})$ are the windows on $[p,q]$.
\end{proof}

\begin{lemma}\label{lem:windowmoves}
Let $U=[a,b)$ be an $\epsilon$-window of $f$, and $a<s<b$. If
$Q\subseteq D_x\cap[a,s)$ is clean, $W\subseteq D_x\cap[s,b)$ and
$c_2(Q)<c_2(W)\le(1+\epsilon)c_2(Q)$, then $Q\preceq_dW$ and
$\Phi(W)\le(1+\epsilon)^2\Phi(Q)$. The same holds for windows of $h$ with
$D_y$, $y$ and $c_3$. If $U$ is an $\epsilon$-window of $\ftil$,
$T\subseteq D_x\cap\{s<x<b\}$, $T'\subseteq D_x\cap\{a\le x\le s\}$ is clean
and $v(T)/(1+\epsilon)\le v(T')<v(T)$, then $T'\preceq_cT$ and
$\Psi(T')\ge\Psi(T)/(1+\epsilon)^2$. The same holds for windows of $\htil$
with $D_y$, $y$ and $u$.
\end{lemma}
\begin{proof}
$Q\preceq_dW$ is (M1). On $D_x\cap\{x\in U\}$,
$f(a)c_2\le\varphi\le f(b-)c_2\le(1+\epsilon)f(a)c_2$, so
$\Phi(W)\le(1+\epsilon)f(a)c_2(W)\le(1+\epsilon)^2f(a)c_2(Q)\le
(1+\epsilon)^2\Phi(Q)$. For credit, $T'\preceq_cT$ is (M3), and
$\Psi(T')\ge\ftil(a)v(T')\ge\ftil(a)v(T)/(1+\epsilon)\ge
\Psi(T)/(1+\epsilon)^2$, since $\Psi(T)\le(1+\epsilon)\ftil(a)v(T)$. On $D_y$
use (M2), (M4) and $\varphi=h(y)c_3$, $\psi=\htil(y)u$.
\end{proof}

\begin{lemma}\label{lem:cells}
For nonempty open intervals $I,J\subseteq(0,\infty)$ the sets
$P_1\cap\{x\in I,\ y\in J\}$, $P_2\cap\{x\in I\}$ and $P_3\cap\{y\in J\}$ are
non-null. Hence, for all large $m$, their parts outside $Z_m$ have positive
$c_2$-, $c_3$-, $v$-, $u$-, $\Phi$- and $\Psi$-mass.
\end{lemma}
\begin{proof}
The coordinate laws have a.e.\ positive densities (Section~\ref{sec:setting}),
the six densities are a.e.\ positive, and $\mu(Z_m)\to0$.
\end{proof}

\begin{theorem}\label{thm:debtprimary}
Let $B,C\subseteq H'$ be disjoint, $B$ bounded and clean, and
$\Phi(B)<\Phi(C)$, where $\Phi(C)=\infty$ is allowed. Then $B\prec_dC$.
\end{theorem}
\begin{proof}
\emph{Step 0.} Choose a bounded $C_0\subseteq C$ with
$\Phi(B)<\Phi(C_0)<\infty$ and $C\setminus C_0$ non-null. This is possible:
$\Phi(C\cap R_b)\uparrow\Phi(C)$ as $b\to\infty$; if $\Phi(C)=\infty$ take
$C_0=C\cap R_b$ for large $b$, and otherwise remove from a suitable $C\cap
R_b$ a non-null piece of small $\Phi$-mass. By (F2), $C_0\prec_dC$, so it
suffices to prove $B\preceq_dC_0$. If $B$ is null this is (F2). Otherwise put
$\Omega\coloneqq\Phi(C_0)-\Phi(B)>0$ and fix $b_0$ with $B\cup C_0\subseteq
R_{b_0}$. Every set chosen below is chosen clean.

\emph{Step 1: windows and hubs.} Fix $\epsilon\in(0,1)$ with
$(1+\epsilon)^9\Phi(B)<\Phi(B)+\Omega/4$. By Lemma~\ref{lem:windows} on
$[b_0+1,b_0+2)$, fix an $\epsilon$-window $U_*=[\xi_a,\xi_b)$ of $f$ and an
$\epsilon$-window $U_H=[y_a,y_b)$ of $h$, both inside $[b_0+1,b_0+2)$, and
interior points $\xi_m\in U_*$, $y_m\in U_H$. Put
\[
 HL\coloneqq P_1\cap\{y_a\le y<y_m\},\qquad HU\coloneqq P_1\cap\{y_m\le
 y<y_b\}.
\]
Then choose $a_*,\theta>0$ small (Step~5) and, by Lemma~\ref{lem:windows},
finite families $\mathcal U_x$ of $\epsilon$-windows of $f$ and $\mathcal U_y$
of $\epsilon$-windows of $h$ in $[a_*,b_0)$, each covering $[a_*,b_0)$
up to Lebesgue measure less than $\theta$. Give each window $U=[a,b)$ a
slice width $\theta_U\in(0,(b-a)/2)$, also fixed in Step~5.

\emph{Step 2: classes.} Put $B_x\coloneqq B\cap(P_2\cup(P_1\cap\{x<b_0\}))$
and $B_y\coloneqq B\setminus B_x$. Since $B\subseteq R_{b_0}$, $B_x\subseteq
D_x\cap\{x<b_0\}$ and $B_y\subseteq D_y\cap\{y<b_0\}$. The \emph{source
pieces} are $B_x\cap\{a\le x<b-\theta_U\}$ for $U=[a,b)\in\mathcal U_x$ (type
X) and $B_y\cap\{a\le y<b-\theta_U\}$ for $U\in\mathcal U_y$ (type Y). The
rest of $B_x$ is the residue $R_x$, and the rest of $B_y$ is $R_y$: points
whose coordinate lies in $(0,a_*)$, in the uncovered set, or in a top slice
$[b-\theta_U,b)$. Define $C_{0,x}$, $C_{0,y}$ in the same way. The
\emph{target pieces} are $C_{0,x}\cap\{a+\theta_U\le x<b\}$ (type X$'$) and
$C_{0,y}\cap\{a+\theta_U\le y<b\}$ (type Y$'$); the rest of $C_0$ is the
target residue $R'$, which is not used.

\emph{Step 3: from a source to the hub.} Let $Q$ be a subset of a source
piece or of the residue. We route $Q$ by a debt chain to a set
$W_1\subseteq HL$, and record a bound $e(Q)\ge\Phi(W_1)$.
\begin{itemize}[leftmargin=*,itemsep=2pt]
\item Type X, window $[a,b)$: $W_1\subseteq HL\cap\{b-\theta_U\le x<b\}$
with $c_2(W_1)\in(c_2(Q),(1+\epsilon)c_2(Q)]$, by Lemma~\ref{lem:windowmoves}
at $s=b-\theta_U$. Put $e(Q)\coloneqq(1+\epsilon)^2\Phi(Q)$.
\item Type Y, window $[a,b)$ of $h$: first $W_0\subseteq P_1\cap\{b-\theta_U\le
y<b\}\cap\{\xi_a\le x<\xi_m\}$ with $c_3(W_0)\in(c_3(Q),(1+\epsilon)c_3(Q)]$,
by Lemma~\ref{lem:windowmoves} at $s=b-\theta_U$; then $W_1\subseteq
HL\cap\{\xi_m\le x<\xi_b\}$ with $c_2(W_1)\in(c_2(W_0),(1+\epsilon)c_2(W_0)]$,
by Lemma~\ref{lem:windowmoves} in $U_*$ at $s=\xi_m$. Put
$e(Q)\coloneqq(1+\epsilon)^4\Phi(Q)$.
\item $Q\subseteq R_x$: $W_1\subseteq HL\cap\{\xi_m\le x<\xi_b\}$ with
$c_2(W_1)\in(c_2(Q),(1+\epsilon)c_2(Q)]$, by (M1) at $E_2(\xi_m)$, since
$Q\subseteq\{x<b_0\}$. Then $\Phi(W_1)\le f(\xi_b)c_2(W_1)$; put
$e(Q)\coloneqq(1+\epsilon)f(\xi_b)c_2(Q)$.
\item $Q\subseteq R_y$: $W_1\subseteq HL$ with
$c_3(W_1)\in(c_3(Q),(1+\epsilon)c_3(Q)]$, by (M2) at $E_3(y_a)$, since
$Q\subseteq\{y<b_0\}$. Then $\Phi(W_1)\le h(y_m)c_3(W_1)$; put
$e(Q)\coloneqq(1+\epsilon)h(y_b)c_3(Q)$.
\end{itemize}
The first link uses only that $Q$ is clean, as a subset of $B$.

\emph{Step 4: from the hub to a target.} Let $W_1\subseteq HL$ and let $T'$
be a subset of a target piece.
\begin{itemize}[leftmargin=*,itemsep=2pt]
\item Type X$'$, window $[a,b)$: $W_2\subseteq HU\cap\{a\le x<a+\theta_U\}$
with $c_3(W_2)\in(c_3(W_1),(1+\epsilon)c_3(W_1)]$, by
Lemma~\ref{lem:windowmoves} in $U_H$ at $s=y_m$; then $W_2\to T'$ by (M1) at
$E_2(a+\theta_U)$.
\item Type Y$'$, window $[a,b)$ of $h$: $W_2\subseteq HU\cap\{\xi_a\le
x<\xi_m\}$ as before; $W_3\subseteq P_1\cap\{\xi_m\le x<\xi_b\}\cap\{a\le
y<a+\theta_U\}$ with $c_2(W_3)\in(c_2(W_2),(1+\epsilon)c_2(W_2)]$, by
Lemma~\ref{lem:windowmoves} in $U_*$ at $s=\xi_m$; then $W_3\to T'$ by (M2) at
$E_3(a+\theta_U)$.
\end{itemize}
The last intermediate $S$ ($W_2$ or $W_3$) and $T'$ lie in the same window
$[a,b)$ of $f$ (resp.\ $h$). The last link needs $c_2(T')>c_2(S)$ (resp.\
$c_3$), and $c_2(S)\le\Phi(S)/f(a)$, $c_2(T')\ge\Phi(T')/((1+\epsilon)f(a))$.
Each window move multiplies $\Phi$ by at most $(1+\epsilon)^2$. So the route
exists whenever $\Phi(T')>(1+\epsilon)^5\Phi(W_1)$, and in particular
whenever $\Phi(T')>(1+\epsilon)^5e(Q)$. The last link uses only that its
source is clean; $T'$ may be dirty.

\emph{Capacities.} All intermediate sets lie in finitely many fixed cells,
each with positive clean $c_2$- and $c_3$-mass (Lemma~\ref{lem:cells}). Their
prescribed masses are bounded by a constant, depending only on the windows,
times $e(Q)$; for instance $c_3(W_1)\le\Phi(W_1)/h(y_a)$ on $HL$. So every
route exists once $e(Q)$ is small, which we arrange by subdividing. Different
routes may use overlapping intermediate sets, since every link is a separate
statement in the context $V$.

\emph{Step 5: budget.} Summing over a partition of $B$ into the sets $Q$,
\[
 (1+\epsilon)^5\sum_Qe(Q)\le(1+\epsilon)^9\Phi(B)
 +(1+\epsilon)^6\bigl(f(\xi_b)c_2(R_x)+h(y_b)c_3(R_y)\bigr),
\]
and the target pieces have total $\Phi$-mass $\Phi(C_0)-\Phi(R')$. The
$x$-marginals of $c_2\,\mu$ and of $\varphi\,\mu$ on $D_x\cap\{x<b_0\}$, and
the $y$-marginals of $c_3\,\mu$ and $\varphi\,\mu$ on $D_y\cap\{y<b_0\}$, are
finite measures absolutely continuous with respect to Lebesgue measure. So
there is $\lambda>0$ such that sets of these coordinates of Lebesgue measure
less than $\lambda$ carry less than $\Omega/(8(1+\epsilon)^6(f(\xi_b)+
h(y_b)+1))$ of each. The residues $R_x,R_y,R'$ have coordinates in
$(0,a_*)$, in the uncovered sets or in slices. Choose $a_*+2\theta<
\lambda/2$, then the windows, then slice widths with $\sum_U\theta_U<
\lambda/2$. This order is not circular, because the bound holds for all sets
of small Lebesgue measure. Then the residue term above and $\Phi(R')$ are both
less than $\Omega/4$, and
\[
 (1+\epsilon)^5\sum_Qe(Q)<\Phi(B)+\tfrac{\Omega}{2}<\Phi(C_0)-\Phi(R').
\]

\emph{Matching.} Lay out the sets $Q$ as consecutive intervals of lengths
$(1+\epsilon)^5e(Q)$ on $[0,A]$, and the target pieces as consecutive
intervals of lengths $\Phi(T')$ on $[0,A']$; $A<A'$ by Step~5. Rescale the
second layout by $A/A'$ and take the common refinement of the two partitions
of $[0,A]$. Each interval $J$ of it lies in the interval of one $Q$ and in the
rescaled interval of one target piece. Since $e(Q)$ is a constant multiple of
$\Phi(Q)$, $c_2(Q)$ or $c_3(Q)$, and these measures are atomless, we can split
$Q$ into pieces $Q_J$ with $(1+\epsilon)^5e(Q_J)=|J|$, and the target piece
into pieces $T'_J$ with $\Phi(T'_J)=|J|A'/A>|J|$. Then
$\Phi(T'_J)>(1+\epsilon)^5e(Q_J)$, and Steps~3 and~4 give $Q_J\preceq_dT'_J$.

\emph{Step 6.} The $Q_J$ are pairwise disjoint and cover $B$, the $T'_J$ are
pairwise disjoint subsets of $C_0$, and $B\cap C_0=\varnothing$. By (F4) and
(F2), $B\preceq_d\bigcup_JT'_J\preceq_dC_0$.
\end{proof}

\begin{theorem}\label{thm:creditprimary}
Let $B,C\subseteq H'$ be disjoint, $B$ bounded and clean, and
$\Psi(B)<\Psi(C)$, where $\Psi(C)=\infty$ is allowed. Then $B\prec_cC$.
\end{theorem}
\begin{proof}
Credit moves \emph{down} in $x$ and $y$ ((M3), (M4)), and a credit chain from
$T$ to $T'$ gives $T'\preceq_cT$. So the chains now run from the heavier set
to the lighter set: the sources lie in $C$, and the lighter set $B$ is the
target, which must be covered completely.

\emph{Step 0.} As before, choose a bounded $C_0\subseteq C$ with
$\Psi(B)<\Psi(C_0)<\infty$ and $C\setminus C_0$ non-null, so $C_0\prec_cC$;
it suffices to prove $B\preceq_cC_0$. Assume $B$ non-null, put
$\Omega\coloneqq\Psi(C_0)-\Psi(B)$ and fix $b_0$ with $B\cup C_0\subseteq
R_{b_0}$. All intermediate sets are chosen clean.

\emph{Step 1.} Fix $\epsilon\in(0,1)$ with
$(1+\epsilon)^9\Psi(B)<\Psi(B)+\Omega/4$, an $\epsilon$-window
$U_*=[\xi_a,\xi_b)$ of $\ftil$ and an $\epsilon$-window $U_H=[y_a,y_b)$ of
$\htil$ inside $[b_0+1,b_0+2)$, with interior points $\xi_m$, $y_m$, and the
hubs
\[
 HU^c\coloneqq P_1\cap\{y_m<y<y_b\},\qquad HL^c\coloneqq P_1\cap\{y_a\le y\le
 y_m\}.
\]
Windows $\mathcal U_x$ of $\ftil$ and $\mathcal U_y$ of $\htil$ in
$[a_*,b_0)$, and slice widths $\theta_U$, are chosen as in
Theorem~\ref{thm:debtprimary}, with Step~5 below.

\emph{Step 2: classes.} $B_x,B_y$ and $C_{0,x},C_{0,y}$ are as in
Theorem~\ref{thm:debtprimary}. Source pieces: $C_{0,x}\cap\{a+\theta_U<x<b\}$
(type X) and $C_{0,y}\cap\{a+\theta_U<y<b\}$ (type Y), for windows $[a,b)$; the
rest of $C_0$ (coordinates in $(0,a_*)$, uncovered, or in a bottom slice
$[a,a+\theta_U]$) is discarded. Target pieces: $B_x\cap\{a\le x\le
b-\theta_U\}$ (type X$'$) and $B_y\cap\{a\le y\le b-\theta_U\}$ (type Y$'$);
the target residue $R'_x$ (resp.\ $R'_y$) consists of the points of $B_x$
(resp.\ $B_y$) whose coordinate lies in $(0,a_*)$, in the uncovered set,
or in a \emph{top} slice $(b-\theta_U,b)$.

\emph{Step 3: from a source to the hub.} Let $Q$ be a subset of a source
piece. It is routed by a credit chain to $S_1\subseteq HU^c$.
\begin{itemize}[leftmargin=*,itemsep=2pt]
\item Type X, window $[a,b)$ of $\ftil$: $S_1\subseteq HU^c\cap\{a\le x\le
a+\theta_U\}$ with $v(S_1)\in[v(Q)/(1+\epsilon),v(Q))$, by
Lemma~\ref{lem:windowmoves} (base $E_1(a+\theta_U)$). Then
$\Psi(S_1)\ge\Psi(Q)/(1+\epsilon)^2$.
\item Type Y, window $[a,b)$ of $\htil$: $S_0\subseteq P_1\cap\{a\le y\le
a+\theta_U\}\cap\{\xi_m<x<\xi_b\}$ with $u(S_0)\in[u(Q)/(1+\epsilon),u(Q))$
(base $E_1'(a+\theta_U)$); then $S_1\subseteq HU^c\cap\{\xi_a\le x\le\xi_m\}$
with $v(S_1)\in[v(S_0)/(1+\epsilon),v(S_0))$ (base $E_1(\xi_m)$, window
$U_*$). Then $\Psi(S_1)\ge\Psi(Q)/(1+\epsilon)^4$.
\end{itemize}
These links need only that the destinations are clean; $Q$ may be dirty.

\emph{Step 4: from the hub to a target.} Let $S_1\subseteq HU^c$ and let $T'$
be a subset of $B$ of one of the following kinds.
\begin{itemize}[leftmargin=*,itemsep=2pt]
\item Type X$'$, window $[a,b)$: $S_2\subseteq HL^c\cap\{b-\theta_U<x<b\}$
with $u(S_2)\in[u(S_1)/(1+\epsilon),u(S_1))$ (base $E_1'(y_m)$, window
$U_H$); then $S_2\to T'$ by (M3) at $E_1(b-\theta_U)$.
\item Type Y$'$, window $[a,b)$: $S_2\subseteq HL^c\cap\{\xi_m<x<\xi_b\}$
(base $E_1'(y_m)$); $S_3\subseteq P_1\cap\{\xi_a\le x\le\xi_m\}\cap
\{b-\theta_U<y<b\}$ with $v(S_3)\in[v(S_2)/(1+\epsilon),v(S_2))$ (base
$E_1(\xi_m)$, window $U_*$); then $S_3\to T'$ by (M4) at $E_1'(b-\theta_U)$.
\item $T'\subseteq R'_x$: $S_2\subseteq HL^c\cap\{\xi_m<x<\xi_b\}$ as in
type Y$'$; then $S_2\to T'$ by (M3) at $E_1(b_0)$, since
$T'\subseteq\{x<b_0\}$.
\item $T'\subseteq R'_y$: $S_1\to T'$ by (M4) at $E_1'(b_0)$, since
$T'\subseteq\{y<b_0\}$.
\end{itemize}
In the first two cases the last intermediate $S$ and $T'$ lie in the same
window of $\ftil$ (resp.\ $\htil$); the last link needs $v(T')<v(S)$
(resp.\ $u(T')<u(S)$), and
$v(S)\ge\Psi(S)/((1+\epsilon)\ftil(a))$, $v(T')\le\Psi(T')/\ftil(a)$
(resp.\ the same with $u$ and $\htil$). For
$R'_x$, $v(S_2)\ge\Psi(S_2)/\ftil(\xi_b)$; for $R'_y$,
$u(S_1)\ge\Psi(S_1)/\htil(y_b)$. Put
\[
 e(T')\coloneqq
 \begin{cases}
 (1+\epsilon)^5\Psi(T') & \text{(types X$'$, Y$'$)},\\
 (1+\epsilon)^2\ftil(\xi_b)v(T') & (T'\subseteq R'_x),\\
 \htil(y_b)u(T') & (T'\subseteq R'_y).
 \end{cases}
\]
Then the route from $S_1$ to $T'$ exists whenever $\Psi(S_1)>e(T')$, and the
whole route from $Q$ to $T'$ whenever $\Psi(Q)>(1+\epsilon)^4e(T')$. The last
link uses that $T'$ is clean, as a subset of $B$. Capacities are handled as
in Theorem~\ref{thm:debtprimary}.

\emph{Step 5: budget.} Summing over a partition of $B$ into sets $T'$,
\[
 (1+\epsilon)^4\sum_{T'}e(T')\le(1+\epsilon)^9\Psi(B)+(1+\epsilon)^6
 \bigl(\ftil(\xi_b)v(R'_x)+\htil(y_b)u(R'_y)\bigr),
\]
while the source pieces have total $\Psi$-mass $\Psi(C_0)$ minus the
discarded part. By absolute continuity of the coordinate marginals of $v\,\mu$,
$u\,\mu$ and $\psi\,\mu$ on the bounded region, as in
Theorem~\ref{thm:debtprimary}, choose $a_*,\theta$ and the slice widths so
that the residue term and the discarded mass are both less than $\Omega/4$.
Then $(1+\epsilon)^4\sum e(T')<\Psi(C_0)-\Psi(\text{discarded})$.

\emph{Matching} is as in Theorem~\ref{thm:debtprimary}, with the roles
exchanged: lay out the $T'$ with lengths $(1+\epsilon)^4e(T')$ on $[0,A]$ and
the source pieces with lengths $\Psi$ on $[0,A']$, $A<A'$; the common
refinement gives pairs $(Q_J,T'_J)$ with $\Psi(Q_J)>(1+\epsilon)^4e(T'_J)$,
where the $T'_J$ partition $B$ and the $Q_J$ are disjoint subsets of $C_0$.
Then $T'_J\preceq_cQ_J$, and (F4), (F2) give $B\preceq_c\bigcup_JQ_J\preceq_c
C_0$.
\end{proof}

\begin{remark}\label{rem:bounded}
The lighter side must be bounded in both theorems. Debt moves only up and
credit only down, and an unbounded far tail can be neither the start of an
upward chain nor the end of a downward one. Under (H) the lighter side must
also be clean: in the debt chains it is the first source, and in the credit
chains it receives the last destination. The heavier side is arbitrary.
\end{remark}

\begin{corollary}\label{cor:agree}
If $B\preceq_dC$ and $C$ is bounded and clean, then $\Phi(B)\le\Phi(C)$. If
$B\preceq_cC$ and $C$ is bounded and clean, then $\Psi(B)\le\Psi(C)$.
\end{corollary}
\begin{proof}
If $\Phi(B)>\Phi(C)$, then $\Phi(B\setminus C)>\Phi(C\setminus B)$, since
$\Phi(B\cap C)\le\Phi(C)<\infty$. Theorem~\ref{thm:debtprimary} applied to
the bounded clean set $C\setminus B$ gives $C\setminus B\prec_dB\setminus C$,
hence $C\prec_dB$ by (F3), contradicting $B\preceq_dC$. Credit is the same.
\end{proof}

\section{Totals and trade ratios}\label{sec:totals}

A \emph{pair} $(A,T)$ consists of disjoint, bounded, clean sets
$A,T\subseteq H'$. It is \emph{AC} if $V\preceq_1V+\debt(A)+\cred(T)$, and
\emph{WI} if $V+\debt(A)+\cred(T)\preceq_1V$. By (IB) these properties do not
change if $V$ is replaced by any regular model silent on $A\cup T$. Pairwise
disjoint AC pairs compose to an AC pair, by (IB) and transitivity, as in IC,
Lemma~14.

\begin{lemma}\label{lem:totals}
$\Phi(H')<\infty$ and $\Psi(H')<\infty$. Moreover
$\Phi(D_x\cap\{x\ge s\})\to0$ and $\Psi(D_x\cap\{x>s\})\to0$ as $s\to\infty$,
and likewise on $D_y$ in $y$.
\end{lemma}
\begin{proof}
Let $B_0\subseteq D_x\cap\{x<1\}$ be clean and non-null. By
Lemma~\ref{lem:noncomp}(a) there is $S_0\ge1$ with $B_0\not\preceq_d
D_x\cap\{x\ge s\}$ for $s\ge S_0$. If $\Phi(D_x\cap\{x\ge s\})>\Phi(B_0)$ for
such an $s$, monotone convergence gives $s'$ with $\Phi(D_x\cap[s,s'))>
\Phi(B_0)$, and Theorem~\ref{thm:debtprimary} and (F2) give
$B_0\prec_dD_x\cap[s,s')\preceq_dD_x\cap\{x\ge s\}$, a contradiction. So
$\Phi(D_x\cap\{x\ge s\})\le\Phi(B_0)$ for $s\ge S_0$. Since $B_0$ can have
arbitrarily small $\Phi$-mass, the tail tends to $0$. The same argument with
Lemma~\ref{lem:noncomp}(d) and Theorem~\ref{thm:creditprimary} (lighter side
$T_0\subseteq D_x\cap\{x<1\}$, clean) gives the credit tail, and the mirrors
give the tails on $D_y$. Finally $F_s\subseteq(D_x\cap\{x\ge s\})\cup(D_y
\cap\{y\ge s\})$ and $\Phi(R_s),\Psi(R_s)<\infty$.
\end{proof}

\begin{lemma}\label{lem:splus}
$s_+\coloneqq\inf\{\Phi(A)/\Psi(T):(A,T)\text{ an AC pair},\ \Psi(T)>0\}$ is
positive.
\end{lemma}
\begin{proof}
(a) Fix $b\in(0,\infty)$. Let $c>0$; choose a clean non-null bounded
$Z\subseteq D_x\cap\{x\ge b\}$ with $\Psi(Z)<c/3$. By
Lemma~\ref{lem:noncomp}(c) there is $\delta>0$ such that, for
$s\in(b-\delta,b)$ and $\Band\coloneqq D_x\cap[s,b)$, the pair $(\Band,Z)$
satisfies $V\not\preceq_1V+\debt(\Band)+\cred(Z)$. Fix such an $s$ with
$\Psi(\Band)<c/3$ and put $\eta\coloneqq\Phi(\Band)>0$. The order of choices
is $c$, $Z$, $\delta$, $s$, $\eta$.

(b) \emph{No AC pair $(A,T)$ has $\Phi(A)<\eta$ and $\Psi(T)\ge c$.} Let
$T_r\coloneqq T\setminus(\Band\cup Z)$, so $\Psi(T_r)\ge
\Psi(T)-\Psi(\Band)-\Psi(Z)>\Psi(Z)$. Starting from the AC relation
$V\preceq_1V+\debt(A)+\cred(T)$:
\begin{itemize}[leftmargin=*,itemsep=1pt]
\item drop the credit on $T\cap(\Band\cup Z)$: a world-1 worsening, by (G);
\item replace the debt on $A$ by debt on $\Band$: Theorem~\ref{thm:debtprimary}
gives $A\setminus\Band\prec_d\Band\setminus A$, since $A$ is clean and
bounded and $\Phi(A\setminus\Band)<\Phi(\Band\setminus A)$; transfer it by
(F5) to the context $\cred(T_r)+\debt(A\cap\Band)$, which is silent on
$A\triangle\Band$;
\item replace the credit on $T_r$ by credit on $Z$: Theorem~\ref{thm:creditprimary}
gives $Z\prec_cT_r$ ($Z$ clean and bounded); transfer it to the context
$\debt(\Band)$, which is silent on $Z\cup T_r$.
\end{itemize}
Each step is a weak worsening, so $V\preceq_1V+\debt(\Band)+\cred(Z)$,
contrary to (a).

(c) Suppose AC pairs $(A_n,T_n)$ with $\Psi(T_n)>0$ have
$r_n\coloneqq\Phi(A_n)/\Psi(T_n)\to0$. If $\Psi(T_n)\ge c>0$ along a
subsequence, then $\Phi(A_n)\le r_n\Psi(H')\to0$ by Lemma~\ref{lem:totals},
so $\Phi(A_n)<\eta$ eventually, contradicting (b). Otherwise
$\Psi(T_n)\to0$ along a subsequence. Note $A_n$ is non-null: otherwise
$V+\cred(T_n)\prec_1V$ would contradict the AC property. Put
$Q_T\coloneqq P_2\cap\{1\le x<2\}$, $Q_A\coloneqq P_1\cap\{x<1,\ y<1\}$ and
$c\coloneqq\Psi(Q_T)/4$. On the bounded set $Q_T$ the finite measures $\Phi$
and $\Psi$ are mutually absolutely continuous, so $\Psi(Q_T\cap A_n)\to0$.
Fix $\delta\in(0,\frac12)$. For large $n$ let $N_n\coloneqq\lceil
c/\Psi(T_n)\rceil$ and choose, successively for $k\le N_n$, pairwise
disjoint clean sets
\[
 T^*_k\subseteq Q_T\setminus A_n,\ \ \Psi(T^*_k)\in((1-\delta)\Psi(T_n),
 \Psi(T_n)),\qquad
 A^*_k\subseteq Q_A,\ \ \Phi(A^*_k)\in(\Phi(A_n),(1+\delta)\Phi(A_n)).
\]
There is room: the $T^*_k$ need $\Psi$-mass at most $c+\Psi(T_n)\le2c$, and
the $A^*_k$ need $\Phi$-mass at most $(1+\delta)(cr_n+\Phi(A_n))\to0$. Each
$(A^*_k,T^*_k)$ is AC:
\[
 V\preceq_1V+\debt(A_n)+\cred(T_n)\preceq_1V+\debt(A_n)+\cred(T^*_k)
 \preceq_1V+\debt(A^*_k)+\cred(T^*_k).
\]
The second step is $T^*_k\prec_cT_n$ (Theorem~\ref{thm:creditprimary} on
$T^*_k\setminus T_n$ and $T_n\setminus T^*_k$, with (F3)) in the context
$\debt(A_n)$, silent on $T_n\cup T^*_k$; the third is $A_n\prec_dA^*_k$
(Theorem~\ref{thm:debtprimary} on $A_n\setminus A^*_k$ and $A^*_k\setminus
A_n$) in the context $\cred(T^*_k)$, silent on $A_n\cup A^*_k$. The composite
of these disjoint AC pairs is an AC pair with credit mass at least
$(1-\delta)c$ and debt mass at most $(1+\delta)(cr_n+\Phi(A_n))\to0$. This
reduces to the first case, with $(1-\delta)c$ in place of $c$.
\end{proof}

\begin{corollary}\label{cor:kappaV}
$M_f\coloneqq\sup\ftil$ and $M_h\coloneqq\sup\htil$ are finite, and
\[
 s_+\le\kappa_V\coloneqq\min\{f(0+)/M_f,\ h(0+)/M_h\}.
\]
\end{corollary}
\begin{proof}
Let $0<\eta'<X$ and $\epsilon>0$. Choose a clean $T\subseteq
D_x\cap[X,X+1)$ with small $v(T)>0$ and a clean $A\subseteq D_x\cap\{x<\eta'\}$
with $c_2(A)=(1+\epsilon)v(T)$. By (M5), $(A,T)$ is AC. Its ratio is at most
$f(\eta')c_2(A)/(\ftil(X)v(T))=(1+\epsilon)f(\eta')/\ftil(X)$. Hence $s_+\le
f(0+)/\sup\ftil$, and Lemma~\ref{lem:splus} forces $\sup\ftil<\infty$ (and
again $f(0+)>0$). The mirror, with $V_{13}$, gives the bound with $h$.
\end{proof}

\begin{remark}\label{rem:ftil0}
Nothing is claimed about $\ftil(0+)$ or $\htil(0+)$. They can vanish: for
payoff-proportional penalties $\sigma$ is a constant times $x/y$, so
$\ftil(x)$ may be taken proportional to $x$. No argument below uses a lower
bound on $\ftil$ or $\htil$ near $0$.
\end{remark}

The next lemma bounds the ratios of WI pairs by the ratio of small forced AC
pairs. It is not needed for the main results: the proof of Theorem~\ref{thm:I}
uses the shorter credit cap of Lemma~\ref{lem:cap}, and the finitely additive
case uses Lemma~\ref{lem:Vcap}.

\begin{lemma}\label{lem:WI}
If $(A_0,T_0)$ is a WI pair, then $\Phi(A_0)\le\kappa_V\Psi(T_0)$.
\end{lemma}
\begin{proof}
Suppose $\Phi(A_0)>(\kappa_V+2\gamma)\Psi(T_0)$ with $\gamma>0$. If $T_0$ is
null, then $A_0$ is non-null and $V+\debt(A_0)\preceq_1V$ contradicts strict
(G); so $\Psi(T_0)>0$. Assume $\kappa_V=f(0+)/M_f$; otherwise use the mirror.
Fix $b$ with $A_0\cup T_0\subseteq R_b$ and put $m_1\coloneqq
\Psi(R_{b+1}\setminus R_b)>0$. Fix $\delta\in(0,\frac12)$ with
$\delta\Psi(T_0)<m_1/8$ and $(1+\delta)^4\kappa_V<\kappa_V+\gamma$, and then
$0<\eta'<X$ with
$(1+\delta)^4f(\eta')/\ftil(X)<\kappa_V+\gamma$ and $\Psi(N_a)<m_1/4$, where
$N_a\coloneqq D_x\cap\{x<\eta'\}$.

\emph{The forced pair and its relocated copies.} Choose a clean $t'\subseteq
D_x\cap[X,X+1)$ with $0<\Psi(t')<m_1/8$, small enough for the conditions
below, and a clean $a'\subseteq N_a$ with $c_2(a')=(1+\delta)v(t')$. By (M5),
$(a',t')$ is AC, and $\Phi(a')\le(1+\delta)f(\eta')\Psi(t')/\ftil(X)$. Let
$N$ be the least integer with $N(1-\delta)\Psi(t')>\Psi(T_0)$. Choose pairwise
disjoint clean sets
\begin{gather*}
 t^*_k\subseteq\bigl(T_0\cup(R_{b+1}\setminus R_b)\bigr)\setminus N_a,\qquad
 \Psi(t^*_k)\in((1-\delta)\Psi(t'),\Psi(t')),\\
 a^*_k\subseteq A_0,\qquad\Phi(a^*_k)\in(\Phi(a'),(1+\delta)\Phi(a')),
\end{gather*}
for $k\le N$. There is room. The $t^*_k$ need $\Psi$-mass at most
$N\Psi(t')\le(\Psi(T_0)+\Psi(t'))/(1-\delta)<\Psi(T_0)+m_1/2$, while the
ambient set has $\Psi$-mass at least $\Psi(T_0)+m_1-\Psi(N_a)$. The
$a^*_k$ need $\Phi$-mass at most
\[
 N(1+\delta)\Phi(a')\le\frac{(1+\delta)^2f(\eta')}{(1-\delta)\ftil(X)}
 \bigl(\Psi(T_0)+\Psi(t')\bigr)<(\kappa_V+\gamma)\bigl(\Psi(T_0)+\Psi(t')
 \bigr)<\Phi(A_0)
\]
for small $\Psi(t')$, using $1/(1-\delta)\le(1+\delta)^2$. Each
$(a^*_k,t^*_k)$ is AC:
\[
 V\preceq_1V+\debt(a')+\cred(t')\preceq_1V+\debt(a')+\cred(t^*_k)
 \preceq_1V+\debt(a^*_k)+\cred(t^*_k),
\]
by Theorem~\ref{thm:creditprimary} in the context $\debt(a')$ (silent on
$t'\cup t^*_k$, since $t^*_k\cap N_a=\varnothing$), and then
Theorem~\ref{thm:debtprimary} in the context $\cred(t^*_k)$ (silent on
$a'\cup a^*_k$, since $t^*_k\cap A_0=\varnothing$). The same forced pair
$(a',t')$ serves every copy, because each relation is established separately
in the context $V$.

\emph{Conclusion.} Let $a\coloneqq\bigcup_ka^*_k$ and
$t\coloneqq\bigcup_kt^*_k$. Then $(a,t)$ is AC, $a\subseteq A_0$ with
$A_0\setminus a$ non-null, $t\cap A_0=\varnothing$ and $\Psi(t)>\Psi(T_0)$.
Hence
\begin{align*}
 V+\debt(a)+\cred(t)&\prec_1V+\debt(A_0)+\cred(t)\prec_1V+\debt(A_0)+
 \cred(T_0)\\
 &\preceq_1V\preceq_1V+\debt(a)+\cred(t).
\end{align*}
The first step adds debt on the silent set $A_0\setminus a$ (strict (G)); the
second is $T_0\prec_ct$ (Theorem~\ref{thm:creditprimary}, lighter side
$T_0\setminus t$) in the context $\debt(A_0)$; the third is the WI property;
the last is the AC property. This is a contradiction.
\end{proof}

\section{Silent targets: the ordinary case}\label{sec:silent}

In this section we assume (H$_{\rm II}$) and that $V$ is undominated; every
set is clean. Let
$\ell_2\coloneqq c_2\one_{\{r_2\ge1\}}+v\one_{\{r_2\ge2\}}$ and
$\ell_3\coloneqq c_3\one_{\{r_3\ge1\}}+u\one_{\{r_3\ge2\}}$ be the local
penalties, so $L_k(M)=\int\ell_k(g,M(g))\,d\mu$ for regular $M$, and put
$\bar e_2\coloneqq c_2+v$, $\bar e_3\coloneqq c_3+u$, integrable bounds for $\ell_2,
\ell_3$ on $H'$.

Fix $r\coloneqq(d+g_0)/2\in(d,g_0)$ (Proposition~\ref{prop:dg}). For
$\kappa>0$ let $\lambda_\kappa$ be the auxiliary world-1 weight on $H'$:
$\lambda_\kappa=\varphi$ for debt, $0$ for silence and $-\kappa\psi$ for
credit. It is not claimed to represent $\preceq_1$; $\kappa$ is a free
parameter. With $M_f,M_h$ from Corollary~\ref{cor:kappaV} put
\[
 s_0(\kappa)\coloneqq\min\Bigl\{\frac{r-d}{\kappa M_h},\
 \frac{1-r/g_0}{\kappa M_f}\Bigr\}>0.
\]
For $0<s<s_0(\kappa)$ let $p\coloneqq(s,1,r)/(1+s+r)$, so $p_1=sp_2$ and
$p_3=rp_2$, and put
$\ell_p(g,a)\coloneqq p_1\lambda_\kappa(g,a)+p_2\ell_2(g,a)+p_3\ell_3(g,a)$.

\begin{lemma}\label{lem:silent}
Let $0<s<s_0(\kappa)$. Relative to silence, the costs $\ell_p$ of the
decisive options are, divided by $p_2$:
\begin{center}
\begin{tabular}{c|l|l}
 & credit & debt\\\hline
$P_1$ & $v(1-s\kappa\ftil)+ru$ & $c_2\,\bigl[sf-1-rf/h\bigr]$\\
$P_2$ & $c_3\,\bigl[G(1-s\kappa\ftil)-r\bigr]$ & $u\,\bigl[F(sf-1)+r\bigr]$\\
$P_3$ & $u\,\bigl[r-s\kappa\htil-F\bigr]$ & $c_3\,\bigl[sh+G-r\bigr]$
\end{tabular}
\end{center}
with $f=f(x)$, $\ftil=\ftil(x)$ on $D_x$, $h=h(y)$, $\htil=\htil(y)$ on
$D_y$, and $F=F(t)$, $G=G(t)$. All entries except the debt entry on $P_1$
are a.e.\ positive. Hence the a.e.-unique pointwise minimizer $\Sigma_p$ of
$\ell_p$ over regular advice gives debt on
\[
 \Theta_s\coloneqq\{g\in P_1:s<1/f(x)+r/h(y)\}
\]
and silence on the rest of $H'$; uniqueness can fail only on the level set
$\{g\in P_1:s=1/f(x)+r/h(y)\}$. $\Sigma_p$ never gives credit.
\end{lemma}
\begin{proof}
The entries follow from \eqref{eq:table} with densities, the world-1
increments $\varphi$ (debt) and $-\kappa\psi$ (credit), and the identities
$c_2=Fu$, $v=Gc_3$ on $P_2\cup P_3$, $c_3=(f/h)c_2$, $u=(\ftil/\htil)v$ on
$P_1$, $\varphi=fc_2=hc_3$ and $\psi=\ftil v=\htil u$. For example, debt on
$P_1$ costs $p_1fc_2-p_2c_2-p_3c_3=p_2c_2[sf-1-rf/h]$. Positivity uses
$F\le d$, $G\ge g_0$, $\ftil\le M_f$, $\htil\le M_h$ and $c_2,c_3,u,v>0$ a.e.:
\begin{itemize}[leftmargin=*,itemsep=1pt]
\item $P_1$ credit: $v(1-s\kappa M_f)+ru>0$, since $s\kappa M_f<1-r/g_0<1$;
\item $P_2$ credit: $1-s\kappa\ftil>0$, so $G(1-s\kappa\ftil)-r\ge
g_0(1-s\kappa M_f)-r>0$;
\item $P_2$ debt: $F(sf-1)+r\ge r-F\ge r-d>0$;
\item $P_3$ credit: $r-s\kappa\htil-F\ge r-d-s\kappa M_h>0$;
\item $P_3$ debt: $sh+G-r\ge g_0-r>0$.
\end{itemize}
The debt entry on $P_1$ is negative exactly on $\Theta_s$. On unanimous
problems regular advice is fixed.
\end{proof}

For a regular $M$ put $\delta_M(g)\coloneqq\ell_p(g,M(g))-\ell_p(g,\Sigma_p(g))$.
Then $\delta_M\ge0$ on $H'$, and $|\delta_M|\le p_1(\varphi+\kappa\psi)+p_2\bar e_2
+p_3\bar e_3$, which is integrable on every $R_b$.

\begin{lemma}\label{lem:corner}
Let $\kappa>0$. Suppose $f$ jumps at $x_1$ and $h$ jumps at $y_1$, and
$1/f(x_1+)+r/h(y_1+)<s_0(\kappa)$. Put $a_-\coloneqq1/f(x_1+)<a_+\coloneqq
1/f(x_1-)$, $b_-\coloneqq1/h(y_1+)<b_+\coloneqq1/h(y_1-)$ and
\[
 I\coloneqq\bigl(a_-+rb_-,\ \min\{a_++rb_-,\ a_-+rb_+,\ s_0(\kappa)\}\bigr)
 \ne\varnothing .
\]
For all but countably many $s\in I$, with $p$ as above and
$b_0\coloneqq2\max(x_1,y_1)+2$,
\[
 \eta\coloneqq\inf\Bigl\{\int_{R_{b_0}}\delta_C\,d\mu:C\text{ coherent}\Bigr\}>0 .
\]
\end{lemma}
\begin{proof}
(a) \emph{Coherent debt sets are convex in a chart.} On $P_1$ use the chart
$(\xi_1,\xi_2)\coloneqq(1/x,1/y)=(|g_2|/g_1,|g_3|/g_1)$. A coherent
$\Dscr_K$ gives debt at $g\in P_1$ iff $\max_Km\cdot g<0$, that is, iff
$m_2\xi_1+m_3\xi_2>m_1$ for all $m\in K$. Since $K$ is compact, this set
$O_K\subseteq(0,\infty)^2$ is open, and as an intersection of open half-planes
it is convex.

(b) \emph{The debt set of $\Sigma_p$ is not convex near the corner.} In the
chart, $\Theta_s$ is $U_s\coloneqq\{\alpha(\xi_1)+r\beta(\xi_2)>s\}$ with
$\alpha(\xi)\coloneqq1/f(1/\xi)$ and $\beta(\xi)\coloneqq1/h(1/\xi)$, both
nondecreasing. Put $X_1\coloneqq1/x_1$, $Y_1\coloneqq1/y_1$. Then
$\alpha\le a_-$ on $(0,X_1)$, $\alpha\ge a_+$ on $(X_1,\infty)$ and
$\alpha(\xi)\to a_-$ as $\xi\uparrow X_1$; likewise $\beta$ with $b_\pm$ and
$Y_1$. Since $s<a_++rb_-$ and $s<a_-+rb_+$, there is
$\delta\in(0,\min(X_1,Y_1)/4)$ with $a_++r\beta(Y_1-2\delta)>s$ and
$\alpha(X_1-2\delta)+rb_+>s$. Put $\epsilon\coloneqq\delta/4$ and
\[
 B_P\coloneqq(X_1+\delta\pm\epsilon)\times(Y_1-\delta\pm\epsilon),\qquad
 B_Q\coloneqq(X_1-\delta\pm\epsilon)\times(Y_1+\delta\pm\epsilon),
\]
where $(z\pm\epsilon)$ denotes $(z-\epsilon,z+\epsilon)$. On $B_P$,
$\alpha\ge a_+$ and $\beta\ge\beta(Y_1-2\delta)$, so $B_P\subseteq U_s$;
likewise $B_Q\subseteq U_s$. Their midpoint set
$B_M\coloneqq\frac12(B_P+B_Q)=(X_1\pm\epsilon)\times(Y_1\pm\epsilon)$ contains
$B_M^-\coloneqq(X_1-\epsilon,X_1)\times(Y_1-\epsilon,Y_1)$, on which
$\alpha+r\beta\le a_-+rb_-<s$; so $B_M^-\cap U_s=\varnothing$. Suppose an open
convex $O$ equals $U_s$ Lebesgue-a.e.\ on $B_P\cup B_Q\cup B_M$. Then $O\cap
B_P$ and $O\cap B_Q$ have full measure in $B_P$ and $B_Q$. For $z\in B_M$ the
set $O_z\coloneqq\{w\in B_P:2z-w\in B_Q\}$ is open and nonempty, hence of
positive measure. Both $O\cap O_z$ and $\{w\in O_z:2z-w\in O\}$ have full
measure in $O_z$, the second because $w\mapsto2z-w$ preserves Lebesgue
measure. So some $w$ has $w\in O$ and $2z-w\in O$, and by convexity $z\in O$.
Hence $B_M^-\subseteq O$, while $B_M^-$ misses $U_s$ and has positive
measure. This is a contradiction.

(c) \emph{Transfer to $\mu$.} The law of $(x,y)$ under $\mu|_{P_1}$ has the
same null sets as Lebesgue measure (ECS, Lemma~1), and so does that of
$(\xi_1,\xi_2)$, since $(x,y)\mapsto(1/x,1/y)$ is a diffeomorphism. So by
(a) and (b) no coherent model gives debt $\mu$-a.e.\ exactly on $\Theta_s$
within the region $P_1\cap(\xi_1,\xi_2)^{-1}(B_P\cup B_Q\cup B_M)$. In this
region $\xi_1>X_1-5\delta/4>\frac{11}{16}X_1$, so $x<2x_1<b_0$, and the region
lies in $R_{b_0}$.

(d) \emph{Strictness.} By Lemma~\ref{lem:silent}, $\delta_C(g)>0$ wherever
$C(g)\ne\Sigma_p(g)$, except on the level set $\{s=1/f(x)+r/h(y)\}$. These
level sets are disjoint for different $s$, so only countably many of them
have positive measure; exclude those $s$.

(e) \emph{Compactness.} If $\eta=0$, choose coherent $\Dscr_{K_n}$ with
$\int_{R_{b_0}}\delta_{\Dscr_{K_n}}\to0$, and a subsequence with $K_n\to K$
in Hausdorff distance. By CBOI, Lemma~4, the advice of $\Dscr_{K_n}$ converges
a.e.\ to that of $\Dscr_K$. By dominated convergence
$\int_{R_{b_0}}\delta_{\Dscr_K}=0$, so by (d) $\Dscr_K=\Sigma_p$ a.e.\ on
$R_{b_0}$, contradicting (c).
\end{proof}

A pair $(A,T)$ of disjoint bounded sets is \emph{SI} if
$V+\debt(A)+\cred(T)\prec_1V$.

\begin{lemma}\label{lem:kkmsilent}
Let $\kappa$, $s$, $p$, $b_0$ and $\eta$ be as in Lemma~\ref{lem:corner}.
There are an SI pair $(A_3,T)$ and $\theta>0$ with
$\Phi(A_3)>\kappa\Psi(T)+\theta$.
\end{lemma}
\begin{proof}
Put $\epsilon_k(b)\coloneqq \bar e_k(F_b)$; it tends to $0$ as $b\to\infty$. Fix
$b\ge b_0$ with $p_2\epsilon_2(b)+p_3\epsilon_3(b)<\eta/2$. Let $\Sigma^b$
be $\Sigma_p$ on $R_b$ and credit on $F_b$; it is regular. By (K1) there is a
selected $C$ with $C\preceq_i\Sigma^b$ for $i=1,2,3$.

(i) \emph{Accounting.} By \eqref{eq:W}, $L_k(C)\le L_k(\Sigma^b)$ for
$k=2,3$. Splitting the integrals over $R_b$ and $F_b$, and using $0\le\ell_k\le
\bar e_k$,
\[
 \int_{R_b}\bigl(\ell_k(C)-\ell_k(\Sigma_p)\bigr)\,d\mu
 \le\int_{F_b}\bigl(\ell_k(\cred)-\ell_k(C)\bigr)\,d\mu\le\epsilon_k(b).
\]
Since $C$ is coherent, $\delta_C\ge0$ and $b\ge b_0$,
$\int_{R_b}\delta_C\ge\eta$. As
$\delta_C=p_1(\lambda_\kappa(C)-\lambda_\kappa(\Sigma_p))+p_2(\ell_2(C)-
\ell_2(\Sigma_p))+p_3(\ell_3(C)-\ell_3(\Sigma_p))$,
\[
 p_1\int_{R_b}\bigl(\lambda_\kappa(C)-\lambda_\kappa(\Sigma_p)\bigr)\,d\mu
 \ge\eta-p_2\epsilon_2(b)-p_3\epsilon_3(b)>\eta/2 .
\]

(ii) \emph{World 1.} Let $\widetilde\Sigma$ be $\Sigma_p$ on $R_b$ and $C$ on
$F_b$. Credit has world-1 rank $0$, so $\widetilde\Sigma$ is a pointwise weak
world-1 worsening of $\Sigma^b$, and $C\preceq_1\Sigma^b\preceq_1
\widetilde\Sigma$ by (G). The set $D\coloneqq\{C\ne\widetilde\Sigma\}$ lies
in $R_b$, where $\widetilde\Sigma$ gives only debt (on $\Theta_s$) or
silence. Put
\[
 A\coloneqq D\cap\{C\text{ debt}\},\qquad A'\coloneqq D\cap\{\widetilde\Sigma
 \text{ debt}\},\qquad T\coloneqq D\cap\{C\text{ credit}\},
\]
and let $N$ be $\widetilde\Sigma$ silenced on $D$. Then
$\widetilde\Sigma=N+\debt(A')$ and $C=N+\debt(A)+\cred(T)$; $A$ is disjoint
from $A'\cup T$, while $A'\cap T$ may be non-null. Since $\Sigma_p$ gives no
credit, there is no cell where $\widetilde\Sigma$ gives credit, and
\[
 \int_{R_b}\bigl(\lambda_\kappa(C)-\lambda_\kappa(\Sigma_p)\bigr)\,d\mu
 =\Phi(A)-\Phi(A')-\kappa\Psi(T)\ge\frac\eta{2p_1}\eqqcolon2\theta .
\]

(iii) \emph{Reduction to a pure trade.} Since $\Phi$ is finite and atomless
on $R_b$, choose $A_2\subseteq A$ with $\Phi(A')<\Phi(A_2)<\Phi(A')+\theta$.
$A'$ is bounded and disjoint from $A_2$, so Theorem~\ref{thm:debtprimary}
gives $A'\prec_dA_2$, and (F5), with $N$ silent on $A'\cup A_2$, gives
$\widetilde\Sigma=N+\debt(A')\prec_1Y\coloneqq N+\debt(A_2)$. Put
$A_3\coloneqq A\setminus A_2$. $Y$ is silent on $A_3\cup T$, and
$C=Y+\debt(A_3)+\cred(T)$; on $A'\cap T$ both give credit, on $A'\setminus T$
both are silent. Hence $Y+\debt(A_3)+\cred(T)\preceq_1\widetilde\Sigma
\prec_1Y$, and (IB) gives $V+\debt(A_3)+\cred(T)\prec_1V$. Finally
$\Phi(A_3)=\Phi(A)-\Phi(A_2)>\kappa\Psi(T)+\theta$, and $A_3,T\subseteq R_b$.
\end{proof}

\begin{lemma}\label{lem:cap}
Let $Q_1\coloneqq P_2\cap\{\frac12\le x<1\}$, $Q_2\coloneqq P_2\cap\{1\le
x<2\}$, and $\omega(\delta)\coloneqq\sup\{c_2(E):E\subseteq Q_1,\ v(E)\le
\delta\}$, so $\omega(\delta)\to0$ as $\delta\to0$. Choose
$\kappa_*\ge2f(1)/\ftil(1)$ so large that $\tau\coloneqq\Phi(H')/\kappa_*$
satisfies $2\tau<\Psi(Q_2)$ and $\omega(\tau/\ftil(\frac12))+2\tau/\ftil(1)<
c_2(Q_1)$. Then no SI pair $(A,T)$ has $\Phi(A)>\kappa_*\Psi(T)$.
\end{lemma}
\begin{proof}
$\kappa_*$ exists by Lemma~\ref{lem:totals}. Suppose $(A,T)$ is SI with
$\Phi(A)>\kappa_*\Psi(T)$. Then $\Phi(A)\le\Phi(H')$, so $\Psi(T)<\tau$.
Also $T$ is non-null, since otherwise $V+\debt(A)\prec_1V$ with $A$
non-null, contrary to (G). Choose $T'\subseteq Q_2$ with
$\Psi(T)<\Psi(T')<2\Psi(T)$. Since $\psi=\ftil(x)v\ge\ftil(1)v$ on $Q_2$,
$v(T')<2\tau/\ftil(1)$. Since $\psi\ge\ftil(\frac12)v$ on $Q_1$,
$v(Q_1\cap T)<\tau/\ftil(\frac12)$, so $c_2(Q_1\setminus T)>c_2(Q_1)-
\omega(\tau/\ftil(\frac12))>v(T')$. Choose $a\subseteq Q_1\setminus T$ with
$v(T')<c_2(a)<v(T')+\gamma$, where
$\gamma\coloneqq\frac12\min\{c_2(Q_1\setminus T)-v(T'),\
(\Phi(A)-\kappa_*\Psi(T))/f(1)\}$. As $\varphi=f(x)c_2\le f(1)c_2$ on $Q_1$,
\[
 \Phi(a)\le f(1)c_2(a)<\frac{f(1)}{\ftil(1)}\Psi(T')+f(1)\gamma
 \le\kappa_*\Psi(T)+\tfrac12\bigl(\Phi(A)-\kappa_*\Psi(T)\bigr)<\Phi(A).
\]
\begin{enumerate}[label=\arabic*.,leftmargin=*,itemsep=1pt]
\item $V+\debt(a)+\cred(T)\prec_1V+\debt(A)+\cred(T)$: apply
Theorem~\ref{thm:debtprimary} to $a\setminus A$ and $A\setminus a$, and (IB)
in the context $\debt(a\cap A)+\cred(T)$; $T$ is disjoint from $a\cup A$. So
$(a,T)$ is SI.
\item $V+\debt(a)+\cred(T')\prec_1V+\debt(a)+\cred(T)$: apply
Theorem~\ref{thm:creditprimary} to $T\setminus T'$ and $T'\setminus T$, and
(IB) in the context $\cred(T\cap T')+\debt(a)$; $a$ is disjoint from
$T\cup T'$. So $(a,T')$ is SI.
\item $a,T'\subseteq D_x$ are disjoint and $c_2(a)>v(T')$, so (M5) gives
$V\preceq_1V+\debt(a)+\cred(T')$.
\end{enumerate}
Steps 3 and 2 give $V\prec_1V$, a contradiction.
\end{proof}

\begin{theorem}\label{thm:I}
Assume (H$_{\rm I}$): $n=3$; worlds $2,3$ are ranked exactly by ordinary
additive losses with strictly ordered local penalties and an integrable
envelope; world~$1$ is any preorder satisfying (G) and (IB); and (S) holds.
Then $V$ is dominated.
\end{theorem}
\begin{proof}
Suppose $V$ is undominated. All results of Sections~\ref{sec:prelim}
to~\ref{sec:silent} apply. By Theorem~\ref{thm:sep}, $c_3/c_2=f(x)/h(y)$ on
$P_1$ with $f,h$ positive and nondecreasing. If, for some $J<\infty$ and
$a_0<\infty$, $f(a+)/f(a-)\le J$ for all $a>a_0$, or the same holds for $h$,
then DT, Theorem~10, gives a contradiction. Otherwise, for every $J$ and $a_0$
there is $a>a_0$ with $f(a+)>Jf(a-)\ge Jf(1)$; so $f$ is unbounded and has
jumps beyond every point; the same holds for $h$.

Take $\kappa\coloneqq\kappa_*$ from Lemma~\ref{lem:cap}, $r=(d+g_0)/2$ and
$s_0\coloneqq s_0(\kappa_*)$. Choose a jump point $x_1$ of $f$ with
$f(x_1+)>2/s_0$ and a jump point $y_1$ of $h$ with $h(y_1+)>2r/s_0$; then
$1/f(x_1+)+r/h(y_1+)<s_0$. Lemma~\ref{lem:corner} gives $s$, $p$, $b_0$ and
$\eta>0$, and Lemma~\ref{lem:kkmsilent} gives an SI pair $(A_3,T)$ with
$\Phi(A_3)>\kappa_*\Psi(T)$. This contradicts Lemma~\ref{lem:cap}.

The order of choices is: the configuration (which determines
$f,h,\ftil,\htil,F,G,\Phi,\Psi$); then $\kappa_*$, which depends only on
$f(1)$, $\ftil(\frac12)$, $\ftil(1)$, $\Phi(H')$ and $c_2,v$ on $Q_1\cup Q_2$;
then $r$ and $s_0$; then $(x_1,y_1)$; then $s$, $p$, $b_0$, $\eta$; then $b$;
then $C$. There is no circularity.
\end{proof}

\begin{remark}\label{rem:split}
The case split is exhaustive. Its first branch contains the cases in which
$f$ or $h$ is continuous beyond some point, or bounded: a bounded factor has
jump ratios at most $\sup f/f(a_0)$ on $(a_0,\infty)$. In the second branch
both factors have unbounded jump ratios on every $(a_0,\infty)$, and
Lemma~\ref{lem:corner} needs only one large jump of each. Jump points and
jump ratios do not depend on the choice of the factors (ECS, Remark~17).
Completeness of world~$1$ is never used: positive world-1 comparisons come
only from Lemma~\ref{lem:forced}, (G), (IB), transitivity and
Theorems~\ref{thm:debtprimary}, \ref{thm:creditprimary}, which are themselves
chain consequences; closedness of $A_1$ at fixed models produces only
non-comparisons. The silent target has no credit on $R_b$, so the KKM
comparison reduces to a pure trade of debt for credit
(Lemma~\ref{lem:kkmsilent}(iii)), and $\kappa$ is free because non-coherence
of $\Sigma_p$ is shown directly (Lemma~\ref{lem:corner}).
\end{remark}

\begin{remark}\label{rem:A1fails}
In terms of the axioms, the residual branch of the proof shows the following.
Suppose (G) and (IB) hold at world~$1$, (G) and \eqref{eq:W} at worlds $2,3$
with ordinary additive losses $L_2,L_3$ as in (H$_{\rm II}$),
the selection consists of coherent undominated models competitive on their
supports, all sections $A_2(Y),A_3(Y)$ are closed, and $A_1(Y)$ is closed for
$Y=V$ and for the edge models and their perturbations used in
Sections~\ref{sec:prelim}--\ref{sec:totals} (in (K4) and
Lemmas~\ref{lem:edges}, \ref{lem:noncomp} and~\ref{lem:closing}). If $V$ is
undominated and $f,h$ have unbounded jump ratios, then $A_1(\Sigma^b_p)$ is
not closed for the silent targets $\Sigma^b_p$ of Lemma~\ref{lem:kkmsilent}:
otherwise (K1) would apply to $\Sigma^b_p$ and give a selected dominator,
which Lemmas~\ref{lem:kkmsilent} and~\ref{lem:cap} exclude. So the property of
(S) that fails first is closedness of the world-1 section at a silent target.
\end{remark}

\section{One-sided ordinary extensions}\label{sec:onesided}

\begin{theorem}\label{thm:II}
Assume (H$_{\rm II}$): $n=3$; (G) at every world; world~$1$ is any preorder
satisfying (IB); worlds $2,3$ are preorders that have ordinary additive
losses $L_2,L_3$ with strictly ordered local penalties and an integrable
envelope satisfying \eqref{eq:W}; and (S) holds. Then $V$ is dominated.
\end{theorem}
\begin{proof}
Sections~\ref{sec:prelim}--\ref{sec:totals} were proved under (H), and the
lemmas of Section~\ref{sec:silent} under (H$_{\rm II}$). In all of them the
losses $L_2,L_3$ enter only through Lemma~\ref{lem:forced} and through
$L_k(C)\le L_k(X)$ for selected $C\preceq_kX$, both of which are consequences
of \eqref{eq:W}. The proof of Theorem~\ref{thm:I} uses exact representation at
one point only: DT, Theorem~10, which is stated under exact representation.
So it suffices to check that DT, Theorem~10, holds under (H$_{\rm II}$) when
$V$ is undominated. Its proof uses DT, Theorem~2 (dense thresholds), DT,
Lemma~9 and the local debt cycle of DT, Lemma~5, and the facts (i)--(vi) of
DT, Section~1. We go through every use of worlds $2,3$.
\begin{enumerate}[label=(\arabic*),leftmargin=*,itemsep=2pt]
\item DT's fact~(iii) is IC, Lemma~7, whose proof derives the selection of the
edge models from continuity of $q\mapsto L_2(C_q)$ (EDL, Lemma~4(c)); that
needs $L_2$ to be a primary at world~$2$, that is, exactness. It is replaced
by Lemma~\ref{lem:edges}.
\item DT's fact~(vi) is IC, Lemma~13, whose proof identifies the limit
selection through continuity and injectivity of $L_3$ on the edge. It is
replaced by Lemma~\ref{lem:closing}.
\item The second sentence of DT's fact~(ii) is IC, Lemma~9 (tied levels),
which exhibits a dominator of $V$ through exact losses. It is used only in
DT, Lemma~1, to obtain $G>\kappa$ a.e.\ for $\kappa=F(\theta+)$. By
Proposition~\ref{prop:dg}, $G\ge g_0>d\ge\kappa$ everywhere, so it is not
needed. Proposition~\ref{prop:dg} is proved from forced comparisons only.
\item DT's fact~(iv): IC, Lemma~10, is Lemma~\ref{lem:noboth}. IC, Lemma~8, is
used in DT, Theorem~2, only to show that a selection with credal set in
$\{m_1=0\}$ containing neither $e_2$ nor $e_3$ does not exist; that is
Lemma~\ref{lem:noedge}. (IC's proof of Lemma~8(c) ends in IC, Lemma~9.)
\item In DT, Theorem~2, the KKM dominators $C_j\preceq_iX_j$ are selected, so
\eqref{eq:W} gives $L_k(C_j)\le L_k(X_j)$. This is all that the proof uses: if
$s>b_j$ or $s<a_j$, then $L_2(B_s)>L_2(X_j)$, resp.\ $L_3(B_s)>L_3(X_j)$, so
$C_j\ne B_s$; if $a_j\le s\le b_j$, then $X_j\prec_1B_s$ by strict (G),
while $C_j\preceq_1X_j$, so again $C_j\ne B_s$. The limit step uses only the
continuity of ordinary losses along Hausdorff-convergent coherent models
(CBOI, Lemma~4), and DT, Lemma~1, is a statement about the densities, applied
to the coherent limit.
\item In DT, Lemmas~5 and~9, every step of the chains is a move (M1) at
$E_2(\ell)$, a move (M2) at $E_3(\theta_0)$, a change at a good base or at
$V$ that decreases both $L_2$ and $L_3$ (Lemma~\ref{lem:forced} with
$S=\{1,2,3\}$), or a world-1 worsening by (G); the chains close by
Lemma~\ref{lem:closing}.
\item DT, Theorem~10, itself uses only the above and properties of $F$ and of
the good bases.
\end{enumerate}
No step derives a positive comparison at world~$2$ or~$3$ from the losses,
and no step exhibits a dominator of $V$ through the losses at worlds $2,3$:
every contradiction is with strict (G) at world~$1$ or with closedness. The
proof of Theorem~\ref{thm:I} then applies verbatim.
\end{proof}

\begin{remark}\label{rem:exactness}
In the working notes on which Sections~\ref{sec:routing}--\ref{sec:silent}
are based, exactness entered in four places: IC, Lemma~7, and IC, Lemma~13
(both replaced here from the start by Lemmas~\ref{lem:edges}
and~\ref{lem:closing}); IC, Lemma~9 (not used); and the continuity of
$q\mapsto L_k(C_q)$, which was used only for a reformulation like
Remark~\ref{rem:A1fails}. At worlds $2,3$ the proof uses transitivity (in
(K2)), strict (G) (in CCF, Lemma~3, and the face lemma), closedness of the
sections $A_2,A_3$ (in CCF, Lemma~3, inside (K1) and (K2), including
their use at $V$ for (K4), and at the coherent Bayes models of Step~1 of the
proof of EDL, Theorem~8, and of FAED, Proposition~6, which enter through
Lemma~\ref{lem:FG}) and \eqref{eq:W}; it uses neither
completeness nor (IB).
\end{remark}

\section{The finitely additive case: what carries over}\label{sec:faclean}

From now on we assume (H) and that $V$ is undominated. All results of
Sections~\ref{sec:prelim}--\ref{sec:totals} hold: they were proved under (H),
with clean sets wherever the charges matter. In particular we have the edge
selections (Lemma~\ref{lem:edges}), the forced moves in charge form
(Lemma~\ref{lem:moves}), $F$, $G$ and $d<g_0$ (Lemma~\ref{lem:FG},
Proposition~\ref{prop:dg}), Lemmas~\ref{lem:noedge} and~\ref{lem:noboth},
the factorizations $f,h,\ftil,\htil$ of the densities
(Theorem~\ref{thm:sep}), $f(0+),h(0+)>0$ (Lemma~\ref{lem:pos}), the routing
primaries with a clean lighter side (Theorems~\ref{thm:debtprimary},
\ref{thm:creditprimary}, Corollary~\ref{cor:agree}),
$\Phi(H'),\Psi(H')<\infty$ (Lemma~\ref{lem:totals}) and $s_+>0$ for clean AC
pairs (Lemma~\ref{lem:splus}). Here $\Phi,\Psi$ are the density weights.

Two steps of the ordinary proof do not carry over directly. First, the KKM
accounting of Lemma~\ref{lem:kkmsilent}(i) used $L_k(C)\le L_k(\Sigma^b)$ to
bound integrals of densities. Under (H), \eqref{eq:W} bounds only the sum of
the ordinary and the purely finitely additive parts, and a coherent $C$ can be
better than $\Sigma^b$ in the pure parts at worlds $2$ and $3$ at once, on
sets of small measure. Second, the limit step in DT, Theorem~2 (used through
DT, Theorem~10): pure parts are not continuous along Hausdorff-convergent
coherent models. Splicing the targets on the concentration sets $Z_m$, as in
CBOI's proof of its Theorem~2, does not repair this: it controls only the
weighted sum $L_2+\kappa L_3$, while DT's argument needs the two losses
separately, and in Lemma~\ref{lem:kkmsilent} it puts dirty sets into the
world-1 comparison. Moreover, Lemma~\ref{lem:cap} needs a clean credit set.
Sections~\ref{sec:germs}--\ref{sec:dense} show that the pure parts are tame
enough for fixed targets to work.

\paragraph{Germs.} Let $\nu$ be a charge on the measurable subsets of $D_x$,
vanishing on null sets. For $b\in(0,\infty)$ its \emph{right germ} at $b$ is
$\lim_{s\downarrow b}\nu(D_x\cap[b,s))$ and its \emph{left germ} is
$\lim_{s\uparrow b}\nu(D_x\cap[s,b))$; its \emph{germ at $0+$} is
$\lim_{s\downarrow0}\nu(D_x\cap\{x<s\})$, and its \emph{germ at $\infty$} is
$\lim_{s\to\infty}\nu(D_x\cap\{x\ge s\})$. The limits exist by monotonicity.
Since level sets of $x$ are null, it does not matter whether the endpoints
are included. The same terms are used on $D_y$ in $y$, and on $P_2\cup P_3$
in $t$, and for restrictions of $\nu$ to subsets.

\begin{lemma}\label{lem:slabs}
If $\nu$ has no germs on $D_x$, then for every $\epsilon>0$ there is a
finite partition of $(0,\infty)$ into intervals $J_1,\dots,J_N$ with
$\nu(D_x\cap\{x\in J_i\})<\epsilon$ for all $i$. More generally, if
$I\subseteq[0,\infty]$ is a compact interval and the germs of $\nu$ vanish at
the points of $I$ (at $0+$ if $0\in I$, at $\infty$ if $\infty\in I$), then
$I\cap(0,\infty)$ has such a finite partition.
\end{lemma}
\begin{proof}
Every point of $[0,\infty]$ (resp.\ of $I$) has a neighbourhood $N$
in $[0,\infty]$ with $\nu(D_x\cap\{x\in N\})<\epsilon$: for $b\in(0,\infty)$ the mass of $(b-\delta,b+\delta)$ is at
most the sum of the left and right masses near $b$, because $\{x=b\}$ is null
and hence $\nu$-null. Compactness of $[0,\infty]$ (resp.\ $I$) gives
finitely many such neighbourhoods; take the partition of $(0,\infty)$
(resp.\ $I\cap(0,\infty)$) they generate.
\end{proof}

\begin{lemma}\label{lem:tails}
(a) $\nu_{21}(D_x\cap\{x\ge s\})\to0$ and $\nu_{22}(D_x\cap\{x>s\})\to0$ as
$s\to\infty$. (b) For every $b\in(0,\infty)$, the left germs of $\nu_{21}$
and $\nu_{22}$ on $D_x$ at $b$ vanish. The mirror statements hold for
$\nu_{31},\nu_{32}$ on $D_y$ in $y$.
\end{lemma}
\begin{proof}
(a) Let $B_0\subseteq D_x\cap\{x<1\}$ be clean and non-null. By
Lemma~\ref{lem:noncomp}(a), $B_0\not\preceq_dD_x\cap\{x\ge s\}$ for $s\ge
S_0$. If $\nu_{21}(D_x\cap\{x\ge s\})>c_2(B_0)$ for some $s\ge\max(S_0,1)$,
(M1) at $E_2(1)$ (clean source $B_0$) gives the opposite. So
$\nu_{21}(D_x\cap\{x\ge s\})\le c_2(B_0)$ for large $s$, and $c_2(B_0)$ is
arbitrarily small. For $\nu_{22}$ take $T_0\subseteq D_x\cap\{x<1\}$ clean
and non-null; by Lemma~\ref{lem:noncomp}(d), $T_0\not\preceq_cD_x\cap\{x>s\}$
for large $s$, while (M3) at $E_1(1)$ (clean destination $T_0$) gives
$T_0\preceq_cD_x\cap\{x>s\}$ whenever $\nu_{22}(D_x\cap\{x>s\})>v(T_0)$.

(b) Let $0<a<b$ and $B_0\subseteq D_x\cap\{x<a\}$ clean and non-null. By
Lemma~\ref{lem:noncomp}(b), $B_0\not\preceq_dD_x\cap[s,b)$ for $s$ near $b$,
while (M1) at $E_2(s)$ gives $B_0\preceq_dD_x\cap[s,b)$ for $s>a$ whenever
$\nu_{21}(D_x\cap[s,b))>c_2(B_0)$. For $\nu_{22}$ use
Lemma~\ref{lem:noncomp}(e) with a clean $T_0\subseteq D_x\cap\{x<a\}$ and
(M3) at $E_1(s)$.
\end{proof}

\section{Germs on the edges \texorpdfstring{$[e_1,e_2]$}{[e1,e2]} and
\texorpdfstring{$[e_1,e_3]$}{[e1,e3]}}\label{sec:germs}

The non-comparisons of Lemma~\ref{lem:noncomp} come from closedness of
world-1 sections. They exclude left germs, but not right germs: approaching
$E_2(b)$ from the right, $E_2(s)$ is world-1 worse than $E_2(b)$, and nothing
follows. The missing tool is closedness of the section at the \emph{other}
world of the edge, fed by a world-1 non-comparison produced by forced moves.

\begin{lemma}\label{lem:otherworld}
Let $\{i,j\}$ be the support of an edge, let $q_n\to q^*$ with $q_n$ in the
open edge, and let $Y$ be regular with $C_{q_n}\not\preceq_iY$ for all $n$.
Then $C_{q^*}\preceq_jY$; if $j\in\{2,3\}$, then $L_j(C_{q^*})\le L_j(Y)$.
\end{lemma}
\begin{proof}
$C_{q_n}$ is competitive on $\{i,j\}$, so $C_{q_n}\preceq_jY$, that is,
$q_n\in A_j(Y)$. This set is closed, so $q^*\in A_j(Y)$. For $j\in\{2,3\}$
apply \eqref{eq:W} to the selected $C_{q^*}$.
\end{proof}

\begin{theorem}\label{thm:rightgerm}
For every $b\in(0,\infty)$ the right germ of $\nu_{21}$ on $D_x$ at $b$
vanishes.
\end{theorem}
\begin{proof}
Let $\delta_0$ be the right germ and $\nu(\epsilon)\coloneqq
\nu_{21}(D_x\cap[b,b+\epsilon))-\delta_0\ge0$; then $\nu(\epsilon)\to0$ as
$\epsilon\downarrow0$. Fix a small $\epsilon\in(0,1)$. Choose clean sets
\begin{gather*}
 Q'\subseteq D_x\cap[b+1,b+2),\qquad \nu(\epsilon)<c_2(Q')<\nu(\epsilon)+
 \epsilon,\\
 B_0\subseteq D_x\cap[\tfrac b2,b),\qquad c_2(B_0)=\frac{2f(b+2)
 (\nu(\epsilon)+\epsilon)}{f(b/2)};
\end{gather*}
this is possible for small $\epsilon$. Then $\Phi(Q')\le f(b+2)c_2(Q')<
\frac12f(\frac b2)c_2(B_0)\le\frac12\Phi(B_0)$.

For $s\in(b,b+\epsilon)$ put $\Band_s\coloneqq D_x\cap[s,b+\epsilon)$. Since
$\nu_{21}(D_x\cap[b,s))\ge\delta_0$,
\[
 \nu_{21}(\Band_s)=\nu_{21}(D_x\cap[b,b+\epsilon))-\nu_{21}(D_x\cap[b,s))
 \le\nu(\epsilon)<c_2(Q')=\nu_{21}(Q').
\]
(M1) at $E_2(b+\epsilon)$ gives $\Band_s\preceq_dQ'$; its source $\Band_s$
may be dirty, since (M1) compares charges. Theorem~\ref{thm:debtprimary},
with the clean bounded lighter side $Q'$, gives $Q'\prec_dB_0$. Hence
$\Band_s\prec_dB_0$, and so $B_0\not\preceq_d\Band_s$.

Let $Y$ be $E_2(b+\epsilon)$ silenced on $B_0$. For $s\in(b,b+\epsilon)$,
$E_2(s)$ and $Y$ differ exactly on $B_0$ (debt, silence) and on $\Band_s$
(silence, debt). By (IB), $E_2(s)\preceq_1Y$ if and only if
$B_0\preceq_d\Band_s$; so $E_2(s)\not\preceq_1Y$. By
Lemma~\ref{lem:edges}(c), with $s_n\downarrow b$ and after passing to a
subsequence, there are open-edge indices
$q_n\to q^*$ with $C_{q_n}=E_2(s_n)$ and $C_{q^*}=E_2(b)$.
Lemma~\ref{lem:otherworld} (with $i=1$, $j=2$) gives
$L_2(E_2(b))\le L_2(Y)$. By Lemma~\ref{lem:diff}, $Y$ and $E_2(b)$ differ on
$B_0$, where $Y$ raises the world-2 rank from $0$ to $1$, and on
$D_x\cap[b,b+\epsilon)$, where it lowers it from $1$ to $0$. Hence
\[
 \delta_0\le\nu_{21}(D_x\cap[b,b+\epsilon))\le c_2(B_0)
 =\frac{2f(b+2)(\nu(\epsilon)+\epsilon)}{f(b/2)}\longrightarrow0 .\qedhere
\]
\end{proof}

The model $Y$ carries the germ's world-2 value, because the germ lies in
$[b,s)$ for every $s>b$, while its extra debt $\Band_s$ misses the germ. This
is why the world-2 section, and not the world-1 section, does the work.

\begin{theorem}\label{thm:zerogerm}
The germ of $\nu_{21}$ on $D_x$ at $0+$ vanishes.
\end{theorem}
\begin{proof}
Let $\delta_0$ be the germ and $\nu(\epsilon)\coloneqq\nu_{21}(D_x\cap
\{x<\epsilon\})-\delta_0\to0$. For small $\epsilon\in(0,\frac12)$ choose clean
sets
\[
 Q'\subseteq D_x\cap[\epsilon,1),\ \ \nu(\epsilon)<c_2(Q')<\nu(\epsilon)+
 \epsilon,\qquad
 T_0\subseteq D_x\cap[1,2),\ \ \Psi(T_0)=\frac{2f(1)(\nu(\epsilon)+\epsilon)}
 {s_+},
\]
with $s_+>0$ from Lemma~\ref{lem:splus}. Then $\Phi(Q')<f(1)(\nu(\epsilon)+
\epsilon)$, so $\Phi(Q')/\Psi(T_0)<s_+/2$ and the clean pair $(Q',T_0)$ is
not AC. For $s\in(0,\epsilon)$ put $\Band_s\coloneqq D_x\cap[s,\epsilon)$.
As before $\nu_{21}(\Band_s)\le\nu(\epsilon)<c_2(Q')$, so (M1) at
$E_2(\epsilon)$ gives $\Band_s\preceq_dQ'$. If
$V\preceq_1V+\debt(\Band_s)+\cred(T_0)$, then (F5) in the context
$\cred(T_0)$, which is silent on $\Band_s\cup Q'$, gives
$V\preceq_1V+\debt(Q')+\cred(T_0)$, which is false. Let
$Y\coloneqq E_2(\epsilon)+\cred(T_0)$. For $s\in(0,\epsilon)$, $E_2(s)$ and
$Y$ differ exactly on $\Band_s$ (silence, debt) and $T_0$ (silence, credit),
so $E_2(s)\not\preceq_1Y$. With $s_n\downarrow0$, Lemma~\ref{lem:edges}(c)
gives, along a subsequence, $q_n\to q^*$ with $C_{q^*}=E_2(0)=V_{12}$, and
Lemma~\ref{lem:otherworld} gives $L_2(V_{12})\le L_2(Y)$. Now
$L_2(Y)-L_2(V_{12})=-\nu_{21}(D_x\cap\{x<\epsilon\})+v(T_0)$, and
$v(T_0)\le\Psi(T_0)/\ftil(1)$. So $\delta_0\le2f(1)(\nu(\epsilon)+\epsilon)/
(s_+\ftil(1))\to0$.
\end{proof}

\begin{theorem}\label{thm:rightgerm22}
For every $b\in(0,\infty)$ the right germ of $\nu_{22}$ on $D_x$ at $b$
vanishes.
\end{theorem}
\begin{proof}
Let $\delta_0$ be the germ and $\nu(\epsilon)\coloneqq\nu_{22}(D_x\cap
\{b<x\le b+\epsilon\})-\delta_0\to0$. Choose clean sets
\begin{gather*}
 T''\subseteq D_x\cap\{b+1<x\le b+2\},\qquad\nu(\epsilon)<v(T'')<
 \nu(\epsilon)+\epsilon,\\
 T_0\subseteq D_x\cap[\tfrac b2,b],\qquad\Psi(T_0)=2\ftil(b+2)(\nu(\epsilon)+
 \epsilon).
\end{gather*}
For $s\in(b,b+\epsilon)$ let $\Band'_s\coloneqq D_x\cap\{s<x\le b+\epsilon\}$;
then $\nu_{22}(\Band'_s)\le\nu(\epsilon)<v(T'')$. (M3) at $E_1(b+\epsilon)$,
moving credit from $T''$ down to $\Band'_s$, gives $\Band'_s\preceq_cT''$;
here the destination $\Band'_s$ may be dirty because (M3) compares charges.
Theorem~\ref{thm:creditprimary}, with the clean lighter side $T''$
($\Psi(T'')\le\ftil(b+2)v(T'')<\Psi(T_0)$), gives $T''\prec_cT_0$. Hence
$\Band'_s\prec_cT_0$, so $T_0\not\preceq_c\Band'_s$. Let
$Y\coloneqq E_1(b+\epsilon)+\cred(T_0)$. For $s\in(b,b+\epsilon)$, $E_1(s)$
and $Y$ differ exactly on $\Band'_s$ (credit, silence) and on $T_0$
(silence, credit), so $E_1(s)\preceq_1Y$ iff $T_0\preceq_c\Band'_s$, which is
false. Lemma~\ref{lem:edges}(c) ($E_1(s)\to E_1(b)$ as $s\downarrow b$,
along a subsequence) and
Lemma~\ref{lem:otherworld} give $L_2(E_1(b))\le L_2(Y)=L_2(E_1(b))-
\nu_{22}(D_x\cap\{b<x\le b+\epsilon\})+v(T_0)$, since credit has world-2 rank
$2$ on $D_x$. So $\delta_0\le v(T_0)\le\Psi(T_0)/\ftil(\frac b2)\to0$.
\end{proof}

\begin{corollary}\label{cor:mirrorgerms}
On $D_y$ in $y$: $\nu_{31}$ has no right germ at any $b\in(0,\infty)$ and no
germ at $0+$; $\nu_{32}$ has no right germ at any $b\in(0,\infty)$.
\end{corollary}
\begin{proof}
These are the mirrors of Theorems~\ref{thm:rightgerm}, \ref{thm:zerogerm}
and~\ref{thm:rightgerm22}, on the edge $[e_1,e_3]$ with $E_3$, $E_1'$,
$V_{13}$, (M2), (M4), and $s_+$ as before (its definition is symmetric). In
the analogue of Theorem~\ref{thm:zerogerm} one uses
$u(T_0)\le\Psi(T_0)/\htil(1)$.
\end{proof}

\begin{corollary}\label{cor:strong}
$\nu_{21}$ has no germs on $D_x$, and $\nu_{31}$ has none on $D_y$. The
charges $\nu_{22}$ on $D_x$ and $\nu_{32}$ on $D_y$ have no germs except
possibly at $0+$. Hence Lemma~\ref{lem:slabs} applies to $\nu_{21}$ on $D_x$
and to $\nu_{31}$ on $D_y$, and to $\nu_{22}$, $\nu_{32}$ on every set
$\{x\ge\eta\}$, resp.\ $\{y\ge\eta\}$, $\eta>0$.
\end{corollary}
\begin{proof}
Left germs: Lemma~\ref{lem:tails}(b); right germs:
Theorems~\ref{thm:rightgerm}, \ref{thm:rightgerm22} and
Corollary~\ref{cor:mirrorgerms}; $0+$: Theorem~\ref{thm:zerogerm} and
Corollary~\ref{cor:mirrorgerms}; $\infty$: Lemma~\ref{lem:tails}(a).
\end{proof}

\begin{theorem}\label{thm:crossgerm}
(a) For every $y_0\in[0,\infty)$, $\nu_{21}(P_1\cap\{|y-y_0|<\delta\})\to0$
as $\delta\downarrow0$ (for $y_0=0$ read $\{y<\delta\}$). (b) For every
$x_0\in[0,\infty)$, $\nu_{31}(P_1\cap\{|x-x_0|<\delta\})\to0$ as
$\delta\downarrow0$.
\end{theorem}
\begin{proof}
(a) Let $\epsilon>0$. By Theorem~\ref{thm:zerogerm} and
Lemma~\ref{lem:tails}(a) choose $0<a'<B<\infty$ with
$\nu_{21}(D_x\cap\{x<a'\})<\epsilon$ and $\nu_{21}(D_x\cap\{x\ge B\})<
\epsilon$. Put $W_\delta\coloneqq P_1\cap\{|y-y_0|<\delta\}\cap\{a'\le
x<B\}$ and $\mathrm{cap}\coloneqq\frac12c_2(P_2\cap[\frac{a'}2,a')\setminus
Z_m)$ for a large $m$; $\mathrm{cap}>0$. Let $e'\in(0,1)$. Choose a clean
$Q\subseteq P_2\cap[\frac{a'}2,a')$ with $c_2(Q)=(1-e')\min\{\nu_{21}(
W_\delta),\mathrm{cap}\}$; if $\nu_{21}(W_\delta)>0$, (M1) at $E_2(a')$ gives
$Q\preceq_dW_\delta$ (dirty destination). Choose a clean $Q'\subseteq
P_3\cap[y_0+\delta,y_0+1)$ with $\nu_{31}(W_\delta)<c_3(Q')<\nu_{31}(
W_\delta)+e'$; (M2) at $E_3(y_0+\delta)$ gives $W_\delta\preceq_dQ'$ (dirty
source). So $Q\preceq_dQ'$, and Corollary~\ref{cor:agree} ($Q'$ clean and
bounded) gives $\Phi(Q)\le\Phi(Q')$, that is,
\[
 f(\tfrac{a'}2)(1-e')\min\{\nu_{21}(W_\delta),\mathrm{cap}\}\le
 h(y_0+1)\bigl(\nu_{31}(W_\delta)+e'\bigr).
\]
Let $e'\to0$. As $\delta\to0$, $\nu_{31}(W_\delta)\le\nu_{31}(D_y\cap
\{|y-y_0|<\delta\})\to0$ by Corollary~\ref{cor:strong} (no germs of
$\nu_{31}$ at $y_0$; at $y_0=0$ the germ at $0+$). Hence
$\nu_{21}(W_\delta)\to0$, and $\nu_{21}(P_1\cap\{|y-y_0|<\delta\})<
2\epsilon+\nu_{21}(W_\delta)$. (b) is the mirror of (a).
\end{proof}

Consequently, for every compact interval $J\subseteq[0,\infty)$ and
$\epsilon>0$, $P_1\cap\{y\in J\}$ has a finite partition into $y$-slabs of
$\nu_{21}$-mass less than $\epsilon$, and $P_1\cap\{x\in J\}$ one into
$x$-slabs of $\nu_{31}$-mass less than $\epsilon$ (compactness, as in
Lemma~\ref{lem:slabs}).

\begin{remark}\label{rem:notcontrolled}
Two kinds of mass are not controlled by these results: $\nu_{21}$ on
$P_1\cap\{x\text{ bounded},\ y\to\infty\}$ and $\nu_{31}$ on
$P_1\cap\{y\text{ bounded},\ x\to\infty\}$ (cross mass at infinity, for which
only the one-sided bounds of Corollary~\ref{cor:onesidedinfty} are proved),
and the germs of $\nu_{22}$ at $x\to0+$ and of $\nu_{32}$ at $y\to0+$. The
approach $E_1(a)\to D_1$ ends at a vertex, and every fixed $Y$ with
$L_2(Y)<L_2(D_1)$ removes credit of $D_1$ on a fixed set, so no
non-comparison with all $E_1(a)$, $a$ small, is available. These masses are
harmless below: credit is only ever bounded from above
(Lemma~\ref{lem:Vcap}), and the cross mass at infinity enters only through
one-sided bounds.
\end{remark}

\section{The edge \texorpdfstring{$[e_2,e_3]$}{[e2,e3]}}\label{sec:edge23}

For $s\in[0,\infty]$ let $n_s\coloneqq(0,\frac1{1+s},\frac s{1+s})$ (with
$n_\infty=e_3$). Let $B_s$ be the coherent model with credal set
$\operatorname{conv}\{e_2,n_s\}$, and $B'_s$ the one with credal set
$\operatorname{conv}\{e_3,n_s\}$. For $m=(0,m_2,m_3)$ we have $m\cdot
g=g_3(m_3-m_2t)$ on $P_2$ and $m\cdot g=|g_3|(m_2t-m_3)$ on $P_3$. Hence both
give debt ($o_{23}$) on $P_1$; $B_s$ gives $\omega_2$ on $(s,\infty)$ and is
silent on $(0,s)$; $B'_s$ gives $\omega_3$ on $(0,s)$ and is silent on
$(s,\infty)$. Here and below $(a,b)$ denotes $(P_2\cup P_3)\cap\{a<t<b\}$.
Then $B_0=D_2$, $B_\infty=V_{23}\coloneqq\Dscr_{[e_2,e_3]}=B'_0$ and
$B'_\infty=D_3$. In DT's notation $B_\theta$ is the same model.

Index both families by $\sigma\in[0,2]$: $N_\sigma\coloneqq B_{\sigma/(1-
\sigma)}$ for $\sigma\in[0,1)$, $N_1\coloneqq V_{23}$, and
$N_\sigma\coloneqq B'_{(\sigma-1)/(2-\sigma)}$ for $\sigma\in(1,2]$. If
$\sigma<\sigma'$, then on a non-null $t$-band $N_\sigma$ gives $\omega_2$
where $N_{\sigma'}$ is silent, or is silent where $N_{\sigma'}$ gives
$\omega_3$, and they agree elsewhere. At world~$2$ the ranks of $\omega_2$,
silence and $\omega_3$ are $0,1,2$, and at world~$3$ they are $2,1,0$. By
strict (G),
\begin{equation}
 \sigma<\sigma'\ \Longrightarrow\ N_\sigma\prec_2N_{\sigma'}\quad\text{and}
 \quad N_{\sigma'}\prec_3N_\sigma .
 \label{eq:order23}
\end{equation}

\begin{lemma}\label{lem:sel23}
There is a continuous map $\sigma\colon[e_2,e_3]\to[0,2]$ with
$C_q=N_{\sigma(q)}$, $\sigma(e_2)=0$ and $\sigma(e_3)=2$. Every $N_\sigma$ is
selected; for $\sigma\in(0,2)$ only at points of the open edge, so that
$B_s$ ($0<s\le\infty$) and $B'_s$ ($0\le s<\infty$) are competitive on
$\{2,3\}$. Parts (c) and (d) of Lemma~\ref{lem:edges} hold for this family,
with worlds $2,3$ in place of $1,2$.
\end{lemma}
\begin{proof}
By the face lemma and Lemma~\ref{lem:noedge}, every selection on the edge is
some $B_s$ or $B'_s$, hence some $N_\sigma$. The proofs of CCF, Lemmas~3
and~4, apply verbatim with worlds $(2,3)$ in place of $(1,j)$ and
\eqref{eq:order23} in place of CCF's (2): they use only (G) at the two edge
worlds, the endpoint property, and closedness of the sections $A_2$, $A_3$ at
the fixed models $N_\sigma$. For $\sigma\in(0,2)$, $N_\sigma\notin\{D_2,D_3\}$,
so it is not selected at a vertex.
\end{proof}

On $P_2\cup P_3$ write
\[
 \alpha\coloneqq\nu_{21},\quad\beta\coloneqq\nu_{32},\quad
 \alpha'\coloneqq\nu_{22},\quad\beta'\coloneqq\nu_{31},
\]
restricted to subsets of $P_2\cup P_3$, with pure parts $\alpha^\pi,\beta^\pi,
\alpha'^\pi,\beta'^\pi$. By \eqref{eq:table}, silence $\to\omega_2$ changes
$(L_2,L_3)$ by $(-\alpha,+\beta)$ and silence $\to\omega_3$ by
$(+\alpha',-\beta')$. The densities are $c_2=Fu$, $u$, $v=Gc_3$ and $c_3$.

\begin{lemma}\label{lem:pareto}
Let $W\subseteq P_2\cup P_3$ be measurable (not necessarily clean), and let
$0<s'<s<\infty$. All compensators below are clean.
\begin{enumerate}[label=(P\arabic*),leftmargin=*,itemsep=2pt]
\item ($B_s$; add $\omega_2$ on $W\subseteq(0,s)$, remove it on
$W_1\subseteq(s,\infty)$.) If $u(W_1)>\beta(W)$, then $\alpha(W)\le
c_2(W_1)$.
\item ($B_s$; remove $\omega_2$ on $W\subseteq(s,\infty)$, add it on
$W_2\subseteq(0,s)$.) If $c_2(W_2)>\alpha(W)$, then $\beta(W)\le u(W_2)$;
if $W_2\subseteq(s',s)$, then $u(W_2)\le c_2(W_2)/F(s')$.
\item ($B'_s$; add $\omega_3$ on $W\subseteq(s,\infty)$, remove it on
$W_4\subseteq(0,s)$.) If $v(W_4)>\alpha'(W)$, then $\beta'(W)\le c_3(W_4)\le
v(W_4)/g_0$.
\item ($B'_s$; remove $\omega_3$ on $W\subseteq(0,s)$, add it on
$W_5\subseteq(s,\infty)$.) If $c_3(W_5)>\beta'(W)$, then $\alpha'(W)\le
v(W_5)$.
\item ($V_{23}$; add $\omega_2$ on $W$ and $\omega_3$ on $W_3$, disjoint
from $W$.) If $c_3(W_3)>\beta(W)$, then $\alpha(W)\le v(W_3)$.
\item ($V_{23}$; add $\omega_3$ on $W$ and $\omega_2$ on $W_6$, disjoint
from $W$.) If $c_2(W_6)>\alpha'(W)$, then $\beta'(W)\le u(W_6)$.
\item ($B_s$; add $\omega_3$ on $W\subseteq(0,s)$, or replace $\omega_2$ by
$\omega_3$ on $W\subseteq(s,\infty)$; add $\omega_2$ on $W_7\subseteq(s',s)$,
disjoint from $W$.) If $c_2(W_7)>\alpha'(W)$, resp.\
$c_2(W_7)>(\alpha+\alpha')(W)$, then $\beta'(W)\le c_2(W_7)/F(s')$, resp.\
$(\beta+\beta')(W)\le c_2(W_7)/F(s')$.
\end{enumerate}
\end{lemma}
\begin{proof}
In each case let $X$ be the named base, which is selected and competitive on
$\{2,3\}$ by Lemma~\ref{lem:sel23}, with the stated changes; the changes are
well defined because each compensator lies where the base has the opposite
advice. By Lemma~\ref{lem:diff} and the clean-side rule, $(L_2,L_3)(X)$ minus
the value at the base is
\begin{gather*}
 \text{(P1)}\ (-\alpha(W)+c_2(W_1),\ \beta(W)-u(W_1));\qquad
 \text{(P2)}\ (\alpha(W)-c_2(W_2),\ -\beta(W)+u(W_2));\\
 \text{(P3)}\ (\alpha'(W)-v(W_4),\ -\beta'(W)+c_3(W_4));\qquad
 \text{(P4)}\ (-\alpha'(W)+v(W_5),\ \beta'(W)-c_3(W_5));\\
 \text{(P5)}\ (-\alpha(W)+v(W_3),\ \beta(W)-c_3(W_3));\qquad
 \text{(P6)}\ (\alpha'(W)-c_2(W_6),\ -\beta'(W)+u(W_6));\\
 \text{(P7)}\ (\alpha'(W)[+\alpha(W)]-c_2(W_7),\ -\beta'(W)[-\beta(W)]+u(W_7)).
\end{gather*}
In each case the hypothesis makes one coordinate negative, and the Pareto
property (Lemma~\ref{lem:forced}) forces the other to be nonnegative. This is
the conclusion. The bounds $u\le c_2/F(s')$ on $(s',\infty)$ and $c_3\le v/g_0$
use $c_2=Fu$, $F$ nondecreasing, and $v=Gc_3$ with $G\ge g_0$.
\end{proof}

\begin{theorem}\label{thm:tgerms}
(a) $\alpha$ and $\beta$ have no left or right $t$-germs at any
$b\in(0,\infty)$; $\alpha$ has no germ at $0+$; $\alpha$ and $\beta$ have no
germ at $\infty$. (b) $\alpha'$ and $\beta'$ have no left or right $t$-germs
at any $b\in(0,\infty)$; $\alpha'$ and $\beta'$ have no germ at $0+$;
$\beta'$ has no germ at $\infty$.
\end{theorem}
\begin{proof}
We use Lemma~\ref{lem:otherworld} on $[e_2,e_3]$, with indices from
Lemma~\ref{lem:edges}(c) for the family $N_\sigma$
(Lemma~\ref{lem:sel23}), and the non-comparison supplied by \eqref{eq:W}:
\begin{itemize}[leftmargin=*,itemsep=1pt]
\item[(KF)] if $N_{\sigma_n}\to N_{\sigma^*}$ along indices $q_n\to q^*$
with $q_n$ in the open edge, and $Y$ is regular with
$L_3(Y)<L_3(N_{\sigma_n})$ for all $n$, then $L_2(N_{\sigma^*})\le L_2(Y)$;
and symmetrically with $2,3$ exchanged.
\end{itemize}
In each item $\gamma$ denotes the germ in question, and $W,W'$ are clean
sets of the indicated small mass. The identities for $L_k(B_s)$ follow from
Lemma~\ref{lem:diff}.

(a) \emph{Left germ of $\beta$ at $b$.} For $s<b$,
$L_3(B_s)=L_3(B_b)+\beta((s,b))\ge L_3(B_b)+\gamma$. If $\gamma>0$, let
$Y\coloneqq B_b+\omega_2(W')$ with $W'\subseteq(0,b)$, $u(W')<\gamma$. Then
$L_3(Y)<L_3(B_s)$ for all $s<b$, and (KF) gives $L_2(B_b)\le L_2(Y)=
L_2(B_b)-c_2(W')$, a contradiction.

\emph{Right germ of $\alpha$ at $b$.} For $s>b$,
$L_2(B_s)=L_2(B_b)+\alpha((b,s))$. If $\gamma>0$, let $Y$ be $B_b$ with
$\omega_2$ removed on $W\subseteq(b,\infty)$, $c_2(W)<\gamma$. Then
$L_2(Y)<L_2(B_s)$ for all $s>b$, and (KF) gives $L_3(B_b)\le L_3(Y)=L_3(B_b)-
u(W)$, a contradiction.

\emph{Left germ of $\alpha$ at $b$.} (P1) at $B_b$ with $W=(b-\eta,b)$ and
$W_1\subseteq(b,b+1)$, $u(W_1)=\beta(W)+\eta'$, gives $\alpha(W)\le
c_2(W_1)\le d(\beta(W)+\eta')$. As $\eta\to0$, $\beta(W)\to0$ by the first
item, and $\eta'>0$ is arbitrary.

\emph{Right germ of $\beta$ at $b$.} (P2) at $B_b$ with $W=(b,b+\eta)$ and
$W_2\subseteq(\frac b2,b)$, $c_2(W_2)=\alpha(W)+\eta'$, gives $\beta(W)\le
(\alpha(W)+\eta')/F(\frac b2)$, which tends to $0$ by the second item.

\emph{$\alpha$ at $0+$.} As $s\downarrow0$, $B_s\to B_0=D_2$, selected at the
vertex $e_2$. $L_2(B_s)=L_2(D_2)+\alpha((0,s))\ge L_2(D_2)+\gamma$. Let $Y$ be
$D_2$ with $\omega_2$ removed on $W$, $c_2(W)<\gamma$; (KF) gives
$L_3(D_2)\le L_3(Y)=L_3(D_2)-u(W)$.

\emph{$\beta$ at $\infty$.} As $s\to\infty$, $B_s\to V_{23}$.
$L_3(B_s)=L_3(V_{23})+\beta((s,\infty))$. Let $Y\coloneqq
V_{23}+\omega_2(W')$, $u(W')<\gamma$; (KF) gives $L_2(V_{23})\le
L_2(V_{23})-c_2(W')$.

\emph{$\alpha$ at $\infty$.} (P5) with $W=(S,\infty)$, $S>2$, and
$W_3\subseteq(1,2)$, $c_3(W_3)=\beta(W)+\eta'$, gives $\alpha(W)\le
v(W_3)\le G(2)(\beta(W)+\eta')$, which tends to $0$ as $S\to\infty$ by the
previous item.

(b) \emph{Left germ of $\beta'$ at $b$.} For $s<b$,
$L_3(B'_s)=L_3(B'_b)+\beta'((s,b))$. Let $Y$ be $B'_b$ with $\omega_3$
removed on $W\subseteq(0,b)$, $c_3(W)<\gamma$; (KF) gives $L_2(B'_b)\le
L_2(Y)=L_2(B'_b)-v(W)$.

\emph{Right germ of $\alpha'$ at $b$.} For $s>b$,
$L_2(B'_s)=L_2(B'_b)+\alpha'((b,s))$. Let $Y\coloneqq B'_b+\omega_3(W)$,
$W\subseteq(b,\infty)$, $v(W)<\gamma$; (KF) gives $L_3(B'_b)\le L_3(Y)=
L_3(B'_b)-c_3(W)$.

\emph{Left germ of $\alpha'$ at $b$.} (P4) at $B'_b$ with $W=(b-\eta,b)$,
$W_5\subseteq(b,b+1)$, $c_3(W_5)=\beta'(W)+\eta'$: $\alpha'(W)\le v(W_5)\le
G(b+1)(\beta'(W)+\eta')\to0$.

\emph{Right germ of $\beta'$ at $b$.} (P3) at $B'_b$ with $W=(b,b+\eta)$,
$W_4\subseteq(\frac b2,b)$, $v(W_4)=\alpha'(W)+\eta'$: $\beta'(W)\le
(\alpha'(W)+\eta')/g_0\to0$.

\emph{$\alpha'$ at $0+$.} As $s\downarrow0$, $B'_s\to B'_0=V_{23}$.
$L_2(B'_s)=L_2(V_{23})+\alpha'((0,s))$. Let $Y\coloneqq V_{23}+\omega_3(W)$,
$v(W)<\gamma$; (KF) gives $L_3(V_{23})\le L_3(V_{23})-c_3(W)$.

\emph{$\beta'$ at $0+$.} (P6) with $W=(0,\eta)$, $W_6\subseteq(1,2)$,
$c_2(W_6)=\alpha'(W)+\eta'$: $\beta'(W)\le u(W_6)\le(\alpha'(W)+\eta')/F(1)
\to0$.

\emph{$\beta'$ at $\infty$.} As $s\to\infty$, $B'_s\to D_3$, selected at the
vertex $e_3$. $L_3(B'_s)=L_3(D_3)+\beta'((s,\infty))$. Let $Y$ be $D_3$ with
$\omega_3$ removed on $W$, $c_3(W)<\gamma$; (KF) gives $L_2(D_3)\le
L_2(Y)=L_2(D_3)-v(W)$.

Each item uses only items proved before it.
\end{proof}

The germs not excluded are those of $\beta$ at $0+$ and of $\alpha'$ at
$\infty$. Exchanging worlds $2$ and $3$ maps $\alpha\leftrightarrow\beta'$
and $\beta\leftrightarrow\alpha'$ and replaces $t$ by $1/t$; it maps this
pair of possible germs onto itself.

\begin{corollary}\label{cor:tstrong}
$\alpha$ and $\beta'$ have no $t$-germs on $P_2\cup P_3$; $\beta$ has none on
$\{t\ge\eta\}$ and $\alpha'$ none on $\{t\le S\}$, for all $\eta>0$, $S<\infty$.
So Lemma~\ref{lem:slabs} applies to them on these sets.
\end{corollary}

\begin{theorem}\label{thm:pureratio}
For every $\epsilon>0$ and every measurable $W\subseteq P_2\cup P_3$:
\begin{enumerate}[label=(\roman*),leftmargin=*,itemsep=1pt]
\item $\alpha^\pi(W)\le d\,\beta^\pi(W)+\epsilon$;
\item $\alpha'^\pi(W)\ge g_0\,\beta'^\pi(W)-\epsilon$.
\end{enumerate}
For $\theta\in(0,\infty)$ and $\kappa\coloneqq F(\theta+)$:
\begin{enumerate}[label=(\roman*),leftmargin=*,itemsep=1pt,start=3]
\item $\alpha^\pi(W)\le\kappa\beta^\pi(W)+\epsilon$ if $W\subseteq(0,\theta)$;
\item $\alpha^\pi(W)\ge\kappa\beta^\pi(W)-\epsilon$ if $W\subseteq(\theta,\infty)$;
\item $(\alpha^\pi+\alpha'^\pi)(W)\ge\kappa(\beta^\pi+\beta'^\pi)(W)-\epsilon$ if
$W\subseteq(\theta,\infty)$.
\end{enumerate}
The auxiliary choices in the proof depend on $\epsilon$ (and $\theta$) only,
not on $W$.
\end{theorem}
\begin{proof}
Pure parts vanish off $Z_m$, so $\alpha^\pi(A)=\alpha^\pi(A\cap Z_m)\le
\alpha(A\cap Z_m)$ for every $A$ and $m$, and $\beta(A\cap Z_m)\le\beta^\pi(A)+
u(Z_m)$; likewise for $\alpha',\beta'$ with $v$, $c_3$.

(i) Choose $S>1$ with $\alpha((S,\infty))<\epsilon/4$ and $\eta_0\in(0,1)$
with $\alpha((0,\eta_0))<\epsilon/4$ (Theorem~\ref{thm:tgerms}(a)). Let
$U\coloneqq u((S,S+1)\setminus Z_{m_0})>0$ for some $m_0$. By
Corollary~\ref{cor:tstrong} cut $[\eta_0,S]$ into finitely many $t$-slabs
$J_1,\dots,J_N$ with $\beta(J_k)<U/4$. Fix $e\in(0,U/2)$ with $dNe<\epsilon/4$
and $m\ge m_0$ with $d\,u(Z_m)<\epsilon/4$. For $W_k\coloneqq W\cap J_k\cap
Z_m$, apply (P1) at $B_S$ with a clean $W_1\subseteq(S,S+1)$ of $u$-mass
$\beta(W_k)+e<U$: $\alpha(W_k)\le c_2(W_1)\le d(\beta(W_k)+e)$, since $F\le d$.
Summing over $k$,
\[
 \alpha^\pi(W)\le\tfrac\epsilon2+\alpha(W\cap[\eta_0,S]\cap Z_m)\le
 \tfrac\epsilon2+d\bigl(\beta^\pi(W)+u(Z_m)\bigr)+dNe\le d\beta^\pi(W)+\epsilon.
\]

(ii) Choose $\eta_0>0$ and $S<\infty$ with $g_0\beta'((0,\eta_0))<
\epsilon/4$ and $g_0\beta'((S,\infty))<\epsilon/4$
(Theorem~\ref{thm:tgerms}(b)). Let $C_0\coloneqq v((\frac{\eta_0}2,\eta_0)
\setminus Z_{m_0})>0$, and cut $[\eta_0,S]$ into $t$-slabs $J_k$, $k\le N$,
with $\alpha'(J_k)<C_0/4$ (Corollary~\ref{cor:tstrong}). Fix $e\in(0,C_0/2)$
with $Ne<\epsilon/4$ and $m\ge m_0$ with $v(Z_m)<\epsilon/4$. For
$W_k\coloneqq W\cap J_k\cap Z_m$, apply (P3) at $B'_{\eta_0}$ with a clean
$W_4\subseteq(\frac{\eta_0}2,\eta_0)$ of $v$-mass $\alpha'(W_k)+e$:
$g_0\beta'(W_k)\le\alpha'(W_k)+e$. Summing,
\[
 g_0\beta'^\pi(W)\le\tfrac\epsilon2+g_0\beta'(W\cap[\eta_0,S]\cap Z_m)\le
 \tfrac\epsilon2+\alpha'^\pi(W)+v(Z_m)+Ne\le\alpha'^\pi(W)+\epsilon.
\]

(iii) Fix $\epsilon_1>0$ with $\epsilon_1\beta^\pi(P_2\cup P_3)<\epsilon/4$,
then $\delta>0$ with $F((\theta+\delta)-)<\kappa+\epsilon_1$, and
$U\coloneqq u((\theta,\theta+\delta)\setminus Z_{m_0})>0$. Choose
$\eta_0\in(0,\theta/2)$ and $\eta_1\in(0,\theta/2)$ with
$\alpha((0,\eta_0))<\epsilon/8$ and $\alpha((\theta-\eta_1,\theta))<
\epsilon/8$ (no germ of $\alpha$ at $0+$, no left germ at $\theta$). Cut
$[\eta_0,\theta-\eta_1]$ into slabs $J_k$, $k\le N$, with $\beta(J_k)<U/4$.
Fix $e$ and $m$ with $(\kappa+\epsilon_1)(Ne+u(Z_m))<\epsilon/4$, $e<U/2$.
(P1) at $B_\theta$ with clean $W_1\subseteq(\theta,\theta+\delta)$ of $u$-mass
$\beta(W_k)+e$ gives $\alpha(W_k)\le(\kappa+\epsilon_1)(\beta(W_k)+e)$.
Summing, $\alpha^\pi(W)\le\frac\epsilon4+(\kappa+\epsilon_1)(\beta^\pi(W)+
u(Z_m)+Ne)\le\kappa\beta^\pi(W)+\epsilon$.

(iv) Choose $\eta_1>0$ with $\kappa\beta((\theta,\theta+\eta_1))<\epsilon/4$
(no right germ of $\beta$ at $\theta$) and $S>\theta+\eta_1$ with
$\kappa\beta((S,\infty))<\epsilon/4$. Put $\theta'\coloneqq\theta+\eta_1/2$
and $C_0\coloneqq c_2((\theta,\theta')\setminus Z_{m_0})>0$; cut
$[\theta+\eta_1,S]$ into slabs $J_k$, $k\le N$, with $\alpha(J_k)<C_0/4$; fix
$e<C_0/2$ and $m$ with $Ne+c_2(Z_m)<\epsilon/4$. (P2) at $B_{\theta'}$, with
$W_k\coloneqq W\cap J_k\cap Z_m\subseteq(\theta',\infty)$ and a clean
$W_2\subseteq(\theta,\theta')$ of $c_2$-mass $\alpha(W_k)+e$, gives
$\beta(W_k)\le u(W_2)\le c_2(W_2)/F(\theta+)$, since $F\ge F(\theta+)=\kappa$
on $(\theta,\theta')$. So $\kappa\beta(W_k)\le\alpha(W_k)+e$, and summing,
$\kappa\beta^\pi(W)\le\frac\epsilon2+\alpha^\pi(W)+c_2(Z_m)+Ne\le
\alpha^\pi(W)+\epsilon$.

(v) Add (iv) and (ii), each with $\epsilon/2$, and use $g_0>d\ge\kappa$.
\end{proof}

\section{Localization: world 1 weighs dirty debt exactly}\label{sec:local}

For a bounded Borel function $a$ of the coordinates and a charge $\nu$, the
integral $\int_Qa\,d\nu$ is defined as the uniform limit of integrals of
simple functions \cite[Chapter~4]{RR}; for bounded measurable $a,b$ and a
charge $\pi$, the charge $b\cdot\pi\colon Q\mapsto\int_Qb\,d\pi$ satisfies
$\int a\,d(b\cdot\pi)=\int ab\,d\pi$ (true for simple $a$, then by uniform
limits). The forced moves (M1), (M2), read with a dirty set, are the
\emph{bracket links} used below: if $Q\subseteq D_x\cap\{x<a\}$ is clean,
$W\subseteq D_x\cap\{x\ge a\}$ and $c_2(Q)<\nu_{21}(W)$, then $Q\preceq_dW$;
if $W\subseteq D_x\cap\{x<a\}$, $Q'\subseteq D_x\cap\{x\ge a\}$ is clean and
$c_2(Q')>\nu_{21}(W)$, then $W\preceq_dQ'$; and the same on $D_y$ with
$\nu_{31}$, $c_3$.

\begin{theorem}\label{thm:puresep}
Let $b,b'<\infty$. For every measurable $Q\subseteq P_1\cap\{x<b,\ y<b'\}$,
\[
 \int_Qf(x)\,d\nu_{21}=\int_Qh(y)\,d\nu_{31},\qquad\text{equivalently}\qquad
 \int_Qf(x)\,d\pi_{21}=\int_Qh(y)\,d\pi_{31}.
\]
\end{theorem}
\begin{proof}
The two forms are equivalent because $f(x)c_2=h(y)c_3$ a.e.\ on $P_1$
(Theorem~\ref{thm:sep}). Put $\mathrm{Rect}\coloneqq P_1\cap\{x<b,\ y<b'\}$
and fix $\epsilon\in(0,1)$ and $\epsilon'>0$.

\emph{Residue.} Choose $\alpha_0\in(0,\min(b,b')/2)$ such that $\nu_{21}$ and
$\nu_{31}$ each give mass less than $\epsilon'$ to $P_1\cap\{x<\alpha_0\}$
and to $P_1\cap\{y<\alpha_0\}$; this uses Theorem~\ref{thm:zerogerm},
Corollary~\ref{cor:mirrorgerms} and Theorem~\ref{thm:crossgerm} at $x_0=0$,
$y_0=0$. Apply the construction of Lemma~\ref{lem:windows} to $f$ on
$[\alpha_0,b)$ and to $h$ on $[\alpha_0,b')$. The large jumps are finitely
many points; choose their deleted neighbourhoods so small that, in total,
$P_1\cap\{x\in N_x\}$ and $P_1\cap\{y\in N_y\}$ (with $N_x,N_y$ the unions of
these neighbourhoods) carry $\nu_{21}$- and $\nu_{31}$-mass less than
$\epsilon'$ each; this is possible by Corollary~\ref{cor:strong} and
Theorem~\ref{thm:crossgerm}. This gives finitely many $\epsilon$-windows
$U_i=[a_i,a_i')$ of $f$ and $V_j=[c_j,c_j')$ of $h$. Then choose $\theta_1>0$
so small that the slices of width $\theta_1$ at the finitely many window
endpoints carry, in total, $\nu_{21}$- and $\nu_{31}$-mass less than
$\epsilon'$ each, in both coordinates. Put $U_i^o\coloneqq[a_i+\theta_1,a_i'-\theta_1)$, $V_j^o
\coloneqq[c_j+\theta_1,c_j'-\theta_1)$, $R_{ij}\coloneqq P_1\cap\{x\in
U_i^o,\ y\in V_j^o\}$, and let $\mathrm{Rz}$ be the rest of
$\mathrm{Rect}$. Then $\nu_{21}(\mathrm{Rz})<6\epsilon'$ and
$\nu_{31}(\mathrm{Rz})<6\epsilon'$.

\emph{Brackets.} Let $\mathrm{cap}>0$ be the minimum, over all $i,j$, of the
clean $c_2$-masses of $P_2\cap[a_i,a_i+\theta_1)$ and
$P_2\cap[a_i'-\theta_1,a_i')$ and the clean $c_3$-masses of
$P_3\cap\{c_j\le y<c_j+\theta_1\}$ and $P_3\cap\{c_j'-\theta_1\le y<c_j'\}$
(Lemma~\ref{lem:cells}). Cut each $Q\cap R_{ij}$ into finitely many pieces
$w$ with $\nu_{21}(w)<\mathrm{cap}/2$ and $\nu_{31}(w)<\mathrm{cap}/2$, using
$x$-slabs for $\nu_{21}$ and $y$-slabs for $\nu_{31}$
(Corollary~\ref{cor:strong}). Fix a piece $w\subseteq R_{ij}$ and $e>0$
small.
\begin{itemize}[leftmargin=*,itemsep=2pt]
\item Lower $x$, upper $y$: a clean $Q_1\subseteq P_2\cap[a_i,a_i+\theta_1)$
with $c_2(Q_1)=(1-e)\nu_{21}(w)$ satisfies $Q_1\preceq_dw$ (bracket at
$a_i+\theta_1$); a clean $Q_2\subseteq P_3\cap\{c_j'-\theta_1\le y<c_j'\}$
with $c_3(Q_2)=\nu_{31}(w)+e$ satisfies $w\preceq_dQ_2$ (bracket at
$c_j'-\theta_1$). So $Q_1\preceq_dQ_2$, and Corollary~\ref{cor:agree} gives
$\Phi(Q_1)\le\Phi(Q_2)$:
\[
 f(a_i)(1-e)\nu_{21}(w)\le h(c_j'-)(\nu_{31}(w)+e)\le(1+\epsilon)h(c_j)
 (\nu_{31}(w)+e).
\]
\item Lower $y$, upper $x$: a clean $Q_3\subseteq P_3\cap\{c_j\le
y<c_j+\theta_1\}$ with $c_3(Q_3)=(1-e)\nu_{31}(w)$ and a clean
$Q_4\subseteq P_2\cap[a_i'-\theta_1,a_i')$ with $c_2(Q_4)=\nu_{21}(w)+e$ give
$Q_3\preceq_dw\preceq_dQ_4$, hence
$h(c_j)(1-e)\nu_{31}(w)\le(1+\epsilon)f(a_i)(\nu_{21}(w)+e)$.
\end{itemize}
Let $e\to0$. On $w$, $f\in[f(a_i),(1+\epsilon)f(a_i)]$ and
$h\in[h(c_j),(1+\epsilon)h(c_j)]$, so
\[
 (1+\epsilon)^{-2}\int_wf\,d\nu_{21}\le\int_wh\,d\nu_{31}\le(1+\epsilon)^2
 \int_wf\,d\nu_{21}.
\]

\emph{Summation.} Summing over the pieces and adding the residue,
\[
 \Bigl|\int_Qf\,d\nu_{21}-\int_Qh\,d\nu_{31}\Bigr|\le
 \bigl((1+\epsilon)^2-1\bigr)(1+\epsilon)^2f(b)\nu_{21}(P_1)+
 f(b)\nu_{21}(\mathrm{Rz})+h(b')\nu_{31}(\mathrm{Rz}).
\]
Since $\epsilon$ and $\epsilon'$ are arbitrary, the two integrals are equal.
\end{proof}

In other words, the charges $f\cdot\pi_{21}$ and $h\cdot\pi_{31}$ coincide on
every bounded rectangle of $P_1$, and so do $f\cdot\nu_{21}$ and
$h\cdot\nu_{31}$.

\begin{corollary}\label{cor:onesidedinfty}
(a) If $Q\subseteq P_1\cap\{x<b\}$ with $b<\infty$, then
$\pi_{31}(Q)\le\int_Q(f/h)\,d\pi_{21}$. (b) If $Q\subseteq P_1\cap\{y<b\}$,
then $\pi_{21}(Q)\le\int_Q(h/f)\,d\pi_{31}$.
\end{corollary}
\begin{proof}
(a) On $Q$, $f/h\le f(b)/h(0+)$, which is finite by Lemma~\ref{lem:pos}. For
$Y<\infty$, Theorem~\ref{thm:puresep} on $P_1\cap\{x<b,\ y<Y\}$, integrated
against the bounded function $1/h$, gives
$\pi_{31}(Q\cap\{y<Y\})=\int_{Q\cap\{y<Y\}}(f/h)\,d\pi_{21}\le\int_Q(f/h)\,
d\pi_{21}$. Also $\pi_{31}(Q\cap\{y\ge Y\})\le\nu_{31}(D_y\cap\{y\ge Y\})\to0$
as $Y\to\infty$ (Lemma~\ref{lem:tails}). (b) is the mirror.
\end{proof}

\begin{definition}\label{def:Phihat}
For $Q\subseteq R_b$ put $Q_x\coloneqq Q\cap(P_2\cup(P_1\cap\{x<b\}))$,
$Q_y\coloneqq Q\setminus Q_x$, and
\[
 \Phihat(Q)\coloneqq\int_{Q_x}f(x)\,d\nu_{21}+\int_{Q_y}h(y)\,d\nu_{31}.
\]
For arbitrary $Q\subseteq H'$ put $\Phihat(Q)\coloneqq\sup_b\Phihat(Q\cap
R_b)$.
\end{definition}

The integrands are bounded: $Q_x\subseteq D_x\cap\{x<b\}$ and $Q_y\subseteq
D_y\cap\{y<b\}$. The value does not depend on the admissible $b$: passing
from $b$ to $b'>b$ moves $Q\cap P_1\cap\{b\le x<b'\}$, which lies in
$\{y<b\}$ because $Q\subseteq R_b$, from the second integral to the first,
and Theorem~\ref{thm:puresep} on $P_1\cap\{x<b',\ y<b\}$ shows that the two
values agree. So $\Phihat$ is a finite charge on the subsets of each $R_b$,
nondecreasing in $b$, and $\Phihat=\Phi+\Phi^\pi$ with $\Phi^\pi$ given by the
same formula with $\pi_{21},\pi_{31}$. On clean sets $\Phihat=\Phi$.

\begin{lemma}\label{lem:ivt}
Let $A$ be bounded and $\delta>0$. For every $c\in[0,\Phihat(A)]$ there is
$A'\subseteq A$ with $|\Phihat(A')-c|<\delta$. The same holds with
$\Phihat$ replaced by $\nu_{21}$ on subsets of $D_x$, or by $\nu_{31}$ on
subsets of $D_y$.
\end{lemma}
\begin{proof}
Let $A\subseteq R_b$. By Corollary~\ref{cor:strong}, $A_x$ has a finite
partition into $x$-slabs with $\Phi^\pi$-mass at most $f(b)\pi_{21}<\delta/2$
each, and $A_y$ one into $y$-slabs with $\Phi^\pi$-mass at most
$h(b)\pi_{31}<\delta/2$ each. Order the pieces $A_1,\dots,A_N$. In each piece
choose an increasing family $(S_r)_{r\in[0,1]}$ with $S_0=\varnothing$,
$S_1=A_i$ and $\Phi(S_r)=r\Phi(A_i)$; it exists because $\Phi$ is atomless.
Then $r\mapsto\Phihat(A_1\cup\dots\cup A_{i-1}\cup S_r)$ is nondecreasing, the
sum of a continuous function and a nondecreasing function of total increase
less than $\delta/2$; so it takes every value between its endpoints up to
an error less than $\delta/2$. Concatenating over $i$ gives the claim. The
same argument applies to $\nu_{21}$ and $\nu_{31}$.
\end{proof}

\begin{lemma}\label{lem:matching}
Let $Q_1,\dots,Q_M$ (sources) and $T_1,\dots,T_N$ (targets) be disjoint
bounded sets, let $\lambda_s$ and $\lambda_t$ be finite nonnegative charges on
their subsets with the approximation property of Lemma~\ref{lem:ivt}
at every accuracy (for every $\delta'>0$, every subset $A$ of a source,
resp.\ target, and every $c\in[0,\lambda(A)]$ there is $A'\subseteq A$ with
$|\lambda(A')-c|<\delta'$), let $\delta>0$, and suppose
$\sum_i\lambda_s(Q_i)+2\delta(M+N)<\sum_j\lambda_t(T_j)$. Then there are
finitely many pairs $(Q'_k,T'_k)$ such that the $Q'_k$ partition
$\bigcup_iQ_i$, each $Q'_k$ lies in some $Q_i$, the $T'_k$ are pairwise
disjoint, each lies in some $T_j$, and $\lambda_t(T'_k)>\lambda_s(Q'_k)$ for
every $k$.
\end{lemma}
\begin{proof}
Treat the sources and targets in order, keeping a current source $Q$ and a
current target $T$ (the unused parts). If
$\lambda_t(T)\ge\lambda_s(Q)+2\delta$, cut $T'\subseteq T$ with
$|\lambda_t(T')-(\lambda_s(Q)+\delta)|<\delta/2$ (accuracy $\delta/2$)
and pair it with $Q$; the
source is finished. If $\lambda_t(T)\le2\delta$, discard $T$. Otherwise
$2\delta<\lambda_t(T)<\lambda_s(Q)+2\delta$; cut $Q'\subseteq Q$ with
$|\lambda_s(Q')-(\lambda_t(T)-\frac32\delta)|<\delta/2$, or take $Q'=Q$ if
$\lambda_t(T)-\frac32\delta>\lambda_s(Q)$, and pair $Q'$ with $T$; then
$\lambda_s(Q')<\lambda_t(T)-\delta$, and the target is finished. Each step finishes a source or a target, so
there are at most $M+N$ steps. A step that finishes a source uses at most
$\lambda_s(Q)+\frac32\delta$ of target mass; a step that finishes or discards
a target uses at most $2\delta$ more target mass than source mass. So the
targets are not exhausted before the sources.
\end{proof}

\begin{theorem}\label{thm:dirtyprimary}
Let $B,C\subseteq H'$ be disjoint, $B$ bounded, and
$\Phihat(B)<\Phihat(C)$. Then $B\prec_dC$. Neither set needs to be clean.
\end{theorem}
\begin{proof}
We follow the proof of Theorem~\ref{thm:debtprimary} with $\Phihat$ in place
of $\Phi$ for $B$ and $C$. Intermediate sets are clean, as before, so their
$\Phi$-values are unchanged. We list every change.

\emph{Step 0.} $\Phihat(C\cap R_b)\uparrow\Phihat(C)$ by definition. Choose a
bounded $C_0\subseteq C$ with $\Phihat(B)<\Phihat(C_0)<\infty$ and
$C\setminus C_0$ non-null (if $\Phihat(C)<\infty$, remove from a suitable
$C\cap R_b$ a non-null piece of small $\Phihat$-mass, which exists by FAED,
Lemma~1; if $\Phihat(C)=\infty$, $C\setminus R_b$ has infinite $\Phihat$-mass
and is non-null). Put $\Omega\coloneqq\Phihat(C_0)-\Phihat(B)$.

\emph{Step 1.} In addition to the requirements on the density masses of the
residues, choose $a_*$, the deleted neighbourhoods of the large jumps and
the slice widths so small that $\pi_{21}$ of the $x$-residue zones and
$\pi_{31}$ of the $y$-residue zones (of $B$ and of $C_0$) are less than
$\gamma\coloneqq\Omega/(16(1+\epsilon)^6(1+f(\xi_b)+h(y_b)+f(b_0)+h(b_0)))$.
This is possible in the order of Theorem~\ref{thm:debtprimary}: $\xi_b,y_b$ are fixed
first, the large jumps depend only on $\epsilon$ (by
Lemma~\ref{lem:pos}, $f(0+),h(0+)>0$, so $\log f$ and $\log h$ have finite
variation on $(0,b_0]$, and only finitely many jumps of size at least
$\log(1+\epsilon)/2$ there, whatever $a_*$ is), and all residue zones are
neighbourhoods of finitely many points or of $0$, which carry small pure mass
by Corollary~\ref{cor:strong}.

\emph{Step 3.} Subdivide the source pieces into sub-pieces $Q$ of small
charge $\nu_{21}(Q)$ (types X and $R_x$) or $\nu_{31}(Q)$ (types Y and $R_y$),
using finer slabs and atomlessness of the densities. The first link of a
type X route is the bracket link from the possibly dirty $Q$, in the window
$[a,b)$, to a clean $W_1\subseteq HL\cap\{b-\theta_U\le x<b\}$ with
$c_2(W_1)\in(\nu_{21}(Q),(1+\epsilon)\nu_{21}(Q)]$. Since $f\ge f(a)$ on
$Q$, $\Phihat(Q)\ge f(a)\nu_{21}(Q)$, and
$\Phi(W_1)\le(1+\epsilon)f(a)c_2(W_1)\le(1+\epsilon)^2\Phihat(Q)$. Put
$e(Q)\coloneqq(1+\epsilon)^2\Phihat(Q)$. Type Y is the same with $h$,
$\nu_{31}$, $c_3$, and $e(Q)\coloneqq(1+\epsilon)^4\Phihat(Q)$. For residue
pieces the first link uses $c_2(W_1)\in(\nu_{21}(Q),(1+\epsilon)\nu_{21}(Q)]$,
resp.\ $c_3$, and $e(Q)\coloneqq(1+\epsilon)f(\xi_b)\nu_{21}(Q)$, resp.\
$(1+\epsilon)h(y_b)\nu_{31}(Q)$. All later links are between clean chosen
sets and are unchanged.

\emph{Step 4.} The last link into a possibly dirty target sub-piece $T'$ in a
window $[a,b)$ of $f$ is the bracket link from a clean $S$; it needs
$c_2(S)<\nu_{21}(T')$. Now $c_2(S)\le\Phi(S)/f(a)$ and $\nu_{21}(T')\ge
\Phihat(T')/((1+\epsilon)f(a))$, so $\Phihat(T')>(1+\epsilon)\Phi(S)$
suffices; likewise in $y$. So the route exists whenever
$\Phihat(T')>(1+\epsilon)^5e(Q)$.

\emph{Step 5.} The budget becomes
\[
 (1+\epsilon)^5\sum_Qe(Q)\le(1+\epsilon)^9\Phihat(B)+(1+\epsilon)^6
 \bigl(f(\xi_b)\nu_{21}(R_x)+h(y_b)\nu_{31}(R_y)\bigr),
\]
against the target mass $\Phihat(C_0)-\Phihat(R')$. With $\epsilon$ chosen
so that $(1+\epsilon)^9\Phihat(B)<\Phihat(B)+\Omega/4$, and with the
density and pure parts of the residue term and of $\Phihat(R')$ each less
than $\Omega/16$ (by Step~1 and the absolute continuity argument of
Theorem~\ref{thm:debtprimary}), the target mass exceeds the source budget by
at least $\Omega/2$. Instead of exact splitting, use
Lemma~\ref{lem:matching} with $\lambda_s\coloneqq(1+\epsilon)^5e$ and
$\lambda_t\coloneqq\Phihat$ and $\delta$ so small that $2\delta(M+N)<
\Omega/4$, where $M,N$ are the numbers of source and target sub-pieces,
fixed before $\delta$. Both charges have the approximation property, by
Lemma~\ref{lem:ivt}.

\emph{Step 6} is unchanged. No other step of the proof of
Theorem~\ref{thm:debtprimary} uses the weight of $B$ or $C$.
\end{proof}

\begin{corollary}\label{cor:dirtyagree}
(a) If $B\preceq_dC$ and $C$ is bounded, then $\Phihat(B)\le\Phihat(C)$.
(b) $\Phihat(H')<\infty$.
\end{corollary}
\begin{proof}
(a) As in Corollary~\ref{cor:agree}, with Theorem~\ref{thm:dirtyprimary}; for
unbounded $B$, $\Phihat(B)=\Phihat(B\setminus C)+\Phihat(B\cap C)$ holds by
passing to the supremum over $b$. (b) Let $B_0\subseteq D_x\cap\{x<1\}$ be
clean and non-null. By Lemma~\ref{lem:noncomp}(a), $B_0\not\preceq_d
D_x\cap\{x\ge s\}$ for $s\ge S_0\ge1$. If $\Phihat(D_x\cap\{x\ge s\})>
\Phi(B_0)$ for such $s$, then some bounded $E\coloneqq D_x\cap\{x\ge s\}\cap
R_{b}$ has $\Phihat(E)>\Phi(B_0)$, and Theorem~\ref{thm:dirtyprimary} gives
$B_0\prec_dE\preceq_dD_x\cap\{x\ge s\}$, a contradiction. So
$\Phihat(D_x\cap\{x\ge s\})\le\Phi(B_0)$, and likewise
$\Phihat(D_y\cap\{y\ge s\})\le\Phi(B'_0)$ for large $s$. For $b\ge s$ and
$b'>b$, $R_{b'}\setminus R_b\subseteq F_b\subseteq(D_x\cap\{x\ge b\})\cup
(D_y\cap\{y\ge b\})$, so $\Phihat(R_{b'})\le f(b)\nu_{21}(H)+h(b)\nu_{31}(H)
+\Phi(B_0)+\Phi(B'_0)$, uniformly in $b'$.
\end{proof}

\section{Silent targets in the finitely additive case}\label{sec:silentfa}

We keep the fixed targets of Section~\ref{sec:silent} and put the pure parts
into the accounting. World~$1$'s debt weight is $\Phihat$, which is tight by
Theorem~\ref{thm:dirtyprimary}; credit gets a crude weight that is only an
upper bound, certified by a change at $V$ (Lemma~\ref{lem:Vcap}).

\paragraph{Credit weight.} Put
\[
 \Psitil(T)\coloneqq h(1)\Bigl[\nu_{32}(T\cap D_y)+\frac{\nu_{22}(T\cap
 D_x)}{F(1)}\Bigr],\qquad
 \psitil\coloneqq h(1)\Bigl[u\one_{D_y}+\frac{v\one_{D_x}}{F(1)}\Bigr],
\]
so $\Psitil$ is a finite charge with density $\psitil$ and pure part
$\Psitil^\pi$ (the same formula with $\pi_{32},\pi_{22}$). For $\kappa>0$ let
$\tilde\lambda_\kappa$ be $\varphi$ for debt, $0$ for silence and
$-\kappa\psitil$ for credit, and with $r=(d+g_0)/2$ put
\[
 s_0'(\kappa)\coloneqq\frac{\min\{F(1)(1-r/g_0),\ r-d\}}{\kappa h(1)}>0 .
\]
For $0<s<s_0'(\kappa)$ let $p\coloneqq(s,1,r)/(1+s+r)$ and
$\tilde\ell_p\coloneqq p_1\tilde\lambda_\kappa+p_2\ell_2+p_3\ell_3$, where
$\ell_2,\ell_3$ are the local penalties of the densities, as in
Section~\ref{sec:silent}.

\begin{lemma}\label{lem:silentfa}
For $0<s<s_0'(\kappa)$ the costs $\tilde\ell_p$ relative to silence are,
divided by $p_2$:
\begin{center}
\begin{tabular}{c|l|l}
 & credit & debt\\\hline
$P_1$ & $v\bigl(1-s\kappa h(1)/F(1)\bigr)+u\bigl(r-s\kappa h(1)\bigr)$ &
$c_2\,[sf-1-rf/h]$\\
$P_2$ & $c_3\,[G(1-s\kappa h(1)/F(1))-r]$ & $u\,[F(sf-1)+r]$\\
$P_3$ & $u\,[r-s\kappa h(1)-F]$ & $c_3\,[sh+G-r]$
\end{tabular}
\end{center}
All entries except the debt entry on $P_1$ are a.e.\ positive. So the
a.e.-unique pointwise minimizer is again $\Sigma_p$: debt on $\Theta_s$,
silence elsewhere, never credit.
\end{lemma}
\begin{proof}
As for Lemma~\ref{lem:silent}, with the credit increment $-\kappa\psitil$.
Now $s\kappa h(1)<F(1)(1-r/g_0)<F(1)$ and $s\kappa h(1)<r-d<r$, so
$1-s\kappa h(1)/F(1)>r/g_0$; hence the $P_2$ credit entry is at least
$c_3[g_0(1-s\kappa h(1)/F(1))-r]>0$, and the $P_3$ credit entry at least
$u(r-d-s\kappa h(1))>0$. The debt entries are as before. No bound on
$\ftil$ or $\htil$ is used.
\end{proof}

\begin{lemma}\label{lem:cornerfa}
Lemma~\ref{lem:corner} holds with $\tilde\lambda_\kappa$, $\tilde\ell_p$,
$s_0'(\kappa)$ and $\delta_M\coloneqq\tilde\ell_p(\cdot,M)-\tilde\ell_p(\cdot,
\Sigma_p)$ in place of $\lambda_\kappa$, $\ell_p$, $s_0(\kappa)$ and the
former $\delta_M$.
\end{lemma}
\begin{proof}
The proof of Lemma~\ref{lem:corner} uses only the shape of the minimizer and
the strict positivity of the other costs, now given by
Lemma~\ref{lem:silentfa}, and the integrable bound
$|\delta_M|\le p_1(\varphi+\kappa\psitil)+p_2\bar e_2+p_3\bar e_3$ on $R_{b_0}$.
\end{proof}

\paragraph{The finitely additive cost.} For a regular $M$ let $\debt_M$ and
$\cred_M$ be its debt and credit sets in $H'$, and for $b<\infty$ put
\[
 \Gamma_b(M)\coloneqq p_1\bigl[\Phihat(\debt_M\cap R_b)-\kappa\Psitil(\cred_M
 \cap R_b)\bigr]+\sum_{k=2,3}p_k\sum_{l=1,2}\nu_{kl}(R^k_l(M)\cap R_b).
\]
Its density part is $\int_{R_b}\tilde\ell_p(g,M(g))\,d\mu$. The rest,
$\Pi_b(M)$, is the same expression with every charge replaced by its pure
part ($\Phi^\pi$, $\Psitil^\pi$, $\pi_{kl}$).

\begin{lemma}\label{lem:pureopt}
Suppose $f$ and $h$ are unbounded. Let $\kappa$, $s$ and $p$ be as in
Lemma~\ref{lem:silentfa}, and $b<\infty$. Then
$\Pi_b(M)\ge\Pi_b(\Sigma_p)-\epsilon$ for every $\epsilon>0$ and every
regular $M$; that is, $\Pi_b(M)\ge\Pi_b(\Sigma_p)$.
\end{lemma}
\begin{proof}
Fix $\epsilon>0$. We choose $\epsilon'>0$ and then $X,Y\ge b$, depending
on $\epsilon,b,p,\kappa$ and the configuration only, not on $M$.
By finite additivity, $\Pi_b(M)-\Pi_b(\Sigma_p)$ is a finite sum over cells:
the sets $E=R_b\cap P\cap\{\Sigma_p=o,\ M=o'\}$, where $P$ runs through the
regions below and $(o,o')$ through the pairs of advice. The contribution of
a cell is $p_1$ times the change of $\Phi^\pi$ (for debt) and of
$-\kappa\Psitil^\pi$ (for credit), plus $p_k$ times the changes of the
$\pi_{kl}$ of the rank sets. We bound every cell from below, using
Theorem~\ref{thm:pureratio}(i), (ii) with an $\epsilon'>0$ fixed at the end;
recall $r-s\kappa h(1)>d$, $g_0(1-s\kappa h(1)/F(1))>r$ and $g_0>r$.

\emph{$P_2\cap R_b$ and $P_3\cap R_b$}, where $\Sigma_p$ is silent. On $P_2$,
$\Psitil^\pi=h(1)\alpha'^\pi/F(1)$; on $P_3$, $\Psitil^\pi=h(1)\beta^\pi$.
Dropping the nonnegative terms $p_1\Phi^\pi$, the four cells contribute at
least
\begin{align*}
 M=\omega_2\text{ on }P_2\ (o_2)&:\ p_2\bigl[r\beta^\pi-\alpha^\pi\bigr]
 \ge p_2\bigl[(r-d)\beta^\pi-\epsilon'\bigr]\ge-p_2\epsilon',\\
 M=\omega_2\text{ on }P_3\ (o_{12})&:\ p_2\bigl[(r-s\kappa h(1))\beta^\pi-
 \alpha^\pi\bigr]\ge-p_2\epsilon',\\
 M=\omega_3\text{ on }P_2\ (o_{13})&:\ p_2\bigl[(1-s\kappa h(1)/F(1))
 \alpha'^\pi-r\beta'^\pi\bigr]\ge-p_2\epsilon',\\
 M=\omega_3\text{ on }P_3\ (o_3)&:\ p_2\bigl[\alpha'^\pi-r\beta'^\pi\bigr]
 \ge-p_2\epsilon',
\end{align*}
all evaluated at the cell. The first two use (i) of
Theorem~\ref{thm:pureratio}, the last two use (ii) and
$g_0(1-s\kappa h(1)/F(1))>r$, $g_0>r$.

\emph{$P_1\cap R_b$.} Here $P_1\cap R_b$ is the disjoint union of
\begin{gather*}
 \mathrm{Rect}_1\coloneqq P_1\cap\{x<b,\ y<Y\},\qquad
 \mathrm{Rect}_2\coloneqq P_1\cap\{b\le x<X,\ y<b\},\\
 \mathrm{Tail}_y\coloneqq P_1\cap\{x<b,\ y\ge Y\},\qquad
 \mathrm{Tail}_x\coloneqq P_1\cap\{x\ge X,\ y<b\}.
\end{gather*}
By Definition~\ref{def:Phihat}, $\Phi^\pi=f\cdot\pi_{21}$ on
$P_1\cap\{x<b\}$ and $\Phi^\pi=h\cdot\pi_{31}$ on $P_1\cap\{x\ge b\}$. Write
$\varrho\coloneqq1/f(x)+r/h(y)$, so $\Theta_s=\{\varrho>s\}$. The cell
contributions are:
\begin{itemize}[leftmargin=*,itemsep=1pt]
\item ($\Sigma_p$ debt, $M$ silent), $E\subseteq\Theta_s$:
$-p_1\Phi^\pi(E)+p_2\pi_{21}(E)+p_3\pi_{31}(E)$;
\item ($\Sigma_p$ silent, $M$ debt), $E\subseteq P_1\setminus\Theta_s$:
$p_1\Phi^\pi(E)-p_2\pi_{21}(E)-p_3\pi_{31}(E)$;
\item ($\Sigma_p$ debt, $M$ credit): the first expression plus
$p_2[(1-s\kappa h(1)/F(1))\pi_{22}+(r-s\kappa h(1))\pi_{32}](E)$, which is at
least the first expression;
\item ($\Sigma_p$ silent, $M$ credit):
$p_2[(1-s\kappa h(1)/F(1))\pi_{22}+(r-s\kappa h(1))\pi_{32}](E)\ge0$.
\end{itemize}
On $\mathrm{Rect}_1$, Theorem~\ref{thm:puresep} gives
$\pi_{31}=(f/h)\cdot\pi_{21}$ (integrate the bounded function $1/h$). So the
first contribution is $p_2\int_Ef(\varrho-s)\,d\pi_{21}\ge0$ and the second
is $p_2\int_Ef(s-\varrho)\,d\pi_{21}\ge0$. On $\mathrm{Rect}_2$,
Theorem~\ref{thm:puresep} on $P_1\cap\{x<X,\ y<b\}$ gives
$\pi_{21}=(h/f)\cdot\pi_{31}$, and the two contributions are
$p_2\int_Eh(\varrho-s)\,d\pi_{31}\ge0$ and $p_2\int_Eh(s-\varrho)\,d\pi_{31}
\ge0$. On $\mathrm{Tail}_y$, Corollary~\ref{cor:onesidedinfty}(a) gives
$\pi_{31}(E)\le\int_E(f/h)\,d\pi_{21}$. The second contribution is then at
least $p_2\int_Ef(s-\varrho)\,d\pi_{21}\ge0$. For the first, drop
$p_3\pi_{31}\ge0$: it is at least $p_2\int_E(1-sf)\,d\pi_{21}$, and on
$\Theta_s$, $1-sf>-rf/h\ge-rf(b)/h(Y)$; so it is at least
$-p_2rf(b)\pi_{21}(E)/h(Y)$. On $\mathrm{Tail}_x$, by
Corollary~\ref{cor:onesidedinfty}(b), the second contribution is at least
$p_2\int_Eh(s-\varrho)\,d\pi_{31}\ge0$, and the first at least
$p_2\int_E(r-sh)\,d\pi_{31}\ge-p_2h(b)\pi_{31}(E)/f(X)$, since $r-sh>-h/f$
on $\Theta_s$.

Summing, $\Pi_b(M)-\Pi_b(\Sigma_p)\ge-4p_2\epsilon'-p_2rf(b)\pi_{21}(P_1)/h(Y)
-p_2h(b)\pi_{31}(P_1)/f(X)$. Choose $\epsilon'$, and then $Y,X\ge b$ with
$h(Y)$, $f(X)$ large, so that this is at least $-\epsilon$. Since
$\epsilon$ is arbitrary, $\Pi_b(M)\ge\Pi_b(\Sigma_p)$.
\end{proof}

\begin{theorem}\label{thm:kkmfa}
Suppose $f$ and $h$ are unbounded. Let $\kappa>0$, and let $s$, $p$, $b_0$
and $\eta$ be as in Lemma~\ref{lem:cornerfa}. There are disjoint bounded sets
$A_3,T$ and $\theta>0$ with $V+\debt(A_3)+\cred(T)\prec_1V$ and
$\Phihat(A_3)>\kappa\Psitil(T)+\theta$.
\end{theorem}
\begin{proof}
Put $\epsilon_2(b)\coloneqq(\nu_{21}+\nu_{22})(D_x\cap\{x\ge b\})$ and
$\epsilon_3(b)\coloneqq(\nu_{31}+\nu_{32})(D_y\cap\{y\ge b\})$; both tend to
$0$ (Lemma~\ref{lem:tails}). Fix $b\ge b_0$ with
$p_2\epsilon_2(b)+p_3\epsilon_3(b)<\eta/8$, and then apply
Lemma~\ref{lem:pureopt} with $\epsilon=\eta/8$ at this $b$. Let $\Sigma^b$ be
$\Sigma_p$ on $R_b$ and credit on $F_b$. By (K1), applied to this fixed
model, there is a selected $C$ with $C\preceq_i\Sigma^b$ for $i=1,2,3$.

(i) By \eqref{eq:W}, $L_k(C)\le L_k(\Sigma^b)$. Split each rank set into its
parts in $R_b$ and $F_b$. On $F_b$, $\Sigma^b$ gives credit, which has
world-2 rank $2$ on $(P_1\cup P_2)\cap F_b\subseteq D_x\cap\{x\ge b\}$ and
rank $0$ on $P_3$, and world-3 rank $2$ on $(P_1\cup P_3)\cap
F_b\subseteq D_y\cap\{y\ge b\}$ and rank $0$ on $P_2$. So the $F_b$-part of
$L_k(C)-L_k(\Sigma^b)$ is at least $-\epsilon_k(b)$, and the $R_b$-part of
$L_k(C)-L_k(\Sigma_p)$ is at most $\epsilon_k(b)$.

(ii) $\Gamma_b(C)-\Gamma_b(\Sigma_p)=\int_{R_b}\delta_C\,d\mu+\Pi_b(C)-
\Pi_b(\Sigma_p)\ge\eta-\eta/8$, by Lemma~\ref{lem:cornerfa} ($C$ is
coherent, $\delta_C\ge0$, $b\ge b_0$) and Lemma~\ref{lem:pureopt}.

(iii) With $\Lambda(M)\coloneqq\Phihat(\debt_M\cap R_b)-\kappa
\Psitil(\cred_M\cap R_b)$, (i) and (ii) give $p_1[\Lambda(C)-
\Lambda(\Sigma_p)]\ge\frac78\eta-\frac18\eta=\frac34\eta$.

(iv) As in Lemma~\ref{lem:kkmsilent}(ii), let $\widetilde\Sigma$ be
$\Sigma_p$ on $R_b$ and $C$ on $F_b$, so $C\preceq_1\widetilde\Sigma$; let
$D$, $A$, $A'$, $T$ and $N$ be as there, so $\widetilde\Sigma=N+\debt(A')$
and $C=N+\debt(A)+\cred(T)$. Since $\Sigma_p$ gives no credit on $R_b$,
$\cred_C\cap R_b=T$; and $\debt_C\cap R_b$ and $\Theta_s\cap R_b$ are the
disjoint unions of $\debt_C\cap\Theta_s\cap R_b$ with $A$ and with $A'$. As
$\Phihat$ is additive on subsets of $R_b$,
\[
 \Lambda(C)-\Lambda(\Sigma_p)=\Phihat(A)-\Phihat(A')-\kappa\Psitil(T)\ge
 \frac{3\eta}{4p_1}\eqqcolon3\theta .
\]

(v) By Lemma~\ref{lem:ivt} choose $A_2\subseteq A$ with
$\Phihat(A')+\theta<\Phihat(A_2)<\Phihat(A')+2\theta$. $A'$ is bounded,
possibly dirty, and disjoint from $A_2$, so Theorem~\ref{thm:dirtyprimary}
gives $A'\prec_dA_2$. Conclude as in Lemma~\ref{lem:kkmsilent}(iii): with
$A_3\coloneqq A\setminus A_2$, $V+\debt(A_3)+\cred(T)\prec_1V$, and
$\Phihat(A_3)=\Phihat(A)-\Phihat(A_2)>\kappa\Psitil(T)+\theta$.
\end{proof}

\begin{lemma}\label{lem:Vcap}
Let $S_{\rm reg}\coloneqq P_3\cap\{1\le t<2\}$, $M_{\rm reg}\coloneqq
P_1\cap\{x<1,\ y<1\}$ and $\omega(\delta)\coloneqq\sup\{c_3(E):E\subseteq
M_{\rm reg},\ u(E)\le\delta\}$. Choose $\tau_c>0$ with
\[
 c_2(S_{\rm reg})-d\tau_c/h(1)>2F(1)\tau_c/h(1),\qquad
 c_3(M_{\rm reg})-\omega(\tau_c/h(1))>3\tau_c/h(1).
\]
Then there is no pair of disjoint sets $A_3,T\subseteq H'$, with $A_3$
bounded, $\Psitil(T)\le\tau_c$, $\Phihat(A_3)>\Psitil(T)$ and
$V+\debt(A_3)+\cred(T)\prec_1V$.
\end{lemma}
\begin{proof}
Suppose there is one, and write $T_k\coloneqq T\cap P_k$. Let $e>0$ be
small. Since $c_2=Fu\le du$ on $P_3$ and $u(T_3)\le\nu_{32}(T_3)\le
\Psitil(T)/h(1)$, $c_2(S_{\rm reg}\cap T)\le d\tau_c/h(1)$; and since
$u(T_1)\le\nu_{32}(T_1)\le\tau_c/h(1)$, $c_3(M_{\rm reg}\cap T)\le
\omega(\tau_c/h(1))$. So we can choose clean sets
\begin{gather*}
 S\subseteq S_{\rm reg}\setminus T,\qquad c_2(S)=\nu_{22}(T_1\cup T_2)+e,\\
 M\subseteq M_{\rm reg}\setminus T,\qquad c_3(M)=\nu_{32}(T_1\cup T_3)+
 c_2(S)/F(1)+e;
\end{gather*}
indeed $\nu_{22}(T_1\cup T_2)\le F(1)\tau_c/h(1)$ and the required $c_3$-mass
is less than $3\tau_c/h(1)$ for small $e$. Let $X\coloneqq
V+\cred(T\cup S)+\debt(M)$. By \eqref{eq:table}, with $S,M$ clean,
\begin{align*}
 L_2(X)-L_2(V)&=\nu_{22}(T_1)+\nu_{22}(T_2)-\nu_{21}(T_3)-c_2(S)-c_2(M)\\
 &\le\nu_{22}(T_1\cup T_2)-c_2(S)=-e,\\
 L_3(X)-L_3(V)&=\nu_{32}(T_1)-\nu_{31}(T_2)+\nu_{32}(T_3)+u(S)-c_3(M)\\
 &\le\nu_{32}(T_1\cup T_3)+c_2(S)/F(1)-c_3(M)=-e,
\end{align*}
using $u=c_2/F\le c_2/F(1)$ on $S\subseteq\{t\ge1\}$. By (K4) and
Lemma~\ref{lem:forced}, $V\preceq_1X$. By strict (G),
$X\prec_1V+\debt(M)+\cred(T)$, since $X$ adds credit on the silent set $S$.
Next, $\varphi=h(y)c_3\le h(1)c_3$ on $M_{\rm reg}$, so
\[
 \Phi(M)\le h(1)c_3(M)=\Psitil(T)+h(1)e\bigl(1+1/F(1)\bigr)<\Phihat(A_3)
\]
for small $e$. Since $M$ is clean, $\Phihat(M\setminus A_3)=\Phi(M)-\Phi(M\cap
A_3)<\Phihat(A_3)-\Phihat(A_3\cap M)=\Phihat(A_3\setminus M)$, and
Theorem~\ref{thm:dirtyprimary} (lighter side $M\setminus A_3$) with (F3) gives
$M\prec_dA_3$. By (F5) in the context $\cred(T)$, which is silent on $M\cup
A_3$, $V+\debt(M)+\cred(T)\prec_1V+\debt(A_3)+\cred(T)$. Hence
\[
 V+\debt(A_3)+\cred(T)\prec_1V\preceq_1X\prec_1V+\debt(M)+\cred(T)\prec_1
 V+\debt(A_3)+\cred(T),
\]
a contradiction.
\end{proof}

The credit set $T$ is never relocated here; it is bounded only through its
own charges. This removes the need for clean credit that
Lemma~\ref{lem:cap} has, and also handles possible germs of $\nu_{22}$,
$\nu_{32}$ at $0+$ (Remark~\ref{rem:notcontrolled}).

\begin{theorem}\label{thm:residualfa}
If $f$ and $h$ are unbounded and have jumps beyond every point, then (H) and
the undominatedness of $V$ are contradictory.
\end{theorem}
\begin{proof}
Let $\tau_c$ be as in Lemma~\ref{lem:Vcap} and $\kappa_c\coloneqq
\max\{1,\Phihat(H')/\tau_c\}$, finite by Corollary~\ref{cor:dirtyagree}.
Let $r=(d+g_0)/2$ and $s_0'\coloneqq s_0'(\kappa_c)$. Choose jump points
$x_1$ of $f$ and $y_1$ of $h$ with $f(x_1+)>2/s_0'$ and $h(y_1+)>2r/s_0'$.
Lemma~\ref{lem:cornerfa} gives $s$, $p$, $b_0$ and $\eta$, and
Theorem~\ref{thm:kkmfa} with $\kappa=\kappa_c$ gives disjoint bounded $A_3,T$
with $V+\debt(A_3)+\cred(T)\prec_1V$ and $\Phihat(A_3)>\kappa_c\Psitil(T)$.
Then $\Psitil(T)<\Phihat(H')/\kappa_c\le\tau_c$ and $\Phihat(A_3)>
\Psitil(T)$, contrary to Lemma~\ref{lem:Vcap}. The order of choices is: the
configuration; $\Phihat(H')$ and $\tau_c$; $\kappa_c$; $r$ and $s_0'$;
$(x_1,y_1)$; $s$, $p$, $b_0$, $\eta$; $b$ (with
$p_2\epsilon_2(b)+p_3\epsilon_3(b)<\eta/8$); then $C$. Lemma~\ref{lem:pureopt}
at this $b$ holds for every regular $M$, and the constants $\epsilon'$, $Y$,
$X$ of its proof are chosen after $b$; since $b$ does not depend on them,
there is no circularity.
\end{proof}

\section{Dense thresholds in the finitely additive case}\label{sec:dense}

Fix $\theta\in(0,\infty)$ and $\kappa\coloneqq F(\theta+)$, so $0<\kappa\le
d<g_0$. Write $\mathcal L\coloneqq L_2+\kappa L_3$, and $\mathcal L^c$, $\mathcal L^s$ for its ordinary
and pure parts:
$\mathcal L^c(M)=\int(\ell_2+\kappa\ell_3)(g,M(g))\,d\mu$ and
$\mathcal L^s(M)=\sum_l\pi_{2l}(R^2_l(M))+\kappa\sum_l\pi_{3l}(R^3_l(M))$ for regular
$M$. Write $L_k^c$, $L_k^s$ for the parts of $L_k$. By DT, Lemma~1, $B_\theta$
minimizes $\mathcal L^c$ over all models, and every coherent minimizer of $\mathcal L^c$ is a
model $B_s$; DT's proof uses its tied-level fact only to obtain $G>\kappa$
a.e., which holds here because $G\ge g_0>d\ge\kappa$.

\begin{lemma}\label{lem:weightedpure}
For every $\epsilon>0$ and every regular $M$,
$\mathcal L^s(M)\ge \mathcal L^s(B_\theta)-\epsilon$.
\end{lemma}
\begin{proof}
On $P_1$, $B_\theta$ gives $o_{23}$, which has rank $0$ at worlds $2$ and
$3$, so $M$'s pure contribution there is not smaller. On $P_2\cup P_3$,
$B_\theta$ gives $\omega_2$ on $(\theta,\infty)$ and silence on
$(0,\theta)$, and $M$ gives $\omega_2$, silence or $\omega_3$. The pure
weighted cost of $M$ relative to $B_\theta$ on a cell $E$ is:
$-\alpha^\pi+\kappa\beta^\pi$ ($E\subseteq(0,\theta)$, $M=\omega_2$);
$\alpha'^\pi-\kappa\beta'^\pi$ ($E\subseteq(0,\theta)$, $M=\omega_3$);
$\alpha^\pi-\kappa\beta^\pi$ ($E\subseteq(\theta,\infty)$, $M$ silent);
$(\alpha^\pi+\alpha'^\pi)-\kappa(\beta^\pi+\beta'^\pi)$ ($E\subseteq(\theta,
\infty)$, $M=\omega_3$). By Theorem~\ref{thm:pureratio} (iii), (ii) (with
$g_0\ge\kappa$), (iv) and (v), each is at least $-\epsilon/4$.
\end{proof}

\begin{lemma}\label{lem:deviation}
Let $s\in[0,\infty]$, let $K_s$ be the credal set of $B_s$, let compact convex
$K_j\to K_s$ in Hausdorff distance $\epsilon_j\to0$, and $C_j\coloneqq
\Dscr_{K_j}$.
\begin{enumerate}[label=(\alph*)]
\item If $s>\theta$, then for large $j$, $\{C_j=\omega_2\}\cap(0,\theta)
\subseteq P_3\cap\{y>Y_j\}$ with $Y_j\to\infty$.
\item If $s<\theta$, then for large $j$,
$\bigl((\theta,\infty)\setminus\{C_j=\omega_2\}\bigr)\cup\bigl(\{C_j=\omega_3\}
\cap(0,\theta)\bigr)\subseteq(P_2\cap\{x>X_j\})\cup(0,\sqrt{\epsilon_j})$
with $X_j\to\infty$.
\end{enumerate}
\end{lemma}
\begin{proof}
On $P_2$, $g=g_3(xt,-t,1)$, so $m\cdot g/g_3=m_1xt-m_2t+m_3$; on $P_3$,
$g=|g_3|(y,t,-1)$, so $m\cdot g/|g_3|=m_1y+m_2t-m_3$. $C_j$ gives $\omega_2$
on $P_2$ iff the first expression is negative for all $m\in K_j$, and
$\omega_3$ iff it is positive for all $m$; on $P_3$, $\omega_2$ iff the second
is positive for all $m$, and $\omega_3$ iff it is negative for all $m$. Every
$m'\in K_s$ has $m'_1=0$, $m'_2\ge1/(1+s)$ and $m'_3\le sm'_2$ (if
$s<\infty$).

(a) Choose $m\in K_j$ within $\epsilon_j$ of $n_s$ (of $e_3$ if $s=\infty$).
On $P_2\cap\{t<\theta\}$, $\omega_2$ would need $m_3<t(m_2-m_1x)\le\theta m_2$,
impossible for large $j$ since $m_3/m_2\to s>\theta$. On
$P_3\cap\{t<\theta\}$, $\omega_2$ needs $m_1y>m_3-m_2\theta\ge c_j$, where
$c_j\coloneqq(s-\theta)/(1+s)-\epsilon_j(1+\theta)$ (or
$1-\epsilon_j(1+\theta)$ if $s=\infty$) is positive for large $j$; as
$m_1\le\epsilon_j$, this forces $y>Y_j\coloneqq c_j/\epsilon_j$.

(b) Every $m\in K_j$ is within $\epsilon_j$ of some $m'\in K_s$. For
$t>\theta$ on $P_2$,
\[
 m_1xt-m_2t+m_3\le\epsilon_jxt-m'_2(t-s)+\epsilon_j(t+1)\le
 t\bigl[\epsilon_jx+\epsilon_j(1+1/\theta)-(1-s/\theta)/(1+s)\bigr],
\]
which is negative unless $x\ge X_j^{(1)}\coloneqq[(1-s/\theta)/(1+s)-
\epsilon_j(1+1/\theta)]/\epsilon_j$. For $t>\theta$ on $P_3$,
$m_1y+m_2t-m_3\ge m'_2(t-s)-\epsilon_j(t+1)>0$ for large $j$. So
$(\theta,\infty)\setminus\{C_j=\omega_2\}\subseteq P_2\cap\{x\ge X_j^{(1)}\}$.
For $t<\theta$, take $m''\in K_j$ within $\epsilon_j$ of $e_2\in K_s$. On
$P_2$, $\omega_3$ needs $\epsilon_jxt>(1-\epsilon_j)t-\epsilon_j$, so either
$t<\sqrt{\epsilon_j}$ or $x>X_j^{(2)}\coloneqq(1-\epsilon_j-
\sqrt{\epsilon_j})/\epsilon_j$. On $P_3$, $\omega_3$ needs
$(1-\epsilon_j)t-\epsilon_j<0$, so $t<\epsilon_j/(1-\epsilon_j)<
\sqrt{\epsilon_j}$ for small $\epsilon_j$. Put
$X_j\coloneqq\min(X_j^{(1)},X_j^{(2)})$.
\end{proof}

\begin{theorem}\label{thm:densefa}
For every $\theta\in(0,\infty)$ there are good bases $C_j=\Dscr_{K_j}$, in
the sense of DT, Section~1, item~(v), with $K_j\to K_\theta$ in Hausdorff
distance. In particular $\tau_j\to\theta$ and $\tau_{2,j}\to\infty$.
\end{theorem}
\begin{proof}
\emph{(a) Targets.} As in DT, Theorem~2, choose $a_j\uparrow\theta$ and
$b_j\downarrow\theta$ with $u(A_j)=u(T_j)=z_j>0$, where $A_j\coloneqq
P_2\cap(\theta,b_j)$ and $T_j\coloneqq P_3\cap(a_j,\theta)$, and a
\emph{clean} non-null $Q_j\subseteq P_1$ with $\mu(Q_j)\to0$ and
\[
 c_2(Q_j)<c_2(P_3\cap(\theta,b_j))+c_2(T_j),\qquad
 c_3(Q_j)<u((a_j,\theta)),
\]
which are DT's conditions (1) and (2). Let $X_j$ be $B_\theta$ silenced on
$Q_j\cup A_j$, with credit ($o_{12}$) added on $T_j$. By (K1) at the fixed
model $X_j$ there is a selected $C_j=\Dscr_{K_j}$ with $C_j\preceq_iX_j$ for
$i=1,2,3$, and \eqref{eq:W} gives $L_k(C_j)\le L_k(X_j)$, $k=2,3$.

\emph{(b) $C_j$ is no $B_{s'}$.} If $s'\in[a_j,b_j]$, $X_j$ is a pointwise
world-1 improvement of $B_{s'}$, strict on $Q_j$: on $P_2$, $X_j$ has debt
only on $\{t\ge b_j\}\subseteq\{t>s'\}$; on $P_3$ it has credit on
$\{t>a_j\}\supseteq\{t>s'\}$; on $P_1$ it removes debt on $Q_j$. So
$X_j\prec_1B_{s'}$, while $C_j\preceq_1X_j$. If $s'>b_j$ (including
$s'=\infty$), the world-2 rank-1 sets give
$L_2^s(B_{s'})-L_2^s(X_j)=\pi_{21}(P_2\cap[b_j,s'))+\pi_{21}(P_3\cap[a_j,s'))
-\pi_{21}(Q_j)\ge0$, as $Q_j$ is clean; and by DT's computation and (1),
$L_2^c(B_{s'})-L_2^c(B_\theta)\ge c_2((\theta,b_j))>L_2^c(X_j)-
L_2^c(B_\theta)$. So $L_2(B_{s'})>L_2(X_j)\ge L_2(C_j)$. If $s'<a_j$, the
world-3 rank sets give $L_3^s(B_{s'})-L_3^s(X_j)=\pi_{32}(P_2\cap(s',b_j))
+\pi_{32}(P_3\cap(s',a_j])-\pi_{31}(Q_j)\ge0$, and by (2)
$L_3^c(B_{s'})-L_3^c(B_\theta)\ge u((a_j,\theta))>c_3(Q_j)=L_3^c(X_j)-
L_3^c(B_\theta)$. So $L_3(B_{s'})>L_3(C_j)$. In all cases $C_j\ne B_{s'}$.

\emph{(c) The targets converge in value.} By Lemma~\ref{lem:diff},
$L_2(X_j)-L_2(B_\theta)=c_2(Q_j)+\alpha(A_j)-\alpha(T_j)$ and
$L_3(X_j)-L_3(B_\theta)=c_3(Q_j)-\beta(A_j)+\beta(T_j)$. The bands $A_j$,
$T_j$ shrink to $\{t=\theta\}$ from the two sides, and $\alpha,\beta$ have no
germs at $\theta$ (Theorem~\ref{thm:tgerms}(a)); so both differences tend to
$0$, and so do their ordinary and pure parts separately.

\emph{(d) Weighted step.} Pass to a subsequence with $K_j\to K_*$ and put
$C_*\coloneqq\Dscr_{K_*}$. By Lemma~\ref{lem:weightedpure},
$\mathcal L^c(C_j)=\mathcal L(C_j)-\mathcal L^s(C_j)\le \mathcal L(X_j)-\mathcal L^s(B_\theta)+\epsilon$, which tends to
$\mathcal L^c(B_\theta)+\epsilon$ by (c). Only ordinary parts are passed to the limit:
$\mathcal L^c(C_j)\to \mathcal L^c(C_*)$ by CBOI, Lemma~4. So $\mathcal L^c(C_*)\le \mathcal L^c(B_\theta)$, and
$C_*$ is a coherent minimizer of $\mathcal L^c$, hence $C_*=B_s$ for some
$s\in[0,\infty]$.

\emph{(e) Separate step: $s=\theta$.} Suppose $s>\theta$. From
$L_2(C_j)\le L_2(X_j)$ and (c),
$L_2^c(C_j)\le L_2^c(B_\theta)+[L_2^s(B_\theta)-L_2^s(C_j)]+o(1)$. The
world-2 rank-1 set of $B_\theta$ is $(0,\theta)$ and its rank-2 set is empty,
so $L_2^s(B_\theta)-L_2^s(C_j)\le\pi_{21}(G_j)\le\alpha(G_j)$ with
$G_j\coloneqq\{C_j=\omega_2\}\cap(0,\theta)$. By
Lemma~\ref{lem:deviation}(a), $G_j\subseteq P_3\cap\{y>Y_j\}$, so
$\beta(G_j)\le\nu_{32}(D_y\cap\{y>Y_j\})\to0$ (Lemma~\ref{lem:tails}), and
(P1) at the selected $B_\theta$, with a clean $W_1\subseteq(\theta,\theta+1)$
of $u$-mass $2\beta(G_j)+1/j$, gives $\alpha(G_j)\le d(2\beta(G_j)+1/j)\to0$.
Hence $L_2^c(C_*)\le L_2^c(B_\theta)$, while $L_2^c(B_s)-L_2^c(B_\theta)=
c_2((\theta,s))>0$. This is a contradiction.

Suppose $s<\theta$. In the same way $L_3^c(C_j)\le L_3^c(B_\theta)+
[L_3^s(B_\theta)-L_3^s(C_j)]+o(1)$. The world-3 rank-$\ge1$ set of
$B_\theta$ is $P_2\cup P_3$ and its rank-2 set is $(\theta,\infty)$. With
$H_j\coloneqq(\theta,\infty)\setminus\{C_j=\omega_2\}$ and
$K'_j\coloneqq\{C_j=\omega_3\}\cap(0,\theta)$, therefore,
\[
 L_3^s(B_\theta)-L_3^s(C_j)\le(\beta+\beta')(H_j)+\beta'(K'_j).
\]
By Lemma~\ref{lem:deviation}(b), $H_j\subseteq\mathcal E_j\coloneqq P_2\cap(\theta,
\infty)\cap\{x>X_j\}$ for large $j$. (P7), second case, at $B_\theta$ with
$W=\mathcal E_j$ and a clean $W_7\subseteq P_3\cap(\frac\theta2,\theta)$ of $c_2$-mass
$(\alpha+\alpha')(\mathcal E_j)+1/j$, gives $(\beta+\beta')(\mathcal E_j)\le c_2(W_7)/F(\frac\theta2)
\to0$, since $(\alpha+\alpha')(\mathcal E_j)\le(\nu_{21}+\nu_{22})(D_x\cap\{x>X_j\})\to0$
(Lemma~\ref{lem:tails}). The part of $K'_j$ in $P_2\cap\{x>X_j\}$ is handled
by (P7), first case, in the same way, and the rest lies in
$(0,\sqrt{\epsilon_j})$, where $\beta'\to0$ because $\beta'$ has no germ at
$0+$ (Theorem~\ref{thm:tgerms}(b)). Hence $L_3^c(C_*)\le L_3^c(B_\theta)$,
while $L_3^c(B_s)-L_3^c(B_\theta)=u((s,\theta))>0$. This is a
contradiction. So every convergent subsequence has limit $K_\theta$, and
$K_j\to K_\theta$.

\emph{(f) Good bases.} For large $j$, $\min_{K_j}m_2>0$ and some $m\in K_j$
has $m_3>0$. If $K_j\subseteq\{m_1=0\}$, then $e_3\notin K_j$, so
Lemma~\ref{lem:noedge} gives $e_2\in K_j$ and $C_j=B_{s'}$, excluded by (b).
So some $m\in K_j$ has $m_1>0$, and by the face lemma the index of $C_j$ has
full support. Every $m\in K_j$ is within $\epsilon_j$ of some $m'\in
K_\theta$, so on $P_2\cap\{x<1,\ 2\theta+1<t<2\theta+3\}$,
$m_1xt-m_2t+m_3\le\epsilon_j(3t+1)-(t-\theta)/(1+\theta)<0$ for large $j$:
$C_j$ gives $o_2$ on this non-null set. By Lemma~\ref{lem:noboth} it gives no
$o_{13}$ on $P_2$. Hence $e_2\in K_j$: otherwise $\min_{K_j}(m_1+m_3)>0$, and
$(1,-\eta,1)$, and an open cone around it, would be accepted for small
$\eta>0$. So $C_j$ is a good base. Finally $\tau_j=\max_{K_j}m_3/m_2\to
\theta$, and $\tau_{2,j}=\min_{m\in K_j,\,m_1>0}m_2/m_1\ge
(1/(1+\theta)-\epsilon_j)/\epsilon_j\to\infty$.
\end{proof}

\begin{corollary}\label{cor:factorfa}
If $f$ or $h$ has bounded multiplicative jump ratios beyond some $a_0$, then
(H) and the undominatedness of $V$ are contradictory.
\end{corollary}
\begin{proof}
Suppose the bound holds for $f$. We run the proof of DT, Theorem~10, with
Theorem~\ref{thm:densefa} in place of DT, Theorem~2. That proof applies DT,
Lemma~9, at the good bases $C_j$; Lemma~9 chooses the sets $Y$, $Y_0'$,
$M_0$, $R$ and $S$ of DT, Lemma~5, by conditions on densities and on the
advice of $C_j$ only, and we choose all of them clean. Then every set chosen
in the rounds of DT, Lemma~5 ($Q_k\subseteq R$, $Y_k\subseteq Y$,
$Y_k'\subseteq Y_0'$, $M_k\subseteq M_0$) is clean. The steps are: the
change at $V$ adding $o_{12}(S)$ and $o_{23}(Q_0)$, which raises $L_3$ only
on the clean $S$; the move (M1) at $E_2(\ell)$ with clean source $Q_k$; the
change at the good base silencing $Y_k$, crediting $Y_k'$ and adding debt on
$M_k$, which raises $L_2$ only on the clean $Y_k$ and $L_3$ only on the clean
$Y_k'$ (Lemma~\ref{lem:forced} with $S=\{1,2,3\}$ and the clean-side rule);
the world-1 worsening by (G); the move (M2) at $E_3(\theta_0)$ with clean
source $M_k$; and the closing move into the set $D_n$ of the proof of DT, Lemma~5
(DT's notation), which may be dirty since only
its source must be clean, followed by Lemma~\ref{lem:closing}. So DT's density
estimates bound the loss changes, and DT's proof goes through. For $h$ use the
mirror. Its proof needs the mirror of Theorem~\ref{thm:densefa}; the
hypotheses are symmetric, and exchanging worlds $2,3$ maps the list of
germs excluded in Theorem~\ref{thm:tgerms} onto itself ($\alpha\leftrightarrow
\beta'$, $\beta\leftrightarrow\alpha'$, $t\leftrightarrow1/t$), so every step
of the proof of Theorem~\ref{thm:densefa} has a valid mirror.
\end{proof}

\section{The main theorem}\label{sec:main}

\begin{theorem}\label{thm:main}
Let $n=3$ and assume (G) at every world and (S). Let world~$1$ be a preorder
satisfying (IB); completeness is not assumed. Let worlds $2,3$ be preorders
with losses $L_2,L_3$ of CBOI's form, with finite charges satisfying
\eqref{eq:positive}, such that \eqref{eq:W} holds for $k=2,3$. Then $V$ is
dominated.
\end{theorem}
\begin{proof}
Suppose $V$ is undominated. All results of
Sections~\ref{sec:prelim}--\ref{sec:totals} and~\ref{sec:faclean}--\ref{sec:dense}
apply. By Theorem~\ref{thm:sep}, $c_3/c_2=f(x)/h(y)$ on $P_1$ with $f,h$
positive and nondecreasing. If $f$ or $h$ has bounded jump ratios beyond some
$a_0$, Corollary~\ref{cor:factorfa} gives a contradiction. Otherwise both
have unbounded jump ratios on every $(a_0,\infty)$, hence both are unbounded
and have jumps beyond every point (as in the proof of
Theorem~\ref{thm:I}), and
Theorem~\ref{thm:residualfa} gives a contradiction.
\end{proof}

Theorems~\ref{thm:I} and~\ref{thm:II} are special cases, with shorter proofs.

\begin{corollary}\label{cor:atmostone}
Let $n=3$, assume (G) and (S), and let every world be a preorder satisfying
(IB). If $V$ is undominated, then at most one world has a loss of CBOI's form
with \eqref{eq:positive} and \eqref{eq:W}.
\end{corollary}
\begin{proof}
If two worlds have such losses, relabel so that they are worlds $2,3$ and
apply Theorem~\ref{thm:main}; world~$1$ then satisfies (IB) by hypothesis.
Relabelling permutes the coordinates of $g$, so problems are then
represented on a half-space $\{g_u>0\}$ in place of $H$. The map
$g\mapsto\operatorname{sgn}(g_1)\,g$ carries it to $H$ off a null set and
preserves $\mu$, so the losses, \eqref{eq:positive} and \eqref{eq:W}
transport to the form of Section~\ref{sec:setting}, as in CCF's proof of
Theorem~16.
\end{proof}

This is CCF, Corollary~20(a), without completeness. By EDL, Theorem~2, the
other two worlds then carry leverage sequences.

\begin{remark}\label{rem:dust}
\emph{Dust.} Germ exclusion does not make the pure parts countably additive
on cones. Let $U_1\supseteq U_2\supseteq\cdots$ be cones in $P_1$, each a
union of finitely many thin cones over a fine grid of the chart
$(\xi_1,\xi_2)$, with $\mu(U_m)=\mu(P_1)/m$ and
$\mu(R\cap U_m)/\mu(U_m)\to\mu(R)/\mu(P_1)$ for every cone $R$ over a
rectangle; and let $\eta(A)\coloneqq\lim_{\mathcal U}\mu(A\cap U_m)/\mu(U_m)$
for a non-principal ultrafilter $\mathcal U$. Then $\eta$ is a pure charge
($\eta(H\setminus U_m)=0$ and $\mu(U_m)\to0$), it agrees with
$\mu/\mu(P_1)$ on cones over rectangles, so it has no germs in $x$ or $y$,
but $\eta(U_m)=1$ for all $m$ while the cones $U_m$ decrease to a null cone.
Such charges are not excluded by Sections~\ref{sec:germs}
and~\ref{sec:edge23}. They are nevertheless handled: the arguments of
Sections~\ref{sec:local}--\ref{sec:dense} use only germ exclusion, the
Pareto bounds and pure separability, and never countable additivity of the
pure parts on any family of sets.
\end{remark}

\begin{remark}\label{rem:harmless}
\emph{Germs that remain.} Not all germs are excluded. The proof leaves open
the germs of $\nu_{22}$ on $D_x$ at $x\to0+$ and of $\nu_{32}$ on $D_y$ at
$y\to0+$ (Remark~\ref{rem:notcontrolled}); the germs of $\beta=\nu_{32}$ on
$P_2\cup P_3$ at $t\to0+$ and of $\alpha'=\nu_{22}$ at $t\to\infty$
(Theorem~\ref{thm:tgerms}); and the cross mass at infinity of
Remark~\ref{rem:notcontrolled}. They are harmless. The first two are pure
credit mass, which enters only Lemma~\ref{lem:Vcap}, through its own charges,
and Lemma~\ref{lem:pureopt}, with nonnegative coefficients. A germ of
$\beta$ at $0+$ makes $\omega_2$ more costly at world~$3$ and a germ of
$\alpha'$ at $\infty$ makes $\omega_3$ more costly at world~$2$; both enter
Theorem~\ref{thm:pureratio} and Lemmas~\ref{lem:pureopt},
\ref{lem:weightedpure} only on the favourable side of the inequalities, and
the proof of Theorem~\ref{thm:pureratio} cuts them off by slabs of small
$\alpha$- resp.\ $\beta'$-mass. The cross mass at infinity enters only
through the one-sided bounds of Corollary~\ref{cor:onesidedinfty}. Whether
these germs can occur at all is moot: by Theorem~\ref{thm:main} no
configuration satisfying the hypotheses has $V$ undominated.
\end{remark}

\section{Remarks}\label{sec:remarks}

\begin{remark}\label{rem:uses}
\emph{Which hypotheses are used where.}
\begin{itemize}[leftmargin=*,itemsep=2pt]
\item \emph{Completeness} is used at no world. Every positive world-1
comparison comes from Lemma~\ref{lem:forced} (competitiveness and
\eqref{eq:W}), (G), (IB) and transitivity; every use of closedness of a
world-1 section produces a non-comparison (Lemmas~\ref{lem:noncomp},
\ref{lem:closing}, \ref{lem:edges}(d)), which is then contradicted by a
positive chain; strictness $X\prec_1Y$ is always used as ``$X\preceq_1Y$ and
not $Y\preceq_1X$''.
\item \emph{(IB)} is used only at world~$1$: in transferring accepted changes
between contexts, in (F1)--(F5), and in compositions.
\item \emph{\eqref{eq:W}} is applied only to selected models (edge models,
$B_s$, $B'_s$, $V_{23}$, $D_2$, $D_3$, KKM dominators) and to $V$, and only in
the form ``$C\preceq_kX\Rightarrow L_k(C)\le L_k(X)$'' or its contrapositive:
in forced comparisons, in KKM accounting at fixed targets, in the Pareto
property, and in Lemma~\ref{lem:otherworld}. No positive comparison at
worlds $2,3$ is derived from the losses.
\item \emph{(S)}: KKM (K1) at fixed targets only ($\Sigma^b$, $X_j$, and
inside DT); (K1) and (K2) at $V$, for (K4); the edge Bayes targets of
Step~1 of the proof of EDL, Theorem~8, and of FAED, Proposition~6, inside
Lemma~\ref{lem:FG}; closedness of $A_1$ at fixed edge models and their
perturbations; closedness of $A_2,A_3$ in CCF, Lemma~3, in
Lemma~\ref{lem:sel23} and in Lemma~\ref{lem:otherworld}, and inside
(K1), (K2) and the EDL/FAED Step~1 just mentioned.
\item At worlds $2,3$, transitivity enters through (K2), strict (G) through
the face lemma, CCF, Lemma~3, and \eqref{eq:order23}. Completeness and (IB)
at worlds $2,3$ are not used.
\end{itemize}
\end{remark}

\begin{remark}\label{rem:relation}
\emph{Relation to CCF and CBOI.} CCF, Theorem~18, proves the conclusion of
Theorem~\ref{thm:main} when the third world is complete, by constructing a
finitely additive weak extension at that world (CCF, Theorem~15) and then
applying CBOI, Theorem~2. Completeness entered there in three places: CCF's
comparison (T1) near an edge selection, the comparative-probability step with
the Savage--Wakker theorem (CCF, Theorem~9; \cite{Savage,Wakker}), and the dichotomy between strict improvements and accepted
changes. Here no extension at world~$1$ is constructed. The routing primaries
replace the comparative probabilities on bounded sets; the silent targets
make the KKM comparison one-sided in credit, so no dichotomy of trades is
needed; and the far tails, where the primaries say nothing, are cut off by
the truncation of the target and by Lemma~\ref{lem:totals}. CBOI, Theorem~2,
needs weak extensions at all worlds (one-sided by EDL, Proposition~1);
Theorem~\ref{thm:main} needs them at two worlds of three, at the price of
(IB) at the third. For $n\ge4$ the method does not apply as it stands: it
uses the geometry of three worlds (the regions $P_1,P_2,P_3$, the two edges
at world~$1$ and the edge $[e_2,e_3]$).
\end{remark}

\begin{remark}\label{rem:notneeded}
\emph{Questions left open because they are not needed.} (a) Whether the
germs listed in Remark~\ref{rem:harmless} can be excluded. (b) A KKM lemma in
which the dominator can be made to agree with a modification of the target
on the concentration sets $Z_m$; the proof uses fixed targets instead. (c)
Whether world~$1$ must have a finitely additive weak extension when worlds
$2,3$ do; Theorem~\ref{thm:main} shows only that the configuration is
impossible.
\end{remark}

\begin{thebibliography}{9}
\bibitem{Sierpinski} W.~Sierpi\'nski, Sur les fonctions d'ensemble additives et
continues, \emph{Fundamenta Mathematicae} 3 (1922), 240--246.
\bibitem{YH} K.~Yosida and E.~Hewitt, Finitely additive measures,
\emph{Transactions of the American Mathematical Society} 72 (1952), 46--66.
\bibitem{PP} T.~Pennanen and A.-P.~Perkki\"o, Topological duals of locally
convex function spaces, \emph{Positivity} 26 (2022), article~2.
\bibitem{RR} K.~P.~S.~Bhaskara Rao and M.~Bhaskara Rao, \emph{Theory of
Charges}, Academic Press, London, 1983.
\bibitem{Savage} L.~J.~Savage, \emph{The Foundations of Statistics},
Wiley, New York, 1954. Chapter~3, \S3.
\bibitem{Wakker} P.~Wakker, Agreeing probability measures for
comparative probability structures, \emph{Annals of Statistics} 9 (1981),
658--662.
\end{thebibliography}
\end{document}
